\subsection{The classification behind Remark 3.11}\label{sec:rem311}
\begingroup
\emergencystretch=3em
\newcommand{\rsq}[1]{(#1,#1)}
\newcommand{\Cyc}{\operatorname{Cyc}}

We will follow the ordinary proof in \path{cgss/proofs/remark311/remark311.tex}.
Its lemma labels are retained. Theorems A, B and C are auxiliary results;
Corollary D is the strengthened form of [CGSS, Theorem 3.8].
All spaces below are finite-dimensional real inner product spaces.
Coordinate vectors $\Zu_i$ are orthonormal and start at index 1;
simple roots $\alpha_i$ start at index 0. A graph called connected is nonempty.
All Lean declarations of ours, as opposed to library declarations, occur in
\verb|\lean| metadata.

\begingroup\raggedright
Every one of these declarations is in Lean. Theorem A and Corollary D
are in \path{LineGraphs/Cartan/}. Theorems B and C are in
\path{LineGraphs/RootSetClassification/}, and Remark (small ranks) and
Lemma (a Gram realization) in \path{LineGraphs/Cartan/SmallRanks.lean} and
\path{LineGraphs/Cartan/GramRealization.lean}. Their
modules are imported by \path{LineGraphs.lean}, and
\path{LineGraphs/RootSetClassification/Audit.lean} prints their axioms.\par
\endgroup

\begin{cgss}{Definition}{3.9}\label{def:roots}\label{def:3.9}
\lean{IsRootSet}\uses{facts:inner}
\assume{$V$ is a finite-dimensional real inner product space.}
A Root System is a finite set $\Phi\subseteq V\setminus\{0\}$ spanning $V$,
such that, for all $\alpha,\beta\in\Phi$,
\[
 \frac{2(\alpha,\beta)}{\rsq\alpha}\in\mathbb Z,
 \qquad s_\alpha\beta:=\beta-
       \frac{2(\alpha,\beta)}{\rsq\alpha}\alpha\in\Phi,
\]
and $t\alpha\in\Phi$ implies $t\in\{1,-1\}$ for every $t\in\mathbb R$.
It is indecomposable if it is nonempty and has no partition into two
nonempty mutually orthogonal parts. A similarity is a linear bijection
$f:V\to W$ carrying the root sets onto one another and satisfying
$(f\Zx,f\Zy)=a(\Zx,\Zy)$ for some $a>0$ and all $\Zx,\Zy$.
\repair\ Nonzero roots, nonemptiness in indecomposability, and the use of
$\Phi$ for the printed $S$ in condition (b) are explicit here.
\end{cgss}

\begin{namedresult}{Lemma}{(composition and inversion of similarities)}\label{lem:rem-similarity-laws}
\lean{IsRootSet.similarity_comp}\lean{IsRootSet.similarity_symm}
\uses{def:roots}
\assume{$\Phi,\Psi,\Omega$ are Root Systems in real inner product spaces.}
\prove{The inverse of any similarity $f:\Phi\to\Psi$ with factor $a>0$
is a similarity with factor $a^{-1}$. Given also a similarity
$g:\Psi\to\Omega$ with factor $b>0$, $g\circ f$ is a similarity with
factor $ab$.}
\end{namedresult}
\begin{proof}
\pfsketch\ Substitute the metric identities and compose or invert the
linear bijections and their equalities of root sets.
\step{rsimcompose}{\assume{$f:\Phi\to\Psi$, $g:\Psi\to\Omega$ are
similarities with factors $a,b>0$.}
\prove{$g\circ f$ is a linear bijection,
$(gf\Zx,gf\Zy)=ab(\Zx,\Zy)$, and $gf(\Phi)=\Omega$.}}
\begin{proof}\pf\ Substitute the two identities of
Definition~\ref{def:roots}; $ab>0$.~\qed\end{proof}
\step{rsiminvert}{\assume{$f:\Phi\to\Psi$ is a similarity with factor $a>0$.}
\prove{$f^{-1}$ is a linear bijection,
$(f^{-1}\Zx,f^{-1}\Zy)=a^{-1}(\Zx,\Zy)$, and $f^{-1}(\Psi)=\Phi$.}}
\begin{proof}\pf\ Evaluate the identity for $f$ on the two inverse
images and divide by $a>0$.~\qed\end{proof}
\qedstep
\begin{proof}\pf\ Composition and inversion of linear bijections,
with \stepref{rsimcompose}, \stepref{rsiminvert}, give the assertions.~\qed\end{proof}
\end{proof}

\begin{namedresult}{Definition}{(bases and Cartan conventions)}\label{def:rem-bases}
\lean{IsRootSet.IsBase}\lean{CartanMatrix.IsFiniteTypeModel}
\uses{def:roots}
\assume{$\Phi$ is a Root System in $V$.}
A root base is a linearly independent finite family in $\Phi$ in which every
root has integer coefficients all nonnegative or all nonpositive. Its matrix is
\begin{equation}\label{eq:pairing}
 M(i,j)=p(i,j):=2(\alpha_i,\alpha_j)/\rsq{\alpha_j}.
\end{equation}
The second root supplies the denominator. A finite integer matrix is finite
Cartan if its diagonal is 2, its other entries are nonpositive, its zeros
are symmetric, and some positive integer diagonal $D$ makes $DM$ symmetric
and positive on all nonzero integer vectors. The permitted models are
\begin{equation}\label{eq:list}
 \begin{gathered}
 \typ A_\ell\ (\ell\geq1),\quad \typ B_\ell\ (\ell\geq2),\quad
 \typ C_\ell\ (\ell\geq3),\quad \typ D_\ell\ (\ell\geq4),\\
 \typ E_6,\ \typ E_7,\ \typ E_8,\ \typ F_4,\ \typ G_2.
 \end{gathered}
\end{equation}
These mean the actual Mathlib matrices, with rank $\ell$ for
\path{CartanMatrix.A} despite its docstring. The model predicate includes
the displayed rank restrictions, and equality after transport to its index type.
Every root base is a real vector-space basis of $V$: its span contains
every root by the coefficient condition, and the roots span $V$ by
Definition~\ref{def:roots}; independence is part of the base definition.
\end{namedresult}

\paragraph{Library interface.}\label{facts:rem-library}
We use Mathlib at commit \path{082e2d37e8b0463410cdb532e111cd43d5a66174}
(4 September 2026), the version pinned in \path{lake-manifest.json}. The following
are the precise library packages used in the proofs; every vector space in
them is over $\mathbb R$ and finite dimensional, and all matrix index sets
are finite with classically chosen decidable equality.
\begingroup\raggedright
\begin{enumerate}
\item \path{Matrix.IsFiniteCartan} and \path{CartanMatrix.A},
\path{CartanMatrix.B}, \path{CartanMatrix.C}, \path{CartanMatrix.D},
\path{CartanMatrix.E}, and the matrices $\typ F_4,\typ G_2$ are in
\path{Mathlib/LinearAlgebra/Matrix/Cartan.lean}.
\path{Matrix.IsIndecomposable} in
\path{Mathlib/LinearAlgebra/Matrix/Block.lean} excludes a reindexing as
$\left(\begin{smallmatrix}A&B\\0&D\end{smallmatrix}\right)$ with both
diagonal block sizes positive. It imposes no nonempty-index hypothesis.
\item \path{Matrix.posDef_iff_dotProduct_mulVec} in
\path{Mathlib/LinearAlgebra/Matrix/PosDef.lean} converts positivity on
finitely supported vectors to the finite quadratic form, over both
$\mathbb Z$ and $\mathbb R$.
\path{Matrix.exists_mulVec_eq_zero_iff} in
\path{Mathlib/LinearAlgebra/Matrix/ToLinearEquiv.lean} identifies a nonzero
kernel vector with zero determinant over any commutative domain.
\path{RingHom.map_det} in
\path{Mathlib/LinearAlgebra/Matrix/Determinant/Basic.lean} commutes
determinants with casts. In the same file, \path{Matrix.det_apply'}
is the signed permutation expansion used for a leaf deletion.
\item \path{isCompact_sphere} in
\path{Mathlib/Topology/MetricSpace/ProperSpace.lean} and
\path{IsCompact.exists_isMaxOn} in
\path{Mathlib/Topology/Order/Compact.lean} give the maximum of a continuous
quadratic form on a nonempty unit sphere. Finite-dimensional Euclidean
spaces are proper. No Perron--Frobenius classification is a premise.
\path{Rat.denseRange_cast}, in
\path{Mathlib/Topology/Algebra/Order/Archimedean.lean}, gives density of
rationals in the reals; \path{DenseRange.piMap}, in
\path{Mathlib/Topology/NhdsWithin.lean}, gives the coordinatewise version.
\item \path{SimpleGraph.IsAcyclic.path_unique} in
\path{Mathlib/Combinatorics/SimpleGraph/Acyclic.lean} gives uniqueness of
paths in an acyclic graph. Connectedness supplies their existence.
The graph definitions are in \path{Connectivity/Connected.lean},
\path{Finite.lean}, and \path{AdjMatrix.lean} under
\path{Mathlib/Combinatorics/SimpleGraph/}.
\item \path{RootPairing} and \path{RootPairing.IsRootSystem} are in
\path{Mathlib/LinearAlgebra/RootSystem/Defs.lean}.
Crystallographic, reduced, and irreducible are defined respectively in
\path{IsValuedIn.lean}, \path{Reduced.lean}, and \path{Irreducible.lean}
in that directory. Irreducibility requires nontrivial root and coroot
spaces and no nonzero proper reflection-invariant submodule in either.
The real bilinear inner product $\texttt{inner}_{\mathtt l}$ is in
\path{Mathlib/Analysis/InnerProductSpace/Basic.lean}.
\path{LinearMap.IsPerfPair.of_injective} in
\path{Mathlib/LinearAlgebra/PerfectPairing/Basic.lean} makes it perfect
from injectivity in both arguments and finite dimensionality.
\item \path{RootPairing.nonempty_base} and \path{RootPairing.Base.mk'} in
\path{Mathlib/LinearAlgebra/RootSystem/BaseExists.lean} require a field
of characteristic zero, finitely many roots, spanning roots and coroots,
crystallographic pairings, and reducedness. The second constructor needs
root independence and membership of each root or its negative in the
additive monoid of the selected roots.
\path{Submodule.span_nat_eq_addSubmonoidClosure} in
\path{Mathlib/LinearAlgebra/Span/Defs.lean} identifies this monoid with
nonnegative integer combinations.
\item \path{RootPairing.Base.support_nonempty} in
\path{Mathlib/LinearAlgebra/RootSystem/Base.lean} needs nonempty roots
and $2\ne0$; \path{RootPairing.Base.toWeightBasis} and
\path{RootPairing.Base.toWeightBasis_apply} need spanning.
\path{Module.Basis.equiv} and \path{Module.Basis.equiv_apply} in
\path{Mathlib/LinearAlgebra/Basis/Defs.lean} extend a bijection of basis
indices to a linear equivalence with the stated values.
\item In \path{Mathlib/LinearAlgebra/RootSystem/CartanMatrix.lean},
\path{RootPairing.Base.cartanMatrix_isFiniteCartan} needs a characteristic
zero domain, finite roots, crystallographic pairings, decidable equality,
and spanning. \path{RootPairing.Base.cartanMatrix_isIndecomposable}
needs a characteristic zero domain, finite roots, crystallographic
pairings, reducedness, and irreducibility.
\path{RootPairing.Base.apply_mem_range_root_of_cartanMatrixEq} needs
finite, spanning, reduced, crystallographic source roots over a
characteristic zero domain; the target needs crystallographic pairings
and a base. Given equality of base matrices and a linear equivalence
taking one base to the other, it maps source roots into target roots.
\path{RootPairing.Base.equivOfCartanMatrixEq} packages both directions
with the basis equivalence above; it additionally requires finite,
spanning, reduced target roots.
\path{RootPairing.Base.algebraMap_cartanMatrixIn_apply} identifies the
integer Cartan entry cast to $\mathbb R$ with the real pairing.
\end{enumerate}
\par\endgroup
The elementary identities of \ref{facts:inner} remain available. Arithmetic
leaves below are substitutions in displayed finite sums or polynomial
identities (the Lean proofs may use \texttt{ring}, \texttt{norm\_num},
or \texttt{linarith}); they do not appeal to Python output.

\begin{namedresult}{Lemma}{(reflection identities)}\label{lem:rem-reflection}
\lean{IsRootSet.reflection_identities}\uses{def:roots,facts:inner}
\assume{$V$ is a real inner product space and $\alpha\ne0$.}
\prove{$s_\alpha$ is an orthogonal linear involution;
 $s_\alpha\alpha=-\alpha$; its $-1$ eigenspace is $\mathbb R\alpha$;
 $s_{-\alpha}=s_\alpha$; and for every linear isometry $f$ onto $W$,
 $s_{f\alpha}=f s_\alpha f^{-1}$. In a Root System these reflections
 permute the roots and $\Phi=-\Phi$.}
\end{namedresult}
\begin{proof}
\pfsketch\ We resolve a vector into its component along $\alpha$ and its
orthogonal component.
\step{rrsplit}{For $t=(\alpha,\Zx)/\rsq\alpha$ and $\Zy=\Zx-t\alpha$,
 $(\alpha,\Zy)=0$.}
\begin{proof}\pf\ Substitute $t$ using bilinearity in \ref{facts:inner}.~\qed\end{proof}
\step{rrformula}{$s_\alpha\Zx=-t\alpha+\Zy$.}
\begin{proof}\pf\ Substitute \stepref{rrsplit} in Definition~\ref{def:roots}.~\qed\end{proof}
\step{rrlinear}{$s_\alpha$ is linear.}
\begin{proof}\pf\ The formula in Definition~\ref{def:roots} is a difference
of linear maps, by \ref{facts:inner}.~\qed\end{proof}
\step{rrinvolution}{$s_\alpha(s_\alpha\Zx)=\Zx$.}
\begin{proof}\pf\ Apply \stepref{rrformula} to the components $-t,\Zy$.~\qed\end{proof}
\step{rrmetric}{$(s_\alpha\Zx,s_\alpha\Zw)=(\Zx,\Zw)$.}
\begin{proof}\pf\ Expand both sides by \stepref{rrsplit},
\stepref{rrformula} and \ref{facts:inner}; both are $tu\rsq\alpha+(\Zy,\Zv)$
for the decompositions $\Zx=t\alpha+\Zy$, $\Zw=u\alpha+\Zv$.~\qed\end{proof}
\step{rreigen}{$s_\alpha\alpha=-\alpha$ and
 $s_\alpha\Zx=-\Zx\iff\Zx\in\mathbb R\alpha$.}
\begin{proof}\pf\ In \stepref{rrformula}, equality to $-\Zx$ is exactly
$2\Zy=0$.~\qed\end{proof}
\step{rrsign}{$s_{-\alpha}=s_\alpha$.}
\begin{proof}\pf\ Substitute $-\alpha$ into Definition~\ref{def:roots}.~\qed\end{proof}
\step{rrconj}{$s_{f\alpha}(f\Zx)=f(s_\alpha\Zx)$.}
\begin{proof}\pf\ Substitute the isometry identities into
Definition~\ref{def:roots}.~\qed\end{proof}
\step{rrroots}{\assume{$\Phi$ is a Root System and $\alpha\in\Phi$.}
\prove{$s_\alpha$ permutes $\Phi$ and $-\alpha\in\Phi$.}}
\begin{proof}\pf\ Restrict \stepref{rrinvolution} and \stepref{rreigen}
using reflection closure in Definition~\ref{def:roots}.~\qed\end{proof}
\qedstep
\begin{proof}\pf\ Combine \stepref{rrlinear}--\stepref{rrroots}.~\qed\end{proof}
\end{proof}

\begin{namedresult}{Lemma}{(a null vector of a nonnegative form)}\label{lem:rem-null}
\lean{Matrix.mulVec_eq_zero_of_quadratic_eq_zero}
\uses{facts:inner}
\assume{\begin{pfenum}
\item $Q$ is real symmetric on a finite index set, and
 $q(\Zx):=\Zx^TQ\Zx\geq0$ for every real vector $\Zx$.
\item $q(\Zy)=0$.
\end{pfenum}}
\prove{$Q\Zy=0$.}
\end{namedresult}
\begin{proof}
\pfsketch\ A quadratic polynomial nonnegative on the line cannot have a
nonzero linear coefficient at a zero.
\step{rnpoly}{For every $\Zz,t$, $0\leq2t\Zz^TQ\Zy+t^2q(\Zz)$.}
\begin{proof}\pf\ Expand $q(\Zy+t\Zz)$ using \asm[1]{0}, \asm[2]{0}.~\qed\end{proof}
\step{rncoeff}{For every $\Zz$, $\Zz^TQ\Zy=0$.}
\begin{proof}
\pflet{$c=\Zz^TQ\Zy$, $d=q(\Zz)\geq0$.}
\step{rnchoice}{\assume{$c\ne0$.}
\prove{For $t=-c/(d+1)$,
 $2tc+t^2d=-c^2(d+2)/(d+1)^2<0$.}}
\begin{proof}\pf\ Substitute $t$, collect the common denominator,
and use $c^2>0$, $d+2>0$, $(d+1)^2>0$.~\qed\end{proof}
\qedstep
\begin{proof}\pf\ The inequality in \stepref{rnchoice} contradicts
\stepref{rnpoly}.~\qed\end{proof}
\end{proof}
\qedstep
\begin{proof}\pf\ Apply \stepref{rncoeff} to each coordinate unit vector. ~\qed\end{proof}
\end{proof}

\begin{namedresult}{Lemma}{(integer positivity)}\label{lem:integer-positive}
\lean{Matrix.posDef_real_of_posDef_int}
\uses{lem:rem-null,facts:rem-library}
\assume{\begin{pfenum}
\item $I$ is finite and $H:I\times I\to\mathbb Z$ is symmetric.
\item $\Zx^TH\Zx>0$ for every nonzero $\Zx\in\mathbb Z^I$.
\end{pfenum}}
\prove{The entrywise real cast of $H$ satisfies
$\Zy^TH\Zy>0$ for every nonzero $\Zy\in\mathbb R^I$.}
\end{namedresult}
\begin{proof}
\pfsketch\ A negative real value would yield a negative integer value
by density and scaling. A real null vector would force an integer null
vector through the determinant. Only the integer matrix is needed here.
\step{riposr}{The form is nonnegative on real vectors.}
\begin{proof}
\step{rinegative}{\assume{Some real vector has negative quadratic value.}
\prove{\textsc{False}.}}
\begin{proof}
\step{ricont}{The negative-value set of $\Zx\mapsto\Zx^TH\Zx$
is nonempty and open.}
\begin{proof}\pf\ A finite sum of products of coordinates is continuous;
use \asm{2} for nonemptiness.~\qed\end{proof}
\step{ridense}{\pick\ a rational vector $\Zv$ with $\Zv^TH\Zv<0$.}
\begin{proof}\pf\ Apply \path{DenseRange.piMap} to
\path{Rat.denseRange_cast} and \stepref{ricont}.~\qed\end{proof}
\step{riclear}{\pick\ an integer $k>0$ with $k\Zv\in\mathbb Z^I$.}
\begin{proof}\pf\ The product of the finitely many positive coordinate
denominators of $\Zv$ is such a $k$.~\qed\end{proof}
\step{riintnegative}{$(k\Zv)^TH(k\Zv)=k^2\Zv^TH\Zv<0$.}
\begin{proof}\pf\ Expand the quadratic form and use
\stepref{ridense}, \stepref{riclear}.~\qed\end{proof}
\qedstep
\begin{proof}\pf\ The vector in \stepref{riintnegative} is nonzero
and integral, contradicting \asm[2]{0} after the order-preserving
integer cast.~\qed\end{proof}
\end{proof}
\qedstep
\begin{proof}\pf\ The contradiction \stepref{rinegative} excludes
every negative value.~\qed\end{proof}
\end{proof}
\step{rinonzero}{\assume{$\Zy\ne0$ and $\Zy^TH\Zy=0$.}
\prove{\textsc{False}.}}
\begin{proof}
\step{rikernel}{$H\Zy=0$.}
\begin{proof}\pf\ Apply Lemma~\ref{lem:rem-null} to \stepref{riposr}
and \asm{1}.~\qed\end{proof}
\step{ridetr}{$\det H=0$ over $\mathbb R$.}
\begin{proof}\pf\ Apply \path{Matrix.exists_mulVec_eq_zero_iff}
to \stepref{rikernel} and $\Zy\ne0$ from \asm{1}.~\qed\end{proof}
\step{ridetz}{$\det H=0$ over $\mathbb Z$.}
\begin{proof}\pf\ Use \path{RingHom.map_det} and injectivity of the
integer cast in \stepref{ridetr}.~\qed\end{proof}
\step{riz}{\pick\ a nonzero $\Zv\in\mathbb Z^I$ with $H\Zv=0$.}
\begin{proof}\pf\ Apply \path{Matrix.exists_mulVec_eq_zero_iff}
over $\mathbb Z$ to \stepref{ridetz}.~\qed\end{proof}
\qedstep
\begin{proof}\pf\ The zero value $\Zv^TH\Zv=0$ from \stepref{riz}
contradicts \asm[2]{0}.~\qed\end{proof}
\end{proof}
\qedstep
\begin{proof}\pf\ Combine nonnegativity \stepref{riposr} and
the exclusion of nonzero null vectors \stepref{rinonzero}.~\qed\end{proof}
\end{proof}

\begin{namedresult}{Lemma}{(connected support)}\label{lem:connected}
\lean{Matrix.isIndecomposable_iff_connected_support}
\uses{facts:rem-library}
\assume{\begin{pfenum}
\item $I$ is finite and nonempty, $R$ is a ring, and $M:I\times I\to R$.
\item $M(i,j)=0\iff M(j,i)=0$ for every $i,j$.
\end{pfenum}}
\prove{The simple graph $i\sim j\iff i\ne j\text{ and }M(i,j)\ne0$
is connected if and only if $M$ is Mathlib-indecomposable.}
\end{namedresult}
\begin{proof}
\pfsketch\ A disconnected component gives two blocks; a block partition
cannot be crossed by a path.
\step{rcforward}{\assume{the support is disconnected.}
\prove{$M$ is decomposable.}}
\begin{proof}
\step{rccomponent}{There is a nonempty proper reachable component $P\subset I$.}
\begin{proof}\pf\ Unfold disconnectedness on the nonempty type $I$.~\qed\end{proof}
\step{rczero}{$M(i,j)=0$ whenever $i\in P$ and $j\notin P$.}
\begin{proof}\pf\ A nonzero entry would give an edge from $i$ to $j$,
contrary to the definition of $P$ in \stepref{rccomponent}.~\qed\end{proof}
\step{rcblocks}{Enumerating $P$ and $I\setminus P$ gives a block diagonal
matrix with both block sizes positive.}
\begin{proof}\pf\ Use \stepref{rczero} and \asm[2]{0} in the four
blocks of the two-set partition. ~\qed\end{proof}
\qedstep
\begin{proof}\pf\ The decomposition in \stepref{rcblocks} violates
\path{Matrix.IsIndecomposable}.~\qed\end{proof}
\end{proof}
\step{rcreverse}{\assume{$M$ has a nontrivial upper block triangular decomposition.}
\prove{Its support is disconnected.}}
\begin{proof}
\step{rcback}{Both cross blocks vanish.}
\begin{proof}\pf\ The lower block vanishes by hypothesis and the upper
one vanishes by \asm[2]{0}.~\qed\end{proof}
\step{rcpath}{Every path starting in the first block stays in that block.}
\begin{proof}\pf\ Induct on path length using \stepref{rcback}.~\qed\end{proof}
\qedstep
\begin{proof}\pf\ Choose one vertex in each positive-size block and apply
\stepref{rcpath}.~\qed\end{proof}
\end{proof}
\qedstep
\begin{proof}\pf\ Contraposition of \stepref{rcforward} and
\stepref{rcreverse} gives the equivalence. ~\qed\end{proof}
\end{proof}

\begin{namedresult}{Definition}{(weighted graphs)}\label{def:rem-weighted}
\lean{Matrix.maxRayleigh}\lean{WeightedGraph}\lean{WeightedGraph.Contains}
\uses{facts:rem-library}
\assume{$I$ is finite and nonempty.}
A weighted graph on $I$ has a real symmetric matrix $B$, zero diagonal,
and nonnegative entries. Its support joins $i,j$ when $B(i,j)>0$.
For any real symmetric matrix on $I$, write
$\lambda(B)=\max_{\rsq\Zx=1}\Zx^TB\Zx$.
For weighted graphs $Y,X$, write $Y\preccurlyeq X$ if an injection $e$
satisfies $B_Y(i,j)\leq B_X(e(i),e(j))$ for every $i,j$.
This specified containment is proper when $e$ is not onto or one inequality
is strict. Edges may be omitted and weights decreased.
\end{namedresult}

\begin{namedresult}{Lemma}{(edge-constant functions)}\label{lem:rem-edge-constant}
\lean{SimpleGraph.Connected.constant_of_edge_eq}
\uses{facts:rem-library}
\assume{\begin{pfenum}
\item $X$ is a connected simple graph on a vertex set $I$.
\item $S$ is any set and $f:I\to S$ satisfies $f(u)=f(v)$ whenever $u\sim v$.
\end{pfenum}}
\prove{For every $u,v\in I$, $f(u)=f(v)$. In particular there is
$c\in S$ with $f(v)=c$ for every $v$.}
\end{namedresult}
\begin{proof}
\pfsketch\ Induct on the walk supplied by connectedness.
\step{recwalk}{For every walk from $u$ to $v$, $f(u)=f(v)$.}
\begin{proof}
\step{recnil}{A walk of length zero has equal endpoint values.}
\begin{proof}\pf\ Its endpoints coincide by the definition of a walk.~\qed\end{proof}
\step{reccons}{An edge $u\sim w$ followed by a walk from $w$ to $v$
with $f(w)=f(v)$ gives $f(u)=f(v)$.}
\begin{proof}\pf\ Compose the edge equality in \asm[2]{0} with
the equality for the remaining walk.~\qed\end{proof}
\qedstep
\begin{proof}\pf\ Walk induction using \stepref{recnil},
\stepref{reccons}.~\qed\end{proof}
\end{proof}
\step{recpair}{For every $u,v$, $f(u)=f(v)$.}
\begin{proof}\pf\ Connectedness in \asm[1]{0} supplies a walk;
apply \stepref{recwalk}.~\qed\end{proof}
\qedstep
\begin{proof}\pf\ Choose $u$ by nonemptiness in connectedness and set
$c=f(u)$ in \stepref{recpair}.~\qed\end{proof}
\end{proof}

\begin{namedresult}{Lemma}{(propagation on a connected support)}\label{lem:rem-propagation}
\lean{WeightedGraph.eq_zero_of_nonneg_eigenvector_zero}
\uses{def:rem-weighted,lem:rem-edge-constant}
\assume{\begin{pfenum}
\item $B$ is a nonnegative real symmetric matrix on a finite nonempty set.
 The simple graph $i\sim j\iff i\ne j\text{ and }B(i,j)>0$ is connected.
\item $B\Zx=t\Zx$, all $\Zx(i)\geq0$, and $\Zx(i_0)=0$.
\end{pfenum}}
\prove{$\Zx=0$.}
\end{namedresult}
\begin{proof}
\pfsketch\ Show that having value zero is constant across every edge.
\step{rpneighbour}{\assume{$\Zx(i)=0$ and $B(i,j)>0$.}
\prove{$\Zx(j)=0$.}}
\begin{proof}
\step{rpsum}{$\sum_k B(i,k)\Zx(k)=0$ with every summand nonnegative.}
\begin{proof}\pf\ Evaluate \asm[2]{0} at $i$.~\qed\end{proof}
\step{rpterm}{$B(i,j)\Zx(j)=0$.}
\begin{proof}\pf\ A nonnegative summand of the zero sum
\stepref{rpsum} is zero. ~\qed\end{proof}
\qedstep
\begin{proof}\pf\ Cancel $B(i,j)>0$ in \stepref{rpterm}.~\qed\end{proof}
\end{proof}
\step{rppath}{Every vertex has value zero.}
\begin{proof}\pf\ By \stepref{rpneighbour} in both directions,
the indicator of $\{i:\Zx(i)=0\}$ agrees across every support edge.
Lemma~\ref{lem:rem-edge-constant} makes it constant; its value at $i_0$
is 1 by \asm[2]{0}.~\qed\end{proof}
\qedstep
\begin{proof}\pf\ Function extensionality applied to \stepref{rppath}.~\qed\end{proof}
\end{proof}

\begin{namedresult}{Lemma}{(Rayleigh quotients)}\label{lem:rayleigh}
\lean{Matrix.IsHermitian.maxRayleigh_eigenvalue}\lean{Matrix.IsHermitian.rayleigh_eq_iff}
\lean{Matrix.IsHermitian.quadratic_le_maxRayleigh}
\lean{Matrix.IsHermitian.exists_nonneg_max_eigenvector}\lean{Matrix.IsHermitian.max_eigenspace_simple}
\uses{def:rem-weighted,lem:rem-null,lem:rem-propagation,facts:rem-library}
\assume{$B$ is real symmetric on a finite nonempty set $I$.}
\prove{For every real vector $\Zy$,
$\Zy^TB\Zy\leq\lambda(B)\rsq\Zy$.
$\lambda(B)$ is an eigenvalue bounding every real eigenvalue.
A nonzero vector attains its Rayleigh quotient exactly when it is a
$\lambda(B)$-eigenvector. If $B\geq0$ there is a nonnegative unit maximizer.
If its support is also connected, this vector is strictly positive and
the $\lambda(B)$-eigenspace is one-dimensional.}
\end{namedresult}
\begin{proof}
\pfsketch\ Use a nonnegative form at the maximum, then absolute values
and the propagation lemma. The nonnegative conclusion does not require
connectedness and is stated for use with a disconnected smaller graph.
\step{rymax}{\pick\ a unit vector $\Zx$ attaining $\lambda=\lambda(B)$.}
\begin{proof}\pf\ Apply the compact maximum facts of
\ref{facts:rem-library} to the continuous quadratic form. ~\qed\end{proof}
\step{rybound}{For every $\Zy$, $\Zy^TB\Zy\leq\lambda\rsq\Zy$.}
\begin{proof}
\step{ryzero}{\case{$\Zy=0$.}}
\begin{proof}\pf\ Both sides vanish by substitution. ~\qed\end{proof}
\step{rynonzero}{\case{$\Zy\ne0$.}}
\begin{proof}\pf\ Apply \stepref{rymax} to $\Zy/\|\Zy\|$ and multiply
by $\|\Zy\|^2>0$. ~\qed\end{proof}
\qedstep
\begin{proof}\pf\ The cases \stepref{ryzero}, \stepref{rynonzero} exhaust
all vectors. ~\qed\end{proof}
\end{proof}
\step{ryattain}{For nonzero $\Zy$, equality in \stepref{rybound} holds
 if and only if $B\Zy=\lambda\Zy$.}
\begin{proof}
\step{rykernel}{Equality implies $(\lambda I-B)\Zy=0$.}
\begin{proof}\pf\ Apply Lemma~\ref{lem:rem-null} to the nonnegative form
\stepref{rybound}.~\qed\end{proof}
\step{rysub}{An eigenvector satisfies $\Zy^TB\Zy=\lambda\rsq\Zy$.}
\begin{proof}\pf\ Substitute the eigenvector equation. ~\qed\end{proof}
\qedstep
\begin{proof}\pf\ Combine \stepref{rykernel}, \stepref{rysub}.~\qed\end{proof}
\end{proof}
\step{ryeigen}{$\lambda$ is an eigenvalue and every real eigenvalue is
at most $\lambda$.}
\begin{proof}
\step{ryexists}{$B\Zx=\lambda\Zx$ with $\Zx\ne0$.}
\begin{proof}\pf\ Apply \stepref{ryattain} to \stepref{rymax}.~\qed\end{proof}
\step{ryother}{\assume{$B\Zv=t\Zv$ and $\Zv\ne0$.}
\prove{$t\leq\lambda$.}}
\begin{proof}\pf\ Substitute in \stepref{rybound} and cancel
$\rsq\Zv>0$.~\qed\end{proof}
\qedstep
\begin{proof}\pf\ Combine \stepref{ryexists}, \stepref{ryother}.~\qed\end{proof}
\end{proof}
\step{ryabs}{\assume{$B\geq0$.}
\prove{$|\Zx|$ is a nonnegative unit maximizer.}}
\begin{proof}
\step{ryabsnorm}{$\rsq{|\Zx|}=1$.}
\begin{proof}\pf\ Use $|x|^2=x^2$ in the coordinate sum of
\stepref{rymax}.~\qed\end{proof}
\step{ryabsform}{$|\Zx|^TB|\Zx|\geq\Zx^TB\Zx$.}
\begin{proof}\pf\ Sum $B(i,j)x(i)x(j)\leq B(i,j)|x(i)||x(j)|$
over ordered pairs, using $B(i,j)\geq0$.~\qed\end{proof}
\qedstep
\begin{proof}\pf\ Sandwich \stepref{ryabsform} between the maximum in
\stepref{rymax} and the bound \stepref{rybound}, using \stepref{ryabsnorm}.~\qed\end{proof}
\end{proof}
\step{rypositive}{\assume{$B\geq0$, its support is connected, and $\Zw\geq0$ is a nonzero maximizing eigenvector.}
\prove{$\Zw$ is positive everywhere.}}
\begin{proof}\pf\ A zero coordinate would force $\Zw=0$ by
Lemma~\ref{lem:rem-propagation}.~\qed\end{proof}
\step{rysimple}{\assume{$B\geq0$ has connected support, $\Zw>0$, and $B\Zw=\lambda\Zw$.}
\prove{Every $\lambda$-eigenvector is a multiple of $\Zw$.}}
\begin{proof}
\pflet{For a $\lambda$-eigenvector $\Zv$, let
$t=\max_i\Zv(i)/\Zw(i)$.}
\step{ryresidual}{$t\Zw-\Zv$ is nonnegative and has a zero coordinate.}
\begin{proof}\pf\ Use the attained finite maximum defining $t$ and
\stepref{rypositive}.~\qed\end{proof}
\step{ryresideigen}{$B(t\Zw-\Zv)=\lambda(t\Zw-\Zv)$.}
\begin{proof}\pf\ Linearity and the two eigenvector equations. ~\qed\end{proof}
\qedstep
\begin{proof}\pf\ Lemma~\ref{lem:rem-propagation} applied to
\stepref{ryresidual}, \stepref{ryresideigen} gives $\Zv=t\Zw$.~\qed\end{proof}
\end{proof}
\qedstep
\begin{proof}\pf\ Combine \stepref{ryattain}, \stepref{ryeigen},
\stepref{ryabs}, \stepref{rypositive}, \stepref{rysimple}.~\qed\end{proof}
\end{proof}

\begin{namedresult}{Lemma}{(Rayleigh quotients under reindexing)}\label{lem:rem-rayleigh-reindex}
\lean{Matrix.IsHermitian.maxRayleigh_reindex}
\uses{def:rem-weighted}
\assume{$B_I,B_J$ are real symmetric matrices on finite nonempty sets
$I,J$, and $e:I\to J$ is a bijection with
$B_J(e(i),e(j))=B_I(i,j)$ for all $i,j$.}
\prove{$\lambda(B_I)=\lambda(B_J)$. In particular isomorphic weighted
graphs have equal largest eigenvalues.}
\end{namedresult}
\begin{proof}
\pfsketch\ Reindex the unit sphere and the quadratic form.
\step{rrinorm}{The bijection $\Zx\mapsto\Zx\circ e$ preserves squared norms.}
\begin{proof}\pf\ Reindex the finite sum of squares along $e$.~\qed\end{proof}
\step{rriform}{$(\Zx\circ e)^TB_I(\Zx\circ e)=\Zx^TB_J\Zx$.}
\begin{proof}\pf\ Reindex the double sum along $e\times e$ using
the entry equality in \asm{0}.~\qed\end{proof}
\qedstep
\begin{proof}\pf\ The sets of values maximized in
Definition~\ref{def:rem-weighted} coincide by \stepref{rrinorm},
\stepref{rriform}.~\qed\end{proof}
\end{proof}

\begin{namedresult}{Lemma}{(strict monotonicity)}\label{lem:monotonicity}
\lean{WeightedGraph.maxRayleigh_mono}\lean{WeightedGraph.maxRayleigh_strictMono}
\lean{Matrix.IsHermitian.eigenvalue_eq_maxRayleigh_of_pos}
\uses{def:rem-weighted,lem:rayleigh,lem:rem-propagation,facts:inner}
\assume{\begin{pfenum}
\item $X$ is a finite connected weighted graph with nonnegative weights.
\item $Y$ is a finite nonempty weighted graph with nonnegative weights,
 and $e$ witnesses $Y\preccurlyeq X$.
\end{pfenum}}
\prove{$\lambda(Y)\leq\lambda(X)$, strictly for proper containment.
For any connected nonnegative real symmetric matrix, a strictly positive
eigenvector has eigenvalue equal to its largest eigenvalue.}
\end{namedresult}
\begin{proof}
\pfsketch\ Extend a nonnegative maximizer by zero; equality precludes
both missing vertices and increased entries.
\step{rmchoose}{\pick\ a nonnegative unit maximizer $\Zv$ of $Y$.}
\begin{proof}\pf\ Apply Lemma~\ref{lem:rayleigh} to $B_Y$.~\qed\end{proof}
\pflet{$\Zx(e(i))=\Zv(i)$ and $\Zx(j)=0$ off the image of $e$.}
\step{rmnorm}{$\rsq\Zx=1$.}
\begin{proof}\pf\ Reindex the sum of squares by the injection $e$ in
\asm[2]{0}, using \stepref{rmchoose}.~\qed\end{proof}
\step{rmcompare}{$\lambda(Y)\leq\Zx^TB_X\Zx\leq\lambda(X)$.}
\begin{proof}
\step{rmlower}{$\lambda(Y)\leq\Zx^TB_X\Zx$.}
\begin{proof}\pf\ Sum the entrywise inequalities in \asm[2]{0}
multiplied by the nonnegative $\Zv(i)\Zv(j)$ from \stepref{rmchoose}.~\qed\end{proof}
\step{rmupper}{$\Zx^TB_X\Zx\leq\lambda(X)$.}
\begin{proof}\pf\ Apply the definition of the maximum to
\stepref{rmnorm}.~\qed\end{proof}
\qedstep
\begin{proof}\pf\ Combine \stepref{rmlower}, \stepref{rmupper}.~\qed\end{proof}
\end{proof}
\step{rmequal}{\assume{$\lambda(Y)=\lambda(X)$.}
\prove{The specified containment is not proper.}}
\begin{proof}
\step{rmmax}{$\Zx$ is a maximizing eigenvector of $X$.}
\begin{proof}\pf\ Apply Lemma~\ref{lem:rayleigh} to equality in
\stepref{rmcompare}.~\qed\end{proof}
\step{rmpos}{$\Zx(j)>0$ for every vertex $j$ of $X$.}
\begin{proof}\pf\ Apply Lemma~\ref{lem:rem-propagation} to
\stepref{rmmax}, \stepref{rmnorm} and nonnegativity of $\Zx$.~\qed\end{proof}
\step{rmsurj}{$e$ is onto.}
\begin{proof}\pf\ A vertex outside the image has $\Zx(j)=0$ by definition,
contradicting \stepref{rmpos}.~\qed\end{proof}
\step{rmentries}{All entry inequalities are equalities.}
\begin{proof}\pf\ In the zero sum
$\sum_{i,j}(B_X(e(i),e(j))-B_Y(i,j))\Zv(i)\Zv(j)=0$
from \stepref{rmcompare}, every summand is nonnegative and both vector
factors are positive by \stepref{rmpos}.~\qed\end{proof}
\qedstep
\begin{proof}\pf\ Use the definition of properness with
\stepref{rmsurj}, \stepref{rmentries}.~\qed\end{proof}
\end{proof}
\step{rmposeigen}{\assume{$B$ is real symmetric, nonnegative, and has connected off-diagonal positive support; $B\Zv=t\Zv$ and $\Zv>0$.}
\prove{$t=\lambda(B)$.}}
\begin{proof}
\step{rmw}{\pick\ $\Zw>0$ with $B\Zw=\lambda(B)\Zw$.}
\begin{proof}\pf\ Apply Lemma~\ref{lem:rayleigh} to this matrix $B$.~\qed\end{proof}
\step{rmorth}{$(t-\lambda(B))(\Zv,\Zw)=0$.}
\begin{proof}\pf\ Use symmetry of $B$ in $(B\Zv,\Zw)=(\Zv,B\Zw)$
and \stepref{rmw}.~\qed\end{proof}
\qedstep
\begin{proof}\pf\ Cancel the positive sum $(\Zv,\Zw)$ in
\stepref{rmorth}.~\qed\end{proof}
\end{proof}
\qedstep
\begin{proof}\pf\ Combine \stepref{rmcompare}, the contrapositive of
\stepref{rmequal}, and \stepref{rmposeigen}.~\qed\end{proof}
\end{proof}

\begin{namedresult}{Lemma}{(symmetrization)}\label{lem:symmetrization}
\lean{Matrix.IsFiniteCartan.symmetrization}
\uses{def:rem-bases,def:rem-weighted,lem:integer-positive,lem:connected,lem:rayleigh,facts:rem-library}
\assume{\begin{pfenum}
\item $I$ is finite and nonempty, $M:I\times I\to\mathbb Z$ is finite
 Cartan and indecomposable.
\item $d:I\to\mathbb Z_{>0}$ symmetrizes $M$ as in its finite Cartan
 hypothesis, and $D=\operatorname{diag}(d)$.
\end{pfenum}}
\prove{$S=D^{1/2}MD^{-1/2}$ is symmetric positive definite and equals
$2I-B$, where $B(i,j)=\sqrt{M(i,j)M(j,i)}$ off the diagonal and $B(i,i)=0$.
The weighted graph of $B$ is connected, has edge weights
$1,\sqrt2,\sqrt3$, and $\lambda(B)<2$.}
\end{namedresult}
\begin{proof}
\pfsketch\ Congruence transfers positivity; a two-coordinate test bounds
the integer edge products.
\step{rsreal}{$DM$ is positive definite over $\mathbb R$.}
\begin{proof}\pf\ Apply Lemma~\ref{lem:integer-positive} and the
quadratic-form interface in \ref{facts:rem-library}.~\qed\end{proof}
\step{rscong}{$S=D^{-1/2}(DM)D^{-1/2}$ is symmetric positive definite.}
\begin{proof}\pf\ Substitute the invertible diagonal change of coordinates
in the form of \stepref{rsreal}.~\qed\end{proof}
\step{rsdiag}{$S(i,i)=2$ for every $i$.}
\begin{proof}\pf\ The diagonal factors cancel and $M(i,i)=2$ by
\asm[1]{0}.~\qed\end{proof}
\step{rsentry}{For $i\ne j$, $S(i,j)\leq0$ and
$S(i,j)^2=M(i,j)M(j,i)$.}
\begin{proof}\pf\ Substitute $S(i,j)=\sqrt{d(i)/d(j)}M(i,j)$ and
$d(i)M(i,j)=d(j)M(j,i)$ from \asm[2]{0}.~\qed\end{proof}
\step{rsbound}{For $i\ne j$, $S(i,j)^2<4$.}
\begin{proof}
\pflet{$b=S(i,j)$ and $\Zx=\Zu_i-(b/2)\Zu_j\ne0$.}
\step{rstest}{$\Zx^TS\Zx=2-b^2/2$.}
\begin{proof}\pf\ Expand the two coordinates using \stepref{rsdiag}
and symmetry in \stepref{rscong}.~\qed\end{proof}
\qedstep
\begin{proof}\pf\ Positivity in \stepref{rscong} and \stepref{rstest}
give $2-b^2/2>0$.~\qed\end{proof}
\end{proof}
\step{rsweights}{$M(i,j)M(j,i)\in\{0,1,2,3\}$.}
\begin{proof}\pf\ Apply integer interval arithmetic to
\stepref{rsentry}, \stepref{rsbound}.~\qed\end{proof}
\step{rsgraph}{$B=2I-S$ has precisely the required entries and weights.}
\begin{proof}\pf\ The nonpositive square root is unique in
\stepref{rsentry}; use \stepref{rsdiag}, \stepref{rsweights}.~\qed\end{proof}
\step{rsconn}{Its support is connected.}
\begin{proof}\pf\ Apply Lemma~\ref{lem:connected} to the identical
nonzero support in \stepref{rsgraph}.~\qed\end{proof}
\step{rslambda}{$\lambda(B)<2$.}
\begin{proof}
\step{rsmax}{\pick\ a unit maximizer $\Zv$ for $B$.}
\begin{proof}\pf\ Apply Lemma~\ref{lem:rayleigh}.~\qed\end{proof}
\qedstep
\begin{proof}\pf\ At \stepref{rsmax},
$\lambda(B)=2-\Zv^TS\Zv<2$ by \stepref{rscong}, \stepref{rsgraph}.~\qed\end{proof}
\end{proof}
\qedstep
\begin{proof}\pf\ Combine \stepref{rscong}, \stepref{rsgraph},
\stepref{rsconn}, \stepref{rslambda}.~\qed\end{proof}
\end{proof}

\subsubsection{Affine obstructions and labelled trees}

\begin{namedresult}{Definition}{(affine obstructions)}\label{def:affine}
\lean{WeightedGraph.IsAffineObstruction}
\lean{SimpleGraph.smithA}\lean{SimpleGraph.smithD}
\lean{SimpleGraph.smithE₆}\lean{SimpleGraph.smithE₇}\lean{SimpleGraph.smithE₈}
\uses{def:rem-weighted}
\assume{All lengths below are integers in the displayed ranges.}
The tree $T(a,b,c)$, $a,b,c\geq1$, has vertex set
\[
 V_{a,b,c}=\{*\}\sqcup(\{1\}\times[1,a])\sqcup
 (\{2\}\times[1,b])\sqcup(\{3\}\times[1,c]).
\]
Its only edges are $*\sim(r,1)$ and consecutive positions on an arm.
The affine obstructions are the following weighted graphs:
\begin{enumerate}
\item Unit cycles $\Cyc_m$, $m\geq3$.
\item $\widetilde{\typ D}_4=K_{1,4}$; and $\widetilde{\typ D}_\ell$,
$\ell\geq5$, a unit path of $\ell-4$ edges with two new leaves at each end.
\item $\widetilde{\typ E}_6=T(2,2,2)$,
$\widetilde{\typ E}_7=T(1,3,3)$, $\widetilde{\typ E}_8=T(1,2,5)$,
all with unit weights.
\item $X^{\rm end}_m$, $m\geq3$: a path on $m$ vertices with first and
last weights $\sqrt2$ and all intervening weights 1.
\item $X^{\rm fork}_k$, $k\geq1$: a path $v_0,\ldots,v_k$ with two
unit leaves at $v_0$, last path weight $\sqrt2$, and other path weights 1.
\item $X^{\rm mid}$: the path with weights $1,1,\sqrt2,1$.
\item $X^{\rm triple}$: the path with weights $\sqrt3,1$.
\end{enumerate}
Thus both edges of $X^{\rm end}_3$ have weight $\sqrt2$, and the weighted
edge of $X^{\rm fork}_1$ meets the branch vertex. The project model
\texttt{smithA} at parameter $\ell\geq2$ means $\Cyc_{\ell+1}$.
\end{namedresult}

\begin{namedresult}{Lemma}{(connected labelled graphs)}\label{lem:rem-labelled-connected}
\lean{WeightedGraph.connected_of_path_or_arm}
\lean{WeightedGraph.IsAffineObstruction.connected}
\uses{def:affine,def:rem-weighted}
\assume{$X$ is a finite weighted graph whose support admits a nonempty
path labelling, an arm labelling $T(a,b,c)$ with $a,b,c\geq1$, or is the
support of an obstruction in Definition~\ref{def:affine}.}
\prove{The support of $X$ is connected.}
\end{namedresult}
\begin{proof}
\pfsketch\ Give a walk from every vertex to one chosen vertex.
\step{rconnpath}{\case{The support has a path labelling.}}
\begin{proof}\pf\ Consecutive positions $i,i-1,\ldots,0$ give a walk
from every vertex to position 0, including the singleton path.~\qed\end{proof}
\step{rconnarm}{\case{The support has an arm labelling.}}
\begin{proof}\pf\ From $(r,t)$ follow $(r,t-1),\ldots,(r,1),*$;
the centre has the length-zero walk to itself.~\qed\end{proof}
\step{rconncycle}{\case{$X$ has the support of $\Cyc_m$.}}
\begin{proof}\pf\ Follow consecutive cycle positions down to 0;
each is an edge by Definition~\ref{def:affine}.~\qed\end{proof}
\step{rconnstar}{\case{$X$ has the support of $\widetilde{\typ D}_4$.}}
\begin{proof}\pf\ Every leaf has an edge to the centre of $K_{1,4}$.~\qed\end{proof}
\step{rconnD}{\case{$X$ has the support of $\widetilde{\typ D}_\ell$,
$\ell\geq5$.}}
\begin{proof}\pf\ A leaf first reaches its attached endpoint; from
any path position, consecutive path edges reach the first endpoint.~\qed\end{proof}
\step{rconnfork}{\case{$X$ has the support of $X^{\rm fork}_k$.}}
\begin{proof}\pf\ Its extra leaves reach $v_0$ in one edge; $v_i$
reaches it along $v_i,v_{i-1},\ldots,v_0$.~\qed\end{proof}
\qedstep
\begin{proof}\pf\ The $\widetilde{\typ E}$ supports are arm trees;
$X^{\rm end},X^{\rm mid},X^{\rm triple}$ have path supports.
Thus \stepref{rconnpath}--\stepref{rconnfork} exhaust
Definition~\ref{def:affine} and the two other hypotheses.~\qed\end{proof}
\end{proof}

\begin{namedresult}{Lemma}{(unit containment and support embeddings)}\label{lem:rem-unit-containment}
\lean{WeightedGraph.unit_contains_iff_support_embedding}
\uses{def:rem-weighted}
\assume{$X$ is a finite weighted graph with $B_X(u,v)\geq1$ on every
support edge, and $Y$ is a finite nonempty simple graph, given unit edge weights.}
\prove{$Y\preccurlyeq X$ if and only if there is an injection
$e:V(Y)\to V(X)$ sending every edge of $Y$ to a support edge of $X$.
The same injection witnesses either assertion; nonedges need not be preserved.}
\end{namedresult}
\begin{proof}
\pfsketch\ Check the containment inequality on edges and nonedges.
\step{rucforward}{\assume{$e$ witnesses $Y\preccurlyeq X$.}
\prove{$e$ preserves edges into the support.}}
\begin{proof}\pf\ On an edge the inequality is
$1\leq B_X(e(u),e(v))$, so that entry is positive.~\qed\end{proof}
\step{rucreverse}{\assume{$e$ is an injective edge-preserving map to the support.}
\prove{$e$ witnesses $Y\preccurlyeq X$.}}
\begin{proof}
\step{rucedge}{On an edge, $B_Y(u,v)=1\leq B_X(e(u),e(v))$.}
\begin{proof}\pf\ Apply the weight lower bound in \asm{0} to
the support edge supplied by \asm{1}.~\qed\end{proof}
\step{rucnonedge}{On a nonedge, $B_Y(u,v)=0\leq B_X(e(u),e(v))$.}
\begin{proof}\pf\ Apply nonnegativity from
Definition~\ref{def:rem-weighted}.~\qed\end{proof}
\qedstep
\begin{proof}\pf\ The cases \stepref{rucedge}, \stepref{rucnonedge}
give every inequality of containment.~\qed\end{proof}
\end{proof}
\qedstep
\begin{proof}\pf\ Combine \stepref{rucforward}, \stepref{rucreverse}.~\qed\end{proof}
\end{proof}

\begin{namedresult}{Lemma}{(a cycle gives a cycle subgraph)}\label{lem:rem-cycle-embedding}
\lean{SimpleGraph.Walk.IsCycle.exists_cycle_embedding}
\uses{def:affine,facts:rem-library}
\assume{$Z$ is a simple graph and $p$ is a closed walk of length $m$
which is a cycle.}
\prove{$m\geq3$, and the map sending $i\in[0,m-1]$ to the $i$th
vertex of $p$ is an injection preserving every edge of $\Cyc_m$.}
\end{namedresult}
\begin{proof}
\pfsketch\ Omit the repeated last vertex from the cycle walk.
\step{rcylength}{$m\geq3$ and the first $m$ vertices of $p$ are distinct.}
\begin{proof}
\step{rcythree}{$m\geq3$.}
\begin{proof}\raggedright\pf\ Apply
\path{SimpleGraph.Walk.isCycle_iff_isPath_tail_and_le_length}
in \path{Mathlib/Combinatorics/SimpleGraph/Paths.lean}.~\qed\end{proof}
\step{rcyinjective}{The first $m$ vertices are distinct.}
\begin{proof}\pf\ Apply \path{SimpleGraph.Walk.IsCycle.getVert_injOn'}
in that file to the positions $0,\ldots,m-1$.~\qed\end{proof}
\qedstep
\begin{proof}\pf\ Combine \stepref{rcythree}, \stepref{rcyinjective}.~\qed\end{proof}
\end{proof}
\step{rcyedges}{Successive positions and the pair $m-1,0$ map to edges.}
\begin{proof}\pf\ These are the walk edges of $p$, including its
last edge back to its start.~\qed\end{proof}
\qedstep
\begin{proof}\pf\ Apply \stepref{rcylength}, \stepref{rcyedges}
to the cycle definition in Definition~\ref{def:affine}.~\qed\end{proof}
\end{proof}

\begin{namedresult}{Lemma}{(affine eigenvectors)}\label{lem:affine}
\lean{WeightedGraph.IsAffineObstruction.maxRayleigh_eq_two}
\uses{def:affine,lem:monotonicity,lem:rem-labelled-connected}
\assume{$X$ is one of the graphs in Definition~\ref{def:affine}.}
\prove{$\lambda(X)=2$. The following positive vectors $\Zw$ satisfy
$B\Zw=2\Zw$: 1 on a cycle; 2 on each nonleaf and 1 on each leaf of
$\widetilde{\typ D}$; centre $h=3,4,6$ respectively on the three
$\widetilde{\typ E}$ trees and value $h(a+1-t)/(a+1)$ at position $t$
of an arm of length $a$; endpoints 1 and interiors $\sqrt2$ for
$X^{\rm end}$; values $2,\ldots,2,\sqrt2$ on the path and 1 on the
two extra leaves for $X^{\rm fork}$; $(1,2,3,2\sqrt2,\sqrt2)$ for
$X^{\rm mid}$; $(\sqrt3,2,1)$ for $X^{\rm triple}$.}
\end{namedresult}
\begin{proof}
\pfsketch\ Check each vertex type. This proves the unbounded families,
including their shortest instances.
\pflet{$\Zw$ is the vector specified in the statement for the graph at hand.}
\step{rafcycle}{On a cycle, each weighted neighbour sum is 2.}
\begin{proof}\pf\ Each vertex has two unit neighbours by
Definition~\ref{def:affine}.~\qed\end{proof}
\step{rafD}{The displayed vector solves the equations on
$\widetilde{\typ D}$.}
\begin{proof}
\step{rafDleaf}{At a leaf the sum is $2=2\cdot1$.}
\begin{proof}\pf\ Its unique neighbour has value 2 by the stated vector. ~\qed\end{proof}
\step{rafDcentre}{At the centre of $K_{1,4}$ the sum is $4=2\cdot2$.}
\begin{proof}\pf\ Add its four leaf values from the definition of $\Zw$.~\qed\end{proof}
\step{rafDfork}{At a fork for $\ell\geq5$ the sum is $2+1+1=2\cdot2$.}
\begin{proof}\pf\ Use its one path neighbour and two leaves in
Definition~\ref{def:affine}.~\qed\end{proof}
\step{rafDinternal}{At a path interior the sum is $2+2=2\cdot2$.}
\begin{proof}\pf\ Use its two path neighbours in Definition~\ref{def:affine}.~\qed\end{proof}
\qedstep
\begin{proof}\pf\ The leaf, centre, fork and interior cases
\stepref{rafDleaf}--\stepref{rafDinternal} exhaust the vertices. ~\qed\end{proof}
\end{proof}
\step{rafE}{The displayed vectors solve the equations on
$\widetilde{\typ E}_6,\widetilde{\typ E}_7,\widetilde{\typ E}_8$.}
\begin{proof}
\step{rafEarm}{At a noncentral vertex the sum is twice its value.}
\begin{proof}\pf\ The arm formula is an arithmetic progression with
value $h$ at position 0 and virtual value 0 at position $a+1$.~\qed\end{proof}
\step{rafEcentre}{For arms $a,b,c$ and centre value $h>0$, the central
eigenvector equation is equivalent to
$1/(a+1)+1/(b+1)+1/(c+1)=1$.}
\begin{proof}\pf\ Its neighbour sum is
$h(a/(a+1)+b/(b+1)+c/(c+1))$; cancel $h$ from equality to $2h$.~\qed\end{proof}
\step{rafE6}{The central equation holds for $(2,2,2)$.}
\begin{proof}\pf\ Substitute $1/3+1/3+1/3=1$ in
\stepref{rafEcentre}.~\qed\end{proof}
\step{rafE7}{The central equation holds for $(1,3,3)$.}
\begin{proof}\pf\ Substitute $1/2+1/4+1/4=1$ in \stepref{rafEcentre}.~\qed\end{proof}
\step{rafE8}{The central equation holds for $(1,2,5)$.}
\begin{proof}\pf\ Substitute $1/2+1/3+1/6=1$ in \stepref{rafEcentre}.~\qed\end{proof}
\qedstep
\begin{proof}\pf\ Combine \stepref{rafEarm} with the three central
cases \stepref{rafE6}--\stepref{rafE8}.~\qed\end{proof}
\end{proof}
\step{rafend}{The equations hold on $X^{\rm end}_m$.}
\begin{proof}
\step{rafendleaf}{At either endpoint the sum is $\sqrt2\sqrt2=2$.}
\begin{proof}\pf\ Substitute the endpoint edge and neighbouring value. ~\qed\end{proof}
\step{rafendinternal}{At an interior vertex each of its two sides
contributes $\sqrt2$.}
\begin{proof}
\step{rafendboundary}{A neighbouring endpoint contributes
$\sqrt2\cdot1=\sqrt2$.}
\begin{proof}\pf\ Substitute the endpoint value and edge weight. ~\qed\end{proof}
\step{rafendmiddle}{A neighbouring interior contributes
$1\cdot\sqrt2=\sqrt2$.}
\begin{proof}\pf\ Substitute the interior value and edge weight. ~\qed\end{proof}
\qedstep
\begin{proof}\pf\ Every neighbour is an endpoint or an interior vertex,
so \stepref{rafendboundary}, \stepref{rafendmiddle} apply to each side. ~\qed\end{proof}
\end{proof}
\qedstep
\begin{proof}\pf\ Add the two contributions in
\stepref{rafendinternal}, together with \stepref{rafendleaf}.~\qed\end{proof}
\end{proof}
\step{raffork}{The equations hold on $X^{\rm fork}_k$.}
\begin{proof}
\step{rafforkleaf}{An extra leaf has sum $2=2\cdot1$.}
\begin{proof}\pf\ Its neighbour is $v_0$ with value 2. ~\qed\end{proof}
\step{rafforklast}{At $v_k$ the sum is $\sqrt2\cdot2=2\sqrt2$.}
\begin{proof}\pf\ Substitute its unique edge. ~\qed\end{proof}
\step{rafforksmall}{For $k=1$, the sum at $v_0$ is
$1+1+\sqrt2\sqrt2=4$.}
\begin{proof}\pf\ Substitute all three incident edges from
Definition~\ref{def:affine}.~\qed\end{proof}
\step{rafforklarge}{For $k>1$, the sum at $v_0$ is $1+1+2=4$.}
\begin{proof}\pf\ The edge $v_0v_1$ has weight 1 and $v_1$ has value 2. ~\qed\end{proof}
\step{rafforkpenultimate}{For $k>1$, the sum at $v_{k-1}$ is
$2+\sqrt2\sqrt2=4$.}
\begin{proof}\pf\ Substitute the two incident path edges. ~\qed\end{proof}
\step{rafforkinterior}{At $v_i$ with $0<i<k-1$, the sum is $2+2=4$.}
\begin{proof}\pf\ Both neighbours have value 2 and both edges weight 1. ~\qed\end{proof}
\qedstep
\begin{proof}\pf\ The six cases \stepref{rafforkleaf}--\stepref{rafforkinterior}
exhaust the vertices for $k\geq1$.~\qed\end{proof}
\end{proof}
\step{raffinite}{The equations hold on $X^{\rm mid},X^{\rm triple}$.}
\begin{proof}
\step{rafmid0}{At position 0 of $X^{\rm mid}$ the sum is $2=2\cdot1$.}
\begin{proof}\pf\ Substitute its sole neighbour. ~\qed\end{proof}
\step{rafmid1}{At position 1 the sum is $1+3=2\cdot2$.}
\begin{proof}\pf\ Substitute its two unit edges. ~\qed\end{proof}
\step{rafmid2}{At position 2 the sum is $2+\sqrt2(2\sqrt2)=2\cdot3$.}
\begin{proof}\pf\ Substitute its unit and double-weight edges. ~\qed\end{proof}
\step{rafmid3}{At position 3 the sum is $3\sqrt2+\sqrt2=2(2\sqrt2)$.}
\begin{proof}\pf\ Substitute its double-weight and unit edges. ~\qed\end{proof}
\step{rafmid4}{At position 4 the sum is $2\sqrt2=2\cdot\sqrt2$.}
\begin{proof}\pf\ Substitute its sole neighbour. ~\qed\end{proof}
\step{raftriple0}{At position 0 of $X^{\rm triple}$ the sum is $2\sqrt3$.}
\begin{proof}\pf\ Its neighbour has value 2 and its edge weight is $\sqrt3$.~\qed\end{proof}
\step{raftriple1}{At position 1 the sum is $\sqrt3\sqrt3+1=2\cdot2$.}
\begin{proof}\pf\ Substitute its two neighbours. ~\qed\end{proof}
\step{raftriple2}{At position 2 the sum is $2=2\cdot1$.}
\begin{proof}\pf\ Substitute its sole neighbour. ~\qed\end{proof}
\qedstep
\begin{proof}\pf\ Positions \stepref{rafmid0}--\stepref{raftriple2}
exhaust both vertex sets. ~\qed\end{proof}
\end{proof}
\step{rafpositive}{Every displayed entry is positive.}
\begin{proof}\pf\ The constants are positive; on an arm $1\leq t\leq a$
gives $a+1-t\geq1$ in the displayed formula. ~\qed\end{proof}
\qedstep
\begin{proof}\pf\ Apply the positive-eigenvector assertion of
Lemma~\ref{lem:monotonicity} to \stepref{rafcycle}--\stepref{raffinite}
and \stepref{rafpositive}, with connectedness from
Lemma~\ref{lem:rem-labelled-connected}.~\qed\end{proof}
\end{proof}

\begin{namedresult}{Lemma}{(path labelling)}\label{lem:path-labelling}
\lean{SimpleGraph.nonempty_iso_pathGraph}\lean{SimpleGraph.exists_path_labelling_from}
\lean{SimpleGraph.path_labelling_reverse}
\uses{facts:rem-library}
\assume{$X$ is a finite nonempty connected acyclic simple graph on $m$
vertices and every degree is at most 2.}
\prove{There is a bijection $e:[0,m-1]\to V(X)$ with
$e(i)\sim e(j)\iff |i-j|=1$. If $w$ is any specified vertex of degree
at most 1, such a bijection can be chosen with $e(0)=w$.
For every path labelling $e$, its reversal $i\mapsto e(m-1-i)$
is also a path labelling.}
\end{namedresult}
\begin{proof}
\pfsketch\ Take a longest simple path; no edge can leave its support.
\step{rlpick}{\pick\ a simple path $v_0,\ldots,v_k$ of maximum length.}
\begin{proof}\pf\ Distinct-vertex lists in a finite set form a finite
set; singleton lists give a nonempty set of paths. ~\qed\end{proof}
\step{rlend0}{No neighbour of $v_0$ lies off this path.}
\begin{proof}\pf\ Prepending that neighbour would extend
\stepref{rlpick} as a simple path. ~\qed\end{proof}
\step{rlendk}{No neighbour of $v_k$ lies off this path.}
\begin{proof}\pf\ Appending that neighbour would extend
\stepref{rlpick} as a simple path. ~\qed\end{proof}
\step{rlinternal}{No neighbour of any $v_i$, $0<i<k$, lies off the path.}
\begin{proof}\pf\ Such a vertex already has the two distinct neighbours
$v_{i-1},v_{i+1}$, so a third violates the degree bound in \asm{0}.~\qed\end{proof}
\step{rlonto}{Every vertex belongs to the path.}
\begin{proof}
\step{rlcross}{A path from a vertex outside the support to $v_0$ would
have an edge with one endpoint inside and one outside.}
\begin{proof}\pf\ Take the first entry into the support in its finite
vertex sequence. ~\qed\end{proof}
\qedstep
\begin{proof}\pf\ Connectedness and \stepref{rlcross} contradict
\stepref{rlend0}, \stepref{rlendk}, \stepref{rlinternal}.~\qed\end{proof}
\end{proof}
\step{rlbijection}{$k+1=m$ and $e(i)=v_i$ is a bijection.}
\begin{proof}\pf\ The path is injective by \stepref{rlpick} and
surjective by \stepref{rlonto}.~\qed\end{proof}
\step{rledges}{The only edges are consecutive pairs.}
\begin{proof}
\step{rlconsecutive}{Every consecutive pair is adjacent.}
\begin{proof}\pf\ Unfold the path in \stepref{rlpick}.~\qed\end{proof}
\step{rlchord}{No two nonconsecutive positions are adjacent.}
\begin{proof}\pf\ Their path segment and their single edge would be
distinct simple paths, contradicting \path{SimpleGraph.IsAcyclic.path_unique}.~\qed\end{proof}
\qedstep
\begin{proof}\pf\ Combine \stepref{rlconsecutive}, \stepref{rlchord}
and the no-loop condition. ~\qed\end{proof}
\end{proof}
\step{rlreversal}{The reversal of every path labelling is a path labelling.}
\begin{proof}\pf\ The involution $i\mapsto m-1-i$ preserves
$|i-j|=1$, so composing with it preserves the defining edge equivalence.~\qed\end{proof}
\step{rlstart}{A prescribed vertex $w$ of degree at most 1 can be first.}
\begin{proof}
\step{rlone}{\case{$m=1$.}}
\begin{proof}\pf\ The bijection \stepref{rlbijection} has sole value $w$.~\qed\end{proof}
\step{rlmany}{\case{$m\geq2$.}}
\begin{proof}
\step{rlendpoint}{$w=v_0$ or $w=v_k$.}
\begin{proof}\pf\ All other positions have two neighbours by
\stepref{rledges}.~\qed\end{proof}
\step{rlforward}{\assume{$w=v_0$.}
\prove{$e$ is a path labelling with $e(0)=w$.}}
\begin{proof}\pf\ Its first value is $w$ by substitution. ~\qed\end{proof}
\step{rlreverse}{\assume{$w=v_k$.}
\prove{$i\mapsto e(m-1-i)$ is a path labelling whose first value is $w$.}}
\begin{proof}\pf\ Use \stepref{rlreversal}; the reversed first
value is $e(m-1)=v_k=w$.~\qed\end{proof}
\qedstep
\begin{proof}\pf\ Use \stepref{rlendpoint} to select
\stepref{rlforward} or \stepref{rlreverse}.~\qed\end{proof}
\end{proof}
\qedstep
\begin{proof}\pf\ Nonemptiness gives $m=1$ or $m\geq2$, the cases
\stepref{rlone}, \stepref{rlmany}.~\qed\end{proof}
\end{proof}
\qedstep
\begin{proof}\pf\ Combine \stepref{rlbijection}, \stepref{rledges},
\stepref{rlstart}, \stepref{rlreversal}.~\qed\end{proof}
\end{proof}

\begin{namedresult}{Lemma}{(three-arm labelling)}\label{lem:arm-labelling}
\lean{SimpleGraph.exists_three_arm_labelling}
\uses{def:affine,lem:path-labelling,facts:rem-library}
\assume{$X$ is a finite connected acyclic simple graph, $v$ has degree 3,
and all its other vertices have degree at most 2.}
\prove{There are positive integers $a\leq b\leq c$ and a bijection
$e:V_{a,b,c}\to V(X)$ with $e(*)=v$ whose edges are exactly the centre
edges and the consecutive arm edges of Definition~\ref{def:affine}.}
\end{namedresult}
\begin{proof}
\pfsketch\ Partition the other vertices by the first edge of the unique
path from $v$, then apply the endpoint version of path labelling to each part.
\step{raunique}{There is a unique simple path from $v$ to each vertex.}
\begin{proof}\pf\ Combine connectedness with
\path{SimpleGraph.IsAcyclic.path_unique}.~\qed\end{proof}
\pflet{Enumerate the neighbours of $v$ as $w_1,w_2,w_3$.
Let $P_r$ consist of vertices whose path from $v$ starts with $vw_r$.}
\step{rapartition}{$P_1,P_2,P_3$ are disjoint nonempty sets partitioning
$V(X)\setminus\{v\}$.}
\begin{proof}\pf\ Each nonconstant path in \stepref{raunique} has one
first edge, and $w_r\in P_r$.~\qed\end{proof}
\step{raprefix}{For every $r$ and $z\in P_r$, the segment from $w_r$ to
$z$ stays inside $P_r$.}
\begin{proof}\pf\ Each prefix is the unique path to its endpoint by
\stepref{raunique}.~\qed\end{proof}
\step{ranocross}{There is no edge between distinct $P_r,P_s$.}
\begin{proof}\pf\ Such an edge would give two paths from $v$ to an
endpoint, contradicting \stepref{raunique}; their prefixes cannot meet
after $v$ by that same uniqueness. ~\qed\end{proof}
\step{raattachment}{The only neighbour of $v$ in $P_r$ is $w_r$.}
\begin{proof}\pf\ The unique path to any neighbour of $v$ is its
single edge, by \stepref{raunique}.~\qed\end{proof}
\step{raarm}{For each $r\in\{1,2,3\}$ the induced graph on $P_r$ admits
a path labelling starting at $w_r$.}
\begin{proof}
\step{raarmconn}{The induced graph on $P_r$ is connected and acyclic.}
\begin{proof}\pf\ Use \stepref{raprefix} for reachability; an induced
cycle would be a cycle in $X$.~\qed\end{proof}
\step{raarmdegree}{All its degrees are at most 2, and $w_r$ has degree
at most 1 within it.}
\begin{proof}\pf\ Deleting $v$ removes the edge $vw_r$ from a vertex
with degree at most 2; every remaining vertex also differs from $v$.~\qed\end{proof}
\qedstep
\begin{proof}\pf\ Apply Lemma~\ref{lem:path-labelling} to
\stepref{raarmconn}, \stepref{raarmdegree}, and $w_r\in P_r$.~\qed\end{proof}
\end{proof}
\step{raassemble}{These three bijections and $*\mapsto v$ give an arm
bijection with all the asserted edges.}
\begin{proof}\pf\ Take their disjoint union using \stepref{rapartition},
\stepref{ranocross}, \stepref{raattachment}, \stepref{raarm}.~\qed\end{proof}
\step{raorder}{The three positive lengths may be ordered increasingly.}
\begin{proof}\pf\ Permute the three labelled parts of
\stepref{raassemble} by a sorting permutation. ~\qed\end{proof}
\qedstep
\begin{proof}\pf\ Use \stepref{raassemble}, \stepref{raorder}.~\qed\end{proof}
\end{proof}

\begin{namedresult}{Lemma}{(unweighted trees without obstructions)}\label{lem:unweighted-trees}
\lean{SimpleGraph.dynkin_labelling_of_no_affine}
\uses{def:affine,lem:path-labelling,lem:arm-labelling,lem:rem-cycle-embedding,facts:rem-library}
\assume{$X$ is a finite connected simple graph containing no cycle,
no $\widetilde{\typ D}_\ell$ ($\ell\geq4$), and no
$\widetilde{\typ E}_6,\widetilde{\typ E}_7,\widetilde{\typ E}_8$ as a
subgraph (the subgraph need not be induced).}
\prove{$X$ has a path labelling, or an arm labelling of $T(1,1,c)$,
$c\geq1$, or of $T(1,2,2)$, $T(1,2,3)$, $T(1,2,4)$.}
\end{namedresult}
\begin{proof}
\pfsketch\ Exclude high degrees and two branch vertices, then truncate
the ordered arms.
\step{rutree}{$X$ is acyclic.}
\begin{proof}\pf\ A closed walk satisfying the cycle predicate would
embed $\Cyc_m$ for $m\geq3$ by Lemma~\ref{lem:rem-cycle-embedding},
contrary to \asm{0}; this is the definition of acyclicity.~\qed\end{proof}
\step{rudegree}{Every degree is at most 3.}
\begin{proof}\pf\ Four distinct neighbours of one vertex would embed
$\widetilde{\typ D}_4$ by Definition~\ref{def:affine}.~\qed\end{proof}
\step{ruone}{\assume{There are distinct degree-3 vertices $u,v$.}
\prove{\textsc{False}.}}
\begin{proof}
\pfsketch\ \new\ We may use any two branch vertices: additional edges
at internal path vertices are omitted from the obstruction subgraph.
\pick\ a simple path between $u,v$, by connectedness in \asm{0}.
\step{ruextra}{Each endpoint has two distinct neighbours off the path.}
\begin{proof}\pf\ An additional neighbour on the path would contradict
path uniqueness from \stepref{rutree}; the degree is 3. ~\qed\end{proof}
\step{rudistinct}{The four selected neighbours are mutually distinct.}
\begin{proof}\pf\ A common neighbour of the endpoints would give a
second simple path between them, contradicting \stepref{rutree}.~\qed\end{proof}
\step{ruembed}{The selected path and four incident edges embed
$\widetilde{\typ D}_{k+4}$, where $k\geq1$ is the path length.}
\begin{proof}\pf\ Use the injective vertex list from
\stepref{ruextra}, \stepref{rudistinct} in Definition~\ref{def:affine}.~\qed\end{proof}
\qedstep
\begin{proof}\pf\ Step \stepref{ruembed} contradicts \asm{0}.~\qed\end{proof}
\end{proof}
\step{rupath}{\case{There is no degree-3 vertex.}}
\begin{proof}\pf\ Apply Lemma~\ref{lem:path-labelling} using
\stepref{rudegree}, \stepref{rutree}.~\qed\end{proof}
\step{rubranch}{\case{There is a degree-3 vertex.}}
\begin{proof}
\step{ruarms}{\pick\ an arm bijection with $1\leq a\leq b\leq c$.}
\begin{proof}\pf\ Apply Lemma~\ref{lem:arm-labelling} to
\stepref{rutree}, \stepref{rudegree}, \stepref{ruone}.~\qed\end{proof}
\step{rua}{$a=1$.}
\begin{proof}\pf\ Otherwise the restriction to the first two positions
of each arm in \stepref{ruarms} embeds $T(2,2,2)$, forbidden by \asm{0}.~\qed\end{proof}
\step{rub}{$b\leq2$.}
\begin{proof}\pf\ If $b\geq3$, restricting the arms to lengths $1,3,3$
in \stepref{ruarms} embeds a forbidden $\widetilde{\typ E}_7$.~\qed\end{proof}
\step{rubone}{\assume{$b=1$.}
\prove{The labelling is $T(1,1,c)$.}}
\begin{proof}\pf\ Substitute in \stepref{ruarms}, using \stepref{rua}.~\qed\end{proof}
\step{rubtwo}{\assume{$b=2$.}
\prove{$c\in\{2,3,4\}$.}}
\begin{proof}\pf\ A value $c\geq5$ embeds $T(1,2,5)$ by restriction;
$2\leq c\leq4$ gives the asserted integer alternatives. ~\qed\end{proof}
\qedstep
\begin{proof}\pf\ The positive integer $b\leq2$ in \stepref{rub} is
1 or 2, so \stepref{rubone}, \stepref{rubtwo} exhaust this case. ~\qed\end{proof}
\end{proof}
\qedstep
\begin{proof}\pf\ The cases \stepref{rupath}, \stepref{rubranch} are
exhaustive. ~\qed\end{proof}
\end{proof}

\begin{namedresult}{Theorem}{(weighted dichotomy)}\label{thm:dichotomy}
\lean{WeightedGraph.affine_or_dynkin}
\uses{def:affine,def:rem-weighted,lem:unweighted-trees,lem:path-labelling,lem:affine,lem:monotonicity,lem:rem-unit-containment,lem:rem-rayleigh-reindex,lem:rem-labelled-connected}
\assume{$X$ is a finite connected weighted graph whose edge weights belong
to $\{1,\sqrt2,\sqrt3\}$.}
\prove{Exactly one alternative holds: $X$ contains an affine obstruction;
or it admits a weight-preserving path or arm bijection to one of
\[
 \begin{gathered}
 \typ A_\ell\ (\ell\geq1),\quad
 \typ{BC}_\ell\ (\ell\geq2),\quad
 T(1,1,\ell-3)=\typ D_\ell\ (\ell\geq4),\\
 T(1,2,\ell-4)=\typ E_\ell\ (\ell=6,7,8),\quad
 \typ F_4,\quad\typ G_2.
 \end{gathered}
\]
Here $\typ A$ is a unit path, $\typ{BC}$ has one terminal $\sqrt2$ edge,
$\typ D,\typ E$ have unit edges, $\typ F_4$ has path weights
$1,\sqrt2,1$, and $\typ G_2$ is one $\sqrt3$ edge.
Every graph in the second alternative has largest eigenvalue less than 2.}
\end{namedresult}
\begin{proof}
\pfsketch\ Retain the labelling of the unit support and specify the
injection for every excluded weighted pattern.
\step{rdrecognize}{\assume{$X$ contains no affine obstruction.}
\prove{$X$ admits one of the listed weight-preserving Dynkin labellings.}}
\begin{proof}
\step{rdsupport}{The support of $X$ has one of the path or arm labellings in
Lemma~\ref{lem:unweighted-trees}.}
\begin{proof}\pf\ By Lemma~\ref{lem:rem-unit-containment}, an unweighted
obstruction embedded in the support would be contained in $X$, since
all its edge weights are at least 1. Thus \asm{1} supplies all hypotheses
of Lemma~\ref{lem:unweighted-trees}.~\qed\end{proof}
\step{rdtriple}{\assume{There is a $\sqrt3$ edge with endpoints $u,v$.}
\prove{Those endpoints give a weight-preserving $\typ G_2$ labelling of $X$.}}
\begin{proof}
\step{rdthird}{\assume{There is a vertex outside $\{u,v\}$.}
\prove{\textsc{False}.}}
\begin{proof}
\step{rdexit}{There is a support edge from $\{u,v\}$ to its complement.}
\begin{proof}\pf\ Take the first exit from $\{u,v\}$ on a connecting
walk to the vertex in \asm{3}, using connectedness in \asm{0}.~\qed\end{proof}
\step{rdtripleembed}{$X^{\rm triple}\preccurlyeq X$.}
\begin{proof}\pf\ Keep the edge in \asm{2} and lower the edge in
\stepref{rdexit} to 1; their three distinct vertices give the injection
in Definition~\ref{def:affine}.~\qed\end{proof}
\qedstep
\begin{proof}\pf\ Step \stepref{rdtripleembed} contradicts \asm{1}.~\qed\end{proof}
\end{proof}
\qedstep
\begin{proof}\pf\ Step \stepref{rdthird} leaves exactly the two distinct
endpoints; their edge has weight $\sqrt3$ by \asm{2}.~\qed\end{proof}
\end{proof}
\step{rdbranch}{\assume{There is no $\sqrt3$ edge, and the support has
an arm labelling from \stepref{rdsupport}.}
\prove{This labelling identifies $X$ with a listed $\typ D$ or $\typ E$ graph.}}
\begin{proof}
\step{rdfork}{\assume{There is a $\sqrt2$ edge.}
\prove{\textsc{False}.}}
\begin{proof}
\pick\ the arm containing that edge; permute the arm labels to call it
arm 1, and let $k\geq1$ be the outer endpoint position of that edge.
\step{rdforkmap}{The map from $v_0,\ldots,v_k$ to
$e(*),e(1,1),\ldots,e(1,k)$, and from the two extra leaves to
$e(2,1),e(3,1)$, is injective.}
\begin{proof}\pf\ These positions are distinct in the arm bijection
of \asm{2}, including when $k=1$.~\qed\end{proof}
\step{rdforkweights}{This map witnesses $X^{\rm fork}_k\preccurlyeq X$.}
\begin{proof}\pf\ Retain its last $\sqrt2$ edge and lower every other
selected edge to 1 in \stepref{rdforkmap}; use the weight bounds in \asm{0}.~\qed\end{proof}
\qedstep
\begin{proof}\pf\ Step \stepref{rdforkweights} contradicts \asm{1}.~\qed\end{proof}
\end{proof}
\step{rdbranchunit}{All edges have weight 1.}
\begin{proof}\pf\ The weight alternatives in \asm{0}, absence of
$\sqrt3$ in \asm{2}, and \stepref{rdfork} exclude both other weights.~\qed\end{proof}
\step{rdbranchD}{\case{The labelling is $T(1,1,c)$, $c\geq1$.}
\prove{It is a weight-preserving $\typ D_{c+3}$ labelling.}}
\begin{proof}\pf\ Set $\ell=c+3\geq4$ in the theorem's definition of
$\typ D_\ell$, using \stepref{rdbranchunit}.~\qed\end{proof}
\step{rdbranchE}{\case{The labelling is $T(1,2,c)$, $c\in\{2,3,4\}$.}
\prove{It is a weight-preserving $\typ E_{c+4}$ labelling.}}
\begin{proof}\pf\ Set $\ell=c+4\in\{6,7,8\}$ in the theorem's definition
of $\typ E_\ell$, using \stepref{rdbranchunit}.~\qed\end{proof}
\qedstep
\begin{proof}\pf\ The two arm alternatives in \stepref{rdsupport}
are exactly \stepref{rdbranchD}, \stepref{rdbranchE}.~\qed\end{proof}
\end{proof}
\step{rdpath}{\assume{There is no $\sqrt3$ edge, and the support has
a path labelling $e:[0,\ell-1]\to V(X)$.}
\prove{$X$ has a weight-preserving $\typ A$, $\typ{BC}$ or $\typ F_4$ labelling.}}
\begin{proof}
\step{rdtwo}{There is at most one $\sqrt2$ edge.}
\begin{proof}
\step{rdtwocontra}{\assume{Two such edges have positions $r<s$,
where position $i$ denotes the edge $e(i)e(i+1)$.}
\prove{\textsc{False}.}}
\begin{proof}
\step{rdendmap}{The map $i\mapsto e(r+i)$,
$0\leq i\leq s+1-r$, is injective.}
\begin{proof}\pf\ Restrict the path bijection of \asm{2} to this interval.~\qed\end{proof}
\step{rdendembed}{This map witnesses $X^{\rm end}_{s-r+2}\preccurlyeq X$.}
\begin{proof}\pf\ Retain the first and last $\sqrt2$ edges of \asm{4}
and lower intervening weights to 1, using \stepref{rdendmap};
$s-r+2\geq3$ includes consecutive multiple edges.~\qed\end{proof}
\qedstep
\begin{proof}\pf\ Step \stepref{rdendembed} contradicts \asm{1}.~\qed\end{proof}
\end{proof}
\qedstep
\begin{proof}\pf\ Two distinct multiple edges would have ordered
positions forbidden by \stepref{rdtwocontra}.~\qed\end{proof}
\end{proof}
\step{rdunit}{\case{There is no $\sqrt2$ edge.}
\prove{$e$ is a weight-preserving $\typ A_\ell$ labelling.}}
\begin{proof}\pf\ Every edge has weight 1 by \asm{0}, \asm{2},
and this case.~\qed\end{proof}
\step{rddouble}{\case{There is a $\sqrt2$ edge.}
\prove{$X$ has a weight-preserving $\typ{BC}$ or $\typ F_4$ labelling.}}
\begin{proof}
\step{rdorient}{\pick\ a path direction in which the two sides of the
multiple edge have $1\leq p\leq q$ vertices; its position is $p-1,p$.}
\begin{proof}\pf\ The reversal assertion of Lemma~\ref{lem:path-labelling}
interchanges the two sides, so one of the two directions orders their sizes.~\qed\end{proof}
\step{rdterminal}{\case{$p=1$.}
\prove{This is a $\typ{BC}_{p+q}$ labelling.}}
\begin{proof}\pf\ The unique multiple edge is terminal by
\stepref{rdorient}; all others are unit by \stepref{rdtwo}, \asm{2}.~\qed\end{proof}
\step{rdF}{\case{$p=q=2$.}
\prove{This is a $\typ F_4$ labelling.}}
\begin{proof}\pf\ Its three weights are $1,\sqrt2,1$ by
\stepref{rdtwo}, \stepref{rdorient}, \asm{2}.~\qed\end{proof}
\step{rdmiddle}{\case{$p\geq2,q\geq3$.}
\prove{\textsc{False}.}}
\begin{proof}
\step{rdmidvertices}{The positions $p+2,p+1,p,p-1,p-2$ are distinct
and lie in $[0,p+q-1]$.}
\begin{proof}\pf\ The least position satisfies $p-2\geq0$, the greatest
satisfies $p+2\leq p+q-1$, and consecutive listed positions differ by 1.~\qed\end{proof}
\step{rdmidweights}{Their path weights in this order are $1,1,\sqrt2,1$.}
\begin{proof}\pf\ Use the unique multiple edge at $p-1,p$ from
\stepref{rdtwo}, \stepref{rdorient}, with no $\sqrt3$ edge by \asm{2}.~\qed\end{proof}
\qedstep
\begin{proof}\pf\ Steps \stepref{rdmidvertices}, \stepref{rdmidweights}
embed $X^{\rm mid}$, contrary to \asm{1}.~\qed\end{proof}
\end{proof}
\qedstep
\begin{proof}\pf\ The integer alternatives \stepref{rdterminal},
\stepref{rdF}, \stepref{rdmiddle} exhaust $1\leq p\leq q$.~\qed\end{proof}
\end{proof}
\qedstep
\begin{proof}\pf\ A $\sqrt2$ edge either exists or does not;
apply \stepref{rddouble} or \stepref{rdunit}.~\qed\end{proof}
\end{proof}
\qedstep
\begin{proof}\pf\ If a $\sqrt3$ edge exists, use \stepref{rdtriple}.
Otherwise \stepref{rdsupport} supplies an arm or a path labelling,
the hypotheses of \stepref{rdbranch} or \stepref{rdpath}.~\qed\end{proof}
\end{proof}
\step{rdbelow}{Every listed Dynkin graph has largest eigenvalue less than 2.}
\begin{proof}
\step{rdAone}{$\typ A_1$ is a proper subgraph of $\Cyc_3$.}
\begin{proof}\pf\ Map its sole vertex to any cycle vertex. ~\qed\end{proof}
\step{rdAmany}{For $\ell\geq2$, $\typ A_\ell$ is a proper subgraph of
$\Cyc_{\ell+1}$.}
\begin{proof}\pf\ Delete one vertex of that cycle. ~\qed\end{proof}
\step{rdDfour}{$\typ D_4$ is a proper subgraph of $\widetilde{\typ D}_4$.}
\begin{proof}\pf\ Delete one leaf of $K_{1,4}$.~\qed\end{proof}
\step{rdDmany}{For $\ell\geq5$, $\typ D_\ell$ is a proper subgraph of
$\widetilde{\typ D}_\ell$.}
\begin{proof}\pf\ Delete one leaf at one fork; the other fork has arms
$1,1,(\ell-4)+1$ by Definition~\ref{def:affine}.~\qed\end{proof}
\step{rdE6}{$\typ E_6$ embeds properly in $T(2,2,2)$.}
\begin{proof}\pf\ Restrict one arm to length 1. ~\qed\end{proof}
\step{rdE7}{$\typ E_7$ embeds properly in $T(1,3,3)$.}
\begin{proof}\pf\ Restrict one length-3 arm to length 2. ~\qed\end{proof}
\step{rdE8}{$\typ E_8$ embeds properly in $T(1,2,5)$.}
\begin{proof}\pf\ Restrict the length-5 arm to length 4. ~\qed\end{proof}
\step{rdBC}{$\typ{BC}_\ell$ embeds properly in $X^{\rm end}_{\ell+1}$.}
\begin{proof}\pf\ Delete its first endpoint, retaining the last
$\sqrt2$ edge; this includes $\ell=2$.~\qed\end{proof}
\step{rdFsub}{$\typ F_4$ embeds properly in $X^{\rm mid}$.}
\begin{proof}\pf\ Delete the first endpoint of the path with weights
$1,1,\sqrt2,1$.~\qed\end{proof}
\step{rdGsub}{$\typ G_2$ embeds properly in $X^{\rm triple}$.}
\begin{proof}\pf\ Delete the endpoint incident only with the unit edge. ~\qed\end{proof}
\qedstep
\begin{proof}\pf\ Apply Lemmas~\ref{lem:affine}, \ref{lem:monotonicity}
to each containment \stepref{rdAone}--\stepref{rdGsub};
the larger obstruction is connected by Lemma~\ref{lem:rem-labelled-connected}.
Lemma~\ref{lem:rem-rayleigh-reindex} transfers the bound along every
weight-preserving labelling.~\qed\end{proof}
\end{proof}
\step{rdexclusive}{No graph admitting a listed Dynkin labelling contains an obstruction.}
\begin{proof}\pf\ Such a containment would give $2\leq\lambda$ by
Lemmas~\ref{lem:affine}, \ref{lem:monotonicity}, contradicting
\stepref{rdbelow} after transfer by Lemma~\ref{lem:rem-rayleigh-reindex};
its labelled support is connected by Lemma~\ref{lem:rem-labelled-connected}.~\qed\end{proof}
\qedstep
\begin{proof}\pf\ Combine \stepref{rdrecognize}, \stepref{rdbelow},
\stepref{rdexclusive} with excluded middle for obstruction containment. ~\qed\end{proof}
\end{proof}

\begin{namedresult}{Definition}{(the indexed matrices)}\label{def:rem-matrices}
\lean{CartanMatrix.modelEntries}\uses{facts:rem-library,def:rem-bases}
\assume{The rank is in the corresponding range of \eqref{eq:list};
indices are $0,\ldots,\ell-1$.}
All diagonals are 2 and unspecified entries are zero. Unit edges have
both entries $-1$. For $\typ A$ the edges are consecutive pairs.
For $\typ B$ they are consecutive pairs except that the last forward
entry is $-2$; $\typ C$ is its transpose.
For $\typ D$ the edges are $\{i,i+1\}$ for $0\leq i\leq\ell-3$,
and $\{\ell-3,\ell-1\}$. For $\typ E$ they are
$\{0,2\},\{1,3\}$ and $\{i,i+1\}$ for $2\leq i\leq\ell-2$.
The two remaining matrices are
\[
 \typ F_4=\begin{pmatrix}2&-1&0&0\\-1&2&-2&0\\0&-1&2&-1\\0&0&-1&2\end{pmatrix},
 \qquad \typ G_2=\begin{pmatrix}2&-3\\-1&2\end{pmatrix}.
\]
The following positive integer diagonal vectors will be used as symmetrizers:
\begin{equation}\label{eq:symmetrizers}
\begin{array}{c|c}
\typ A,\typ D,\typ E&(1,\ldots,1)\\
\typ B&(1,\ldots,1,2)\\
\typ C&(2,\ldots,2,1)\\
\typ F_4&(1,1,2,2)\\
\typ G_2&(1,3).
\end{array}
\end{equation}
These are the defining expressions in the Cartan file in
\ref{facts:rem-library}, not new matrix conventions.
\end{namedresult}

\begin{namedresult}{Lemma}{(supports of the named matrices)}\label{lem:rem-model-supports}
\lean{CartanMatrix.model_support_labelling}
\lean{CartanMatrix.model_support_connected}
\uses{def:rem-matrices,def:affine,def:rem-weighted,lem:rem-labelled-connected,facts:rem-library}
\assume{$N$ is $\typ A_\ell$ with $\ell\geq1$, $\typ B_\ell$ or
$\typ C_\ell$ with $\ell\geq2$, $\typ D_\ell$ with $\ell\geq4$,
$\typ E_\ell$ with $\ell\geq5$, $\typ F_4$, or $\typ G_2$, using the
actual Mathlib definitions on $\mathrm{Fin}\,\ell$.}
\pflet{$W_N(i,i)=0$ and
$W_N(i,j)=\sqrt{N(i,j)N(j,i)}$ for $i\ne j$.}
\prove{$W_N$ is a weighted graph with the same support as the off-diagonal
nonzero entries of $N$. The following maps are weight-preserving
bijections onto it, with exactly the displayed edges, and its support is connected:
\[
\begin{array}{c|c|c}
N&\text{domain and map}&\text{successive weights}\\ \hline
\typ A_\ell&[0,\ell-1],\ i\mapsto i&1,\ldots,1\\
\typ B_\ell,\typ C_\ell&[0,\ell-1],\ i\mapsto i&1,\ldots,1,\sqrt2\\
\typ F_4&[0,3],\ i\mapsto i&1,\sqrt2,1\\
\typ G_2&[0,1],\ i\mapsto i&\sqrt3\\
\typ D_\ell&T(1,1,\ell-3),\ h_D&\text{all }1\\
\typ E_\ell&T(1,2,\ell-4),\ h_E&\text{all }1
\end{array}
\]
Here the arm bijections are
\begin{equation}\label{eq:rem-support-maps}
\begin{aligned}
h_D(*)&=\ell-3,&h_D(1,1)&=\ell-2,&h_D(2,1)&=\ell-1,\\
h_D(3,t)&=\ell-3-t&&&(1\leq t\leq\ell-3),\\
h_E(*)&=3,&h_E(1,1)&=1,&h_E(2,1)&=2,\\
h_E(2,2)&=0,&h_E(3,t)&=3+t&&(1\leq t\leq\ell-4).
\end{aligned}
\end{equation}
All diagonals of $N$ are 2, all unit-edge entries are $(-1,-1)$, and
all nonedges have zero entries in both directions. In the displayed
path orders the exceptional forward/reverse pairs are $(-2,-1)$ on
the last $\typ B$ edge, $(-1,-2)$ on the last $\typ C$ edge,
$(-2,-1)$ on the middle $\typ F$ edge, and $(-3,-1)$ on the $\typ G$ edge.}
\end{namedresult}
\begin{proof}
\pfsketch\ Check the matrix definitions on their own index types.
The arm maps are the inverses of the index orders used to recognize
arbitrary matrices; this lemma establishes their properties independently.
\step{rmsA}{The identity is the asserted unit-path labelling for $\typ A_\ell$.}
\begin{proof}\pf\ The off-diagonal nonzero clause of
\path{CartanMatrix.A} is $i+1=j$ or $j+1=i$, with entry $-1$.~\qed\end{proof}
\step{rmsB}{The identity is the asserted path labelling for $\typ B_\ell$.}
\begin{proof}\pf\ In \path{CartanMatrix.B}, the forward consecutive
entry is $-2$ exactly at $j=\ell-1$; all reverse consecutive entries
are $-1$, and every other off-diagonal entry vanishes.~\qed\end{proof}
\step{rmsC}{The identity is the asserted path labelling for $\typ C_\ell$.}
\begin{proof}\pf\ In \path{CartanMatrix.C}, the reverse consecutive
entry is $-2$ exactly at $i=\ell-1$; the other consecutive entries
are $-1$, and the other off-diagonal entries vanish.~\qed\end{proof}
\step{rmsF}{The identity is the asserted path labelling for $\typ F_4$.}
\begin{proof}\pf\ Read the consecutive pairs $(-1,-1),(-2,-1),(-1,-1)$
and all zero nonedges in \texttt{CartanMatrix.F}\textsubscript{4}.~\qed\end{proof}
\step{rmsG}{The identity is the asserted path labelling for $\typ G_2$.}
\begin{proof}\pf\ Its only off-diagonal pair is $(-3,-1)$ in
\texttt{CartanMatrix.G}\textsubscript{2}.~\qed\end{proof}
\step{rmsD}{$h_D$ is the asserted unit-arm bijection.}
\begin{proof}
\step{rmsDbij}{$h_D$ is bijective.}
\begin{proof}\pf\ Its centre and short arms cover $\ell-3,\ell-2,\ell-1$;
$t\mapsto\ell-3-t$ bijects arm 3 onto $[0,\ell-4]$.
These ranges are disjoint and exhaust $\mathrm{Fin}\,\ell$.~\qed\end{proof}
\step{rmsDentries}{The nonzero off-diagonal unordered pairs of
$\typ D_\ell$ are
$\{i,i+1\}$ for $0\leq i\leq\ell-4$ and
$\{\ell-3,\ell-2\},\{\ell-3,\ell-1\}$, all unit pairs.}
\begin{proof}\pf\ For $i<j$, the clauses of \path{CartanMatrix.D}
reduce to $j=i+1$, $j+2<\ell$, or
$i=\ell-3$, $j\in\{\ell-2,\ell-1\}$.
The reversed clauses give the reversed entries.~\qed\end{proof}
\step{rmsDcentre}{The three centre edges map to
$\{\ell-3,\ell-2\},\{\ell-3,\ell-1\},\{\ell-3,\ell-4\}$.}
\begin{proof}\pf\ Substitute the centre and first arm positions
in \eqref{eq:rem-support-maps}.~\qed\end{proof}
\step{rmsDarm}{The remaining arm edges map exactly to
$\{i,i+1\}$ for $0\leq i$ and $i+5\leq\ell$.}
\begin{proof}\pf\ On arm 3, consecutive $t,t+1$ give consecutive
$\ell-3-t,\ell-4-t$; the interval is empty at $\ell=4$.
The two length-1 arms have no internal edges.~\qed\end{proof}
\qedstep
\begin{proof}\pf\ The edge images in \stepref{rmsDcentre},
\stepref{rmsDarm} exhaust \stepref{rmsDentries}; bijectivity
\stepref{rmsDbij} therefore also excludes every other edge.~\qed\end{proof}
\end{proof}
\step{rmsE}{$h_E$ is the asserted unit-arm bijection for every $\ell\geq5$.}
\begin{proof}
\step{rmsEbij}{$h_E$ is bijective.}
\begin{proof}\pf\ The centre and first two arms cover $3,1,2,0$;
$t\mapsto3+t$ bijects arm 3 onto $[4,\ell-1]$,
disjoint from those four values.~\qed\end{proof}
\step{rmsEentries}{The nonzero off-diagonal unordered pairs are
$\{0,2\},\{1,3\}$ and $\{i,i+1\}$ for $2\leq i\leq\ell-2$,
all unit pairs.}
\begin{proof}\pf\ These are exactly the nonzero clauses in
\path{CartanMatrix.E}, valid in the whole assumed range $\ell\geq5$.~\qed\end{proof}
\step{rmsEcentre}{The centre edges map to $\{3,1\},\{3,2\},\{3,4\}$.}
\begin{proof}\pf\ Substitute the first position on each arm
in \eqref{eq:rem-support-maps}.~\qed\end{proof}
\step{rmsEarm}{The remaining edges map to $\{2,0\}$ and
$\{i,i+1\}$ for $4\leq i\leq\ell-2$.}
\begin{proof}\pf\ The two positions of arm 2 map to $2,0$;
consecutive positions on arm 3 map to consecutive integers starting at 4.
Arm 1 has no internal edge.~\qed\end{proof}
\qedstep
\begin{proof}\pf\ The edge images in \stepref{rmsEcentre},
\stepref{rmsEarm} exhaust \stepref{rmsEentries}; use
\stepref{rmsEbij} for the reverse edge implication.~\qed\end{proof}
\end{proof}
\step{rmsweights}{The weight and support assertions follow from these entries.}
\begin{proof}\pf\ The products of the displayed nonzero pairs are
$1,2,3$; every other off-diagonal product is zero.
Their nonnegative square roots define $W_N$.~\qed\end{proof}
\step{rmsconnected}{Every displayed support is connected.}
\begin{proof}\pf\ Apply Lemma~\ref{lem:rem-labelled-connected}
to the path or arm bijections in \stepref{rmsA}--\stepref{rmsE}.~\qed\end{proof}
\qedstep
\begin{proof}\pf\ Combine \stepref{rmsA}--\stepref{rmsconnected};
the diagonal clause of each named definition is 2.~\qed\end{proof}
\end{proof}

\begin{namedresult}{Definition}{(integer matrix invariants)}\label{def:rem-invariants}
\lean{CartanMatrix.supportInvariants}
\uses{}
\assume{$N:I\times I\to\mathbb Z$ is a matrix on a finite set with
$N(i,j)=0\iff N(j,i)=0$.}
Let $d_N(i)=|\{j:j\ne i,\ N(i,j)\ne0\}|$ be its support degree.
Define the following nonnegative integers, with cardinalities regarded in $\mathbb Z$:
\begin{equation}\label{eq:rem-invariants}
\begin{aligned}
n(N)&=|I|,\\
t(N)&=|\{(i,j):N(i,j)N(j,i)=3\}|,\\
u(N)&=|\{(i,j):N(i,j)N(j,i)=2\}|,\\
b(N)&=|\{i:d_N(i)=3\}|,\\
a_{\rm in}(N)&=|\{(i,j):N(i,j)=-2,\ d_N(j)=1\}|,\\
a_{\rm out}(N)&=|\{(i,j):N(i,j)=-2,\ d_N(i)=1\}|,\\
s(N)&=|\{i:d_N(i)=1,\ \exists j,\ N(i,j)\ne0,\ d_N(j)=3\}|.
\end{aligned}
\end{equation}
Write $\mathcal I(N)=(n,t,u,b,a_{\rm in},a_{\rm out},s)(N)$.
Pairs in $t,u,a_{\rm in},a_{\rm out}$ are ordered; $s$ counts leaves,
not pairs.
\end{namedresult}

\begin{namedresult}{Lemma}{(invariance and values)}\label{lem:rem-invariants}
\lean{CartanMatrix.supportInvariants_reindex}
\lean{CartanMatrix.model_supportInvariants}
\uses{def:rem-invariants,lem:rem-model-supports,def:rem-bases}
\assume{$N$ is an integer matrix on a finite set $I$ with symmetric
zero relation.}
\prove{For every bijection $e:J\to I$, the matrix
$N_e(i,j)=N(e(i),e(j))$ satisfies $\mathcal I(N_e)=\mathcal I(N)$.
For every permitted named matrix, the values are:
\begin{equation}\label{eq:rem-invariant-table}
\begin{array}{c|rrrrrrr}
N&n&t&u&b&a_{\rm in}&a_{\rm out}&s\\ \hline
\typ A_\ell\ (\ell\geq1)&\ell&0&0&0&0&0&0\\
\typ B_2&2&0&2&0&1&1&0\\
\typ B_\ell\ (\ell\geq3)&\ell&0&2&0&1&0&0\\
\typ C_\ell\ (\ell\geq3)&\ell&0&2&0&0&1&0\\
\typ D_4&4&0&0&1&0&0&3\\
\typ D_\ell\ (\ell\geq5)&\ell&0&0&1&0&0&2\\
\typ E_6&6&0&0&1&0&0&1\\
\typ E_7&7&0&0&1&0&0&1\\
\typ E_8&8&0&0&1&0&0&1\\
\typ F_4&4&0&2&0&0&0&0\\
\typ G_2&2&2&0&0&0&0&0
\end{array}
\end{equation}}
\end{namedresult}
\begin{proof}
\pfsketch\ Transport finite counting sets along the bijection, then
count on the independently labelled supports of the named matrices.
\step{rivreindex}{All seven integers are preserved by reindexing.}
\begin{proof}
\step{rivdegree}{$d_{N_e}(i)=d_N(e(i))$ for every $i$.}
\begin{proof}\pf\ The map $j\mapsto e(j)$ bijects the neighbour sets
in Definition~\ref{def:rem-invariants}.~\qed\end{proof}
\step{rivvertices}{The integers $n,b,s$ are preserved.}
\begin{proof}\pf\ Transport their defining vertex sets by $e$,
using \stepref{rivdegree} and the entry equality defining $N_e$.~\qed\end{proof}
\step{rivpairs}{The integers $t,u,a_{\rm in},a_{\rm out}$ are preserved.}
\begin{proof}\pf\ The map $(i,j)\mapsto(e(i),e(j))$ bijects each
defining pair set, by \stepref{rivdegree} and the entry equality.~\qed\end{proof}
\qedstep
\begin{proof}\pf\ Combine \stepref{rivvertices}, \stepref{rivpairs}.~\qed\end{proof}
\end{proof}
\step{rivA}{The $\typ A$ row has only its size nonzero.}
\begin{proof}\pf\ Lemma~\ref{lem:rem-model-supports} gives a unit
path, with no degree-3 vertex or $-2$ entry; use
Definition~\ref{def:rem-invariants}.~\qed\end{proof}
\step{rivBtwo}{The $\typ B_2$ row has $u=2$ and
$(a_{\rm in},a_{\rm out})=(1,1)$.}
\begin{proof}\pf\ The sole double edge has two ordered products 2;
both endpoints of its sole $-2$ entry are leaves, by
Lemma~\ref{lem:rem-model-supports}.~\qed\end{proof}
\step{rivB}{For $\ell\geq3$, the $\typ B_\ell$ row has $u=2$ and
$(a_{\rm in},a_{\rm out})=(1,0)$.}
\begin{proof}\pf\ The sole $-2$ entry points from position $\ell-2$
of degree 2 to the leaf $\ell-1$ in Lemma~\ref{lem:rem-model-supports};
the path has no branch vertex.~\qed\end{proof}
\step{rivC}{For $\ell\geq3$, the $\typ C_\ell$ row has $u=2$ and
$(a_{\rm in},a_{\rm out})=(0,1)$.}
\begin{proof}\pf\ The sole $-2$ entry points from the leaf $\ell-1$
to position $\ell-2$ of degree 2 in Lemma~\ref{lem:rem-model-supports};
the path has no branch vertex.~\qed\end{proof}
\step{rivDfour}{For $\typ D_4$, $b=1$ and $s=3$.}
\begin{proof}\pf\ The unit tree $T(1,1,1)$ in
Lemma~\ref{lem:rem-model-supports} has one degree-3 centre adjacent
to all three leaves.~\qed\end{proof}
\step{rivDmany}{For $\typ D_\ell$, $\ell\geq5$, $b=1$ and $s=2$.}
\begin{proof}\pf\ In $T(1,1,\ell-3)$ the third arm has length at least 2,
so only the first two leaves meet the centre, by
Lemma~\ref{lem:rem-model-supports}.~\qed\end{proof}
\step{rivE}{For each of $\typ E_6,\typ E_7,\typ E_8$, $b=1$ and $s=1$.}
\begin{proof}\pf\ The arm lengths $1,2,\ell-4$ in
Lemma~\ref{lem:rem-model-supports} have exactly one length equal to 1;
all edges are unit.~\qed\end{proof}
\step{rivF}{For $\typ F_4$, $u=2$ and
$(a_{\rm in},a_{\rm out})=(0,0)$.}
\begin{proof}\pf\ The sole $-2$ entry joins the two degree-2 vertices
of the four-vertex path in Lemma~\ref{lem:rem-model-supports}.~\qed\end{proof}
\step{rivG}{For $\typ G_2$, $n=t=2$ and
$u=b=a_{\rm in}=a_{\rm out}=s=0$.}
\begin{proof}\pf\ Lemma~\ref{lem:rem-model-supports} gives two vertices,
one product-3 edge, no $-2$ entry and no branch vertex.~\qed\end{proof}
\qedstep
\begin{proof}\pf\ In the counts \stepref{rivA}--\stepref{rivG}, each
matrix has its stated size, and the listed supports and entries
exclude all uncounted patterns; together with \stepref{rivreindex}
these give \eqref{eq:rem-invariant-table}.~\qed\end{proof}
\end{proof}

\begin{namedresult}{Theorem}{A}\label{thm:A}
\lean{Matrix.IsFiniteCartan.classification}
\lean{CartanMatrix.IsFiniteTypeModel.type_eq_of_reindex_eq}
\lean{CartanMatrix.IsFiniteTypeModel.isFiniteCartan}
\lean{CartanMatrix.IsFiniteTypeModel.isIndecomposable}
\uses{def:rem-matrices,lem:rem-model-supports,lem:rem-invariants,def:rem-invariants,lem:symmetrization,thm:dichotomy,lem:affine,lem:monotonicity,lem:connected,lem:rayleigh,lem:rem-rayleigh-reindex,lem:path-labelling,facts:rem-library,def:rem-bases}
\assume{$I$ is a finite nonempty set and $M:I\times I\to\mathbb Z$
satisfies \path{Matrix.IsFiniteCartan} and \path{Matrix.IsIndecomposable}.}
\prove{There are $\ell>0$ and a bijection $e:[0,\ell-1]\to I$ such that
$M(e(i),e(j))$ is a matrix in \eqref{eq:list}. Its type and rank are
unique; the bijection need not be unique. Conversely every listed matrix
is finite Cartan and indecomposable.}
\end{namedresult}
\noindent\repair\ The empty matrix satisfies both library predicates.
It is excluded here, or may be added as a separate rank-zero alternative.
\begin{proof}
\pfsketch\ Classify the weighted support, orient its one possible multiple
edge, and distinguish the resulting matrices by intrinsic invariants.
\step{rAsupport}{The symmetrized support of $M$ has a Dynkin labelling.}
\begin{proof}
\step{rAbound}{The symmetrization has $\lambda<2$.}
\begin{proof}\pf\ Apply Lemma~\ref{lem:symmetrization}.~\qed\end{proof}
\step{rAnoobstruction}{It contains no affine obstruction.}
\begin{proof}\pf\ Lemmas~\ref{lem:affine}, \ref{lem:monotonicity}
would give $2\leq\lambda$, contradicting \stepref{rAbound}.~\qed\end{proof}
\qedstep
\begin{proof}\pf\ Use the no-obstruction alternative of
Theorem~\ref{thm:dichotomy} with \stepref{rAnoobstruction}.~\qed\end{proof}
\end{proof}
\step{rAentries}{The possible ordered entry pairs are
\[
\begin{array}{c|l}
\text{unit edge}&(-1,-1)\\
\text{double edge}&(-1,-2)\text{ or }(-2,-1)\\
\text{triple edge}&(-1,-3)\text{ or }(-3,-1).
\end{array}
\]}
\begin{proof}
\step{rAunit}{Product 1 forces $(-1,-1)$.}
\begin{proof}\pf\ The two negative integers are negatives of positive
divisors of 1. ~\qed\end{proof}
\step{rAdouble}{Product 2 forces $(-1,-2)$ or $(-2,-1)$.}
\begin{proof}\pf\ The positive divisors of 2 are 1 and 2. ~\qed\end{proof}
\step{rAtriple}{Product 3 forces $(-1,-3)$ or $(-3,-1)$.}
\begin{proof}\pf\ The positive divisors of 3 are 1 and 3. ~\qed\end{proof}
\qedstep
\begin{proof}\pf\ Apply these three arithmetic cases to
Lemma~\ref{lem:symmetrization}.~\qed\end{proof}
\end{proof}
\step{rAlabelling}{There are $\ell>0$, a listed type of rank $\ell$,
and a bijection $e:\mathrm{Fin}\,\ell\to I$ such that
$M(e(i),e(j))$ equals that named matrix at every $i,j$.}
\begin{proof}
\step{rApath}{For a unit path, its path bijection gives $\typ A_\ell$.}
\begin{proof}\pf\ Substitute its edge equivalence and
\stepref{rAentries} in Definition~\ref{def:rem-matrices}.~\qed\end{proof}
\step{rAD}{For $T(1,1,\ell-3)$, use its arm bijection $f$ and set
\[
 \begin{split}
 e(\ell-3)&=f(*),\quad e(\ell-2)=f(1,1),\quad e(\ell-1)=f(2,1),\\
 e(i)&=f(3,\ell-3-i)\quad(0\leq i\leq\ell-4).
 \end{split}
\]
Then $M(e(i),e(j))=\typ D_\ell(i,j)$.}
\begin{proof}
\step{rADbij}{The formula for $e$ is a bijection.}
\begin{proof}\pf\ It is $f\circ h_D^{-1}$, a composition of bijections by
Lemma~\ref{lem:rem-model-supports}. ~\qed\end{proof}
\step{rADedges}{The edges become $\{i,i+1\}$ for $0\leq i\leq\ell-3$
and $\{\ell-3,\ell-1\}$.}
\begin{proof}\pf\ Transport the edges of the arm labelling $f$ along
$h_D^{-1}$ from Lemma~\ref{lem:rem-model-supports}.~\qed\end{proof}
\qedstep
\begin{proof}\pf\ Use \stepref{rADbij}, \stepref{rADedges},
\stepref{rAentries}, and Definition~\ref{def:rem-matrices}.~\qed\end{proof}
\end{proof}
\step{rAE}{For $T(1,2,\ell-4)$, use its arm bijection $f$ and set
\[
 \begin{split}
 e(3)&=f(*),\quad e(1)=f(1,1),\quad e(2)=f(2,1),\quad e(0)=f(2,2),\\
 e(i)&=f(3,i-3)\quad(4\leq i\leq\ell-1).
 \end{split}
\]
Then $M(e(i),e(j))=\typ E_\ell(i,j)$.}
\begin{proof}
\step{rAEbij}{The formula for $e$ is a bijection.}
\begin{proof}\pf\ It is $f\circ h_E^{-1}$, a composition of bijections by
Lemma~\ref{lem:rem-model-supports}. ~\qed\end{proof}
\step{rAEedges}{Its edges become $\{0,2\},\{1,3\}$ and
$\{i,i+1\}$ for $2\leq i\leq\ell-2$.}
\begin{proof}\pf\ Transport the edges of $f$ along $h_E^{-1}$
from Lemma~\ref{lem:rem-model-supports}. ~\qed\end{proof}
\qedstep
\begin{proof}\pf\ Use \stepref{rAEbij}, \stepref{rAEedges},
\stepref{rAentries}, and Definition~\ref{def:rem-matrices}.~\qed\end{proof}
\end{proof}
\step{rABC}{For a $\typ{BC}_\ell$ support, there is a path bijection $e$
for which $M(e(i),e(j))$ is $\typ B_\ell(i,j)$ or, when $\ell\geq3$,
$\typ C_\ell(i,j)$.}
\begin{proof}
\step{rABCorient}{\pick\ a path labelling putting the terminal multiple
edge between positions $\ell-2,\ell-1$.}
\begin{proof}\pf\ Reverse the given path when necessary, by the
reversal assertion of Lemma~\ref{lem:path-labelling}.~\qed\end{proof}
\step{rAB}{\assume{its last forward entry is $-2$.}
\prove{The matrix is $\typ B_\ell$.}}
\begin{proof}\pf\ Use \stepref{rAentries} and Definition~\ref{def:rem-matrices}.~\qed\end{proof}
\step{rAC}{\assume{that entry is $-1$ and $\ell\geq3$.}
\prove{It is $\typ C_\ell$.}}
\begin{proof}\pf\ Use \stepref{rAentries} and Definition~\ref{def:rem-matrices}.~\qed\end{proof}
\step{rAtwo}{\assume{that entry is $-1$ and $\ell=2$.}
\prove{Swapping the indices
gives $\typ B_2$.}}
\begin{proof}\pf\ Transpose the two-entry pair in \stepref{rAentries}.~\qed\end{proof}
\qedstep
\begin{proof}\pf\ The entry orders and rank alternatives in
\stepref{rAB}, \stepref{rAC}, \stepref{rAtwo} are exhaustive. ~\qed\end{proof}
\end{proof}
\step{rAF}{For a $\typ F_4$ support, there is a path bijection $e$
with $M(e(i),e(j))=\typ F_4(i,j)$ for every $i,j$.}
\begin{proof}
\step{rAForient}{\pick\ the direction in which the central forward entry is $-2$.}
\begin{proof}\pf\ Reversal in Lemma~\ref{lem:path-labelling}
interchanges the two entries of the central edge in \stepref{rAentries}.~\qed\end{proof}
\qedstep
\begin{proof}\pf\ The two unit edges and the central pair in
\stepref{rAForient} are exactly the entries in
Lemma~\ref{lem:rem-model-supports}; diagonals and nonedges agree.~\qed\end{proof}
\end{proof}
\step{rAG}{For a $\typ G_2$ support, there is a bijection $e$ with
$M(e(i),e(j))=\typ G_2(i,j)$ for every $i,j$.}
\begin{proof}
\step{rAGorient}{\pick\ the endpoint order with forward entry $-3$.}
\begin{proof}\pf\ One of the two orders has that entry by \stepref{rAentries}.~\qed\end{proof}
\qedstep
\begin{proof}\pf\ Step \stepref{rAGorient} gives the off-diagonal pair
$(-3,-1)$, with diagonal 2, exactly as in
Lemma~\ref{lem:rem-model-supports}.~\qed\end{proof}
\end{proof}
\qedstep
\begin{proof}\pf\ The alternatives in \stepref{rAsupport} are precisely
\stepref{rApath}--\stepref{rAG}.~\qed\end{proof}
\end{proof}
\step{rAconverse}{Every listed matrix is finite Cartan and indecomposable.}
\begin{proof}
\pflet{For a listed matrix $M$ take its positive integer symmetrizer from \eqref{eq:symmetrizers}.}

\step{rAunitbalance}{On each unit edge $d(i)=d(j)$ and both entries are $-1$.}
\begin{proof}\pf\ Read Definition~\ref{def:rem-matrices} and
\eqref{eq:symmetrizers} on the unit edges. ~\qed\end{proof}
\step{rABbalance}{On the $\typ B$ multiple edge the two products are
$1(-2)=2(-1)$.}
\begin{proof}\pf\ Substitute its two entries and symmetrizer values. ~\qed\end{proof}
\step{rACbalance}{On the $\typ C$ multiple edge they are $2(-1)=1(-2)$.}
\begin{proof}\pf\ Substitute its two entries and symmetrizer values. ~\qed\end{proof}
\step{rAFbalance}{On the $\typ F$ multiple edge they are $1(-2)=2(-1)$.}
\begin{proof}\pf\ Substitute its two entries and symmetrizer values. ~\qed\end{proof}
\step{rAGbalance}{On the $\typ G$ edge they are $1(-3)=3(-1)$.}
\begin{proof}\pf\ Substitute its two entries and symmetrizer values. ~\qed\end{proof}
\step{rAnonedge}{On every nonedge both products vanish.}
\begin{proof}\pf\ Use the unspecified-entry rule in
Definition~\ref{def:rem-matrices}.~\qed\end{proof}
\step{rAsymmetric}{$DM$ is symmetric and $D^{1/2}MD^{-1/2}=2I-B$
for the corresponding Dynkin graph.}
\begin{proof}\pf\ Combine the entry equalities
\stepref{rAunitbalance}--\stepref{rAnonedge}, including the diagonal 2,
with the labelled weighted supports of Lemma~\ref{lem:rem-model-supports}. ~\qed\end{proof}
\step{rApositive}{$2I-B$ is positive definite over $\mathbb R$.}
\begin{proof}
\step{rAlambda}{$\lambda(B)<2$.}
\begin{proof}\pf\ Lemma~\ref{lem:rem-model-supports} gives a
weight-preserving Dynkin labelling; transfer the bound in
Theorem~\ref{thm:dichotomy} by Lemma~\ref{lem:rem-rayleigh-reindex}.~\qed\end{proof}
\step{rAquadratic}{For $\Zx\ne0$,
$\Zx^T(2I-B)\Zx\geq(2-\lambda(B))\rsq\Zx>0$.}
\begin{proof}\pf\ Use the quadratic-form bound in Lemma~\ref{lem:rayleigh}
and \stepref{rAlambda}.~\qed\end{proof}
\qedstep
\begin{proof}\pf\ Symmetry from \stepref{rAsymmetric} and
\stepref{rAquadratic} are positive definiteness.~\qed\end{proof}
\end{proof}
\step{rAinteger}{$DM$ is positive on nonzero integer vectors.}
\begin{proof}\pf\ Substitute $D^{1/2}\Zx\ne0$ in
\stepref{rApositive}, using \stepref{rAsymmetric}.~\qed\end{proof}
\step{rAfinite}{The matrix is finite Cartan.}
\begin{proof}\pf\ Assemble its four defining fields from
Definition~\ref{def:rem-matrices}, \stepref{rAsymmetric}, \stepref{rAinteger}
and \path{Matrix.posDef_iff_dotProduct_mulVec}.~\qed\end{proof}
\step{rAindecomp}{The matrix is indecomposable.}
\begin{proof}\pf\ Apply Lemma~\ref{lem:connected} to the symmetric zero relation
and connected support in Lemma~\ref{lem:rem-model-supports}.~\qed\end{proof}
\qedstep
\begin{proof}\pf\ Combine \stepref{rAfinite}, \stepref{rAindecomp}.~\qed\end{proof}
\end{proof}
\step{rAunique}{\assume{$P,Q$ are two permitted named matrices equal
after reindexing by a bijection.}
\prove{Their types and ranks agree.}}
\begin{proof}
\step{rAcounts}{$\mathcal I(P)=\mathcal I(Q)$; in particular their ranks agree.}
\begin{proof}\pf\ Apply Lemma~\ref{lem:rem-invariants} to the
reindexing equality in \asm{1}.~\qed\end{proof}
\step{rAuniqueG}{\case{$t(P)=2$.}
\prove{$P,Q$ both have type $\typ G_2$.}}
\begin{proof}\pf\ Only the $\typ G_2$ row in
\eqref{eq:rem-invariant-table} has $t=2$; use \stepref{rAcounts}.~\qed\end{proof}
\step{rAuniqueF}{\case{$u(P)=2$ and
$(a_{\rm in}(P),a_{\rm out}(P))=(0,0)$.}
\prove{$P,Q$ both have type $\typ F_4$.}}
\begin{proof}\pf\ Among the four double-edge rows of
\eqref{eq:rem-invariant-table}, only $\typ F_4$ has this pair;
use \stepref{rAcounts}.~\qed\end{proof}
\step{rAuniqueBC}{\case{$u(P)=2$ and
$(a_{\rm in}(P),a_{\rm out}(P))\ne(0,0)$.}
\prove{$P,Q$ have the same $\typ B$ or $\typ C$ type and rank.}}
\begin{proof}
\step{rAuniqueBtwo}{\case{$n(P)=2$.}}
\begin{proof}\pf\ Only $\typ B_2$ among the permitted double-edge
rows has this size in \eqref{eq:rem-invariant-table};
apply \stepref{rAcounts}.~\qed\end{proof}
\step{rAuniqueBmany}{\case{$n(P)\geq3$ and $a_{\rm in}(P)=1$.}}
\begin{proof}\pf\ The pair is $(1,0)$ only for the $\typ B_\ell$
row at these ranks in \eqref{eq:rem-invariant-table};
apply \stepref{rAcounts}.~\qed\end{proof}
\step{rAuniqueCmany}{\case{$n(P)\geq3$ and $a_{\rm in}(P)=0$.}}
\begin{proof}\pf\ Under \asm{2}, the pair is $(0,1)$,
which singles out $\typ C_\ell$ in \eqref{eq:rem-invariant-table};
apply \stepref{rAcounts}.~\qed\end{proof}
\qedstep
\begin{proof}\pf\ The three eligible rows give exactly
\stepref{rAuniqueBtwo}--\stepref{rAuniqueCmany}.~\qed\end{proof}
\end{proof}
\step{rAuniqueA}{\case{$t(P)=u(P)=b(P)=0$.}
\prove{$P,Q$ both have type $\typ A$ at the same rank.}}
\begin{proof}\pf\ Only the $\typ A_\ell$ row has these three zero
counts in \eqref{eq:rem-invariant-table}; use \stepref{rAcounts}.~\qed\end{proof}
\step{rAuniqueDE}{\case{$t(P)=u(P)=0$ and $b(P)=1$.}
\prove{$P,Q$ have the same $\typ D$ or $\typ E$ type and rank.}}
\begin{proof}
\step{rAuniqueDfour}{\case{$s(P)=3$.}}
\begin{proof}\pf\ Only the $\typ D_4$ row has $s=3$ in
\eqref{eq:rem-invariant-table}; use \stepref{rAcounts}.~\qed\end{proof}
\step{rAuniqueDmany}{\case{$s(P)=2$.}}
\begin{proof}\pf\ Only the $\typ D_\ell$, $\ell\geq5$, row has $s=2$
in \eqref{eq:rem-invariant-table}; size in \stepref{rAcounts} fixes $\ell$.~\qed\end{proof}
\step{rAuniqueE}{\case{$s(P)=1$.}}
\begin{proof}\pf\ The remaining rows in \eqref{eq:rem-invariant-table}
are $\typ E_6,\typ E_7,\typ E_8$; their sizes $6,7,8$ distinguish
them by \stepref{rAcounts}.~\qed\end{proof}
\qedstep
\begin{proof}\pf\ The branched rows give precisely the three values
in \stepref{rAuniqueDfour}--\stepref{rAuniqueE}.~\qed\end{proof}
\end{proof}
\qedstep
\begin{proof}\pf\ The table exhausts the cases
\stepref{rAuniqueG}--\stepref{rAuniqueDE}; all parameter equalities
follow from the size equality in \stepref{rAcounts}.~\qed\end{proof}
\end{proof}
\qedstep
\begin{proof}\pf\ Combine \stepref{rAlabelling}, \stepref{rAconverse},
\stepref{rAunique}, composing the two labellings when comparing
two models of the original matrix.~\qed\end{proof}
\end{proof}

\begin{namedresult}{Remark}{(small ranks)}\label{rem:small}
\lean{CartanMatrix.small_rank_identifications}
\uses{def:rem-matrices,lem:rem-model-supports,facts:rem-library,lem:affine,lem:symmetrization,def:rem-bases,lem:connected,lem:monotonicity}
\assume{$\ell$ is a natural number.}
Here the symbols $\typ A_\ell,\typ B_\ell,\typ C_\ell,\typ D_\ell,\typ E_\ell$
denote Mathlib's generalized definitions at all natural ranks. An old-index
order $(i_0,\dots,i_{\ell-1})$ lists, for each new index $k$, its old index
$i_k$: the identification with $N$ means $N(k,l)=M(i_k,i_l)$.
\prove{$\typ B_1=\typ C_1=\typ D_1=\typ A_1=[2]$,
$\typ D_2=2I_2$, and the following old-index orders give the identifications:
$\typ C_2\cong\typ B_2$ by $(1,0)$;
$\typ D_3\cong\typ A_3$ by $(1,0,2)$;
$\typ E_1=\typ A_1$, $\typ E_2=2I_2$;
$\typ E_3\cong\typ A_2\oplus\typ A_1$ by $(0,2,1)$;
$\typ E_4\cong\typ A_4$ by $(0,2,3,1)$;
$\typ E_5\cong\typ D_5$ by $(0,2,3,1,4)$.
The matrices $\typ E_\ell$, $\ell\geq9$, are not finite Cartan.
All five parameterized definitions at rank 0 give the empty matrix.}
\end{namedresult}
\begin{proof}
\pfsketch\ The small identities are entrywise substitutions; the long
$\typ E$ arm contains an affine obstruction.
\step{rsmallone}{All stated rank-1 matrices equal $[2]$.}
\begin{proof}\pf\ The only index pair is diagonal in the definitions in
\ref{facts:rem-library}.~\qed\end{proof}
\step{rsmallDtwo}{$\typ D_2=2I_2$.}
\begin{proof}\pf\ Apply \path{CartanMatrix.D_two}.~\qed\end{proof}
\step{rsmallBC}{The order $(1,0)$ sends $\typ C_2$ to $\typ B_2$.}
\begin{proof}\pf\ Substitute in \path{CartanMatrix.B_two},
\path{CartanMatrix.C_two}.~\qed\end{proof}
\step{rsmallDthree}{The order $(1,0,2)$ sends $\typ D_3$ to $\typ A_3$.}
\begin{proof}\pf\ Apply \path{CartanMatrix.D_three'} in
\path{Mathlib/LinearAlgebra/Matrix/Cartan.lean}; its transposition
of indices 0 and 1 is the order $(1,0,2)$.~\qed\end{proof}
\step{rsmallEtwo}{$\typ E_2=2I_2$.}
\begin{proof}\pf\ None of the off-diagonal edge clauses of
\path{CartanMatrix.E} has indices less than 2. ~\qed\end{proof}
\step{rsmallEthree}{At rank 3 its sole edge is $\{0,2\}$.}
\begin{proof}\pf\ Evaluate the edge clauses of \path{CartanMatrix.E}.~\qed\end{proof}
\step{rsmallEfour}{At rank 4 its path order is $(0,2,3,1)$.}
\begin{proof}\pf\ Evaluate those edge clauses at indices less than 4. ~\qed\end{proof}
\step{rsmallEfive}{At rank 5 its edges in order $(0,2,3,1,4)$ are
exactly the $\typ D_5$ edges.}
\begin{proof}\pf\ Substitute that order into Definition~\ref{def:rem-matrices}.~\qed\end{proof}
\step{rsmallElarge}{For $\ell\geq9$ the leading nine vertices contain
$T(1,2,5)$.}
\begin{proof}\pf\ Restrict $h_E$ in Lemma~\ref{lem:rem-model-supports} to
the first five positions of its long arm and retain both short arms.~\qed\end{proof}
\step{rsmallnotfinite}{These $\typ E_\ell$ are not finite Cartan.}
\begin{proof}
\step{rsmallconnected}{Their support is connected.}
\begin{proof}\pf\ Apply the connectedness assertion of
Lemma~\ref{lem:rem-model-supports}, which includes every $\ell\geq9$.~\qed\end{proof}
\step{rsmallindecomposable}{The matrices are indecomposable.}
\begin{proof}\pf\ Apply Lemma~\ref{lem:connected} to
\stepref{rsmallconnected} and their symmetric entries. ~\qed\end{proof}
\step{rsmalllower}{$\lambda\geq2$ for their unit support.}
\begin{proof}\pf\ Apply Lemmas~\ref{lem:affine}, \ref{lem:monotonicity}
to \stepref{rsmallElarge}.~\qed\end{proof}
\qedstep
\begin{proof}\pf\ A finite Cartan hypothesis with
\stepref{rsmallindecomposable} would give $\lambda<2$ by
Lemma~\ref{lem:symmetrization}, contradicting \stepref{rsmalllower}.~\qed\end{proof}
\end{proof}
\step{rsmallempty}{At rank 0 there is only one matrix.}
\begin{proof}\pf\ Its index type has no elements. ~\qed\end{proof}
\qedstep
\begin{proof}\pf\ The entry checks \stepref{rsmallone}--\stepref{rsmallEfive},
\stepref{rsmallnotfinite}, \stepref{rsmallempty} give the claims. ~\qed\end{proof}
\end{proof}

\paragraph{Further linear algebra interfaces.}\label{facts:rem-dimension}
\begingroup\raggedright
\path{Matrix.IsHermitian.eigenvectorBasis} and
\path{Matrix.IsHermitian.mulVec_eigenvectorBasis}, in
\path{Mathlib/Analysis/Matrix/Spectrum.lean}, give an orthonormal
eigenbasis for a real symmetric matrix on a finite index type.
\path{LinearMap.finrank_range_add_finrank_ker}, in
\path{Mathlib/LinearAlgebra/FiniteDimensional/Lemmas.lean}, is rank-nullity;
\path{finrank_span_singleton}, in
\path{Mathlib/LinearAlgebra/FiniteDimensional/Basic.lean}, gives dimension
1 for the span of a nonzero vector.
\path{LinearIndependent.span_eq_top_of_card_eq_finrank'}, in the former
file, makes an independent family of the right cardinality spanning.
\path{linearIndependent_of_top_le_span_of_card_eq_finrank}, in
\path{Mathlib/LinearAlgebra/Dimension/OrzechProperty.lean}, proves the
converse implication for a spanning family of the right cardinality.
All applications below are over a field, with finite-dimensional spaces.
\path{Matrix.inv_def}, in
\path{Mathlib/LinearAlgebra/Matrix/NonsingularInverse.lean}, expresses
the inverse as the inverse determinant times the adjugate.
\par\endgroup

\begin{namedresult}{Lemma}{(a Gram realization)}\label{lem:rem-gram}
\lean{Matrix.exists_gram_realization}
\uses{facts:rem-dimension,facts:inner,lem:rem-null}
\assume{$Q$ is real symmetric on a finite set $I$ and
$\Zx^TQ\Zx\geq0$ for every real vector $\Zx$.}
\prove{There is a finite-dimensional Euclidean space $V$ and spanning
vectors $(\Zv_i)_{i\in I}$ with $(\Zv_i,\Zv_j)=Q(i,j)$.
For every family $(\Zw_i)_{i\in I}$ with Gram matrix $Q$ in any
real inner product space, its relations are exactly $\ker Q$ and
$\dim\operatorname{span}\{\Zw_i:i\in I\}=|I|-\dim\ker Q$.
In particular $\dim V=|I|-\dim\ker Q$ for the constructed realization.}
\end{namedresult}
\begin{proof}
\pfsketch\ \new\ Use the orthonormal eigenbasis to construct the
square root map explicitly. This replaces the ordinary proof's quotient
inner-product construction.
\step{rgdiag}{\pick\ an orthonormal eigenbasis $(\Zb_k)$ with eigenvalues $t_k$.}
\begin{proof}\pf\ Apply the eigenbasis declarations in
\ref{facts:rem-dimension}.~\qed\end{proof}
\step{rgnonneg}{$t_k\geq0$ for every $k$.}
\begin{proof}\pf\ Evaluate the nonnegative form at the unit eigenvector
$\Zb_k$ of \stepref{rgdiag}.~\qed\end{proof}
\pflet{$T\Zx=\sum_k\sqrt{t_k}(\Zb_k,\Zx)\Zb_k$,
$V=\operatorname{range}T$, and $\Zv_i=T\Zu_i$ regarded in $V$.}
\step{rgform}{$(T\Zx,T\Zy)=\Zx^TQ\Zy$.}
\begin{proof}\pf\ Expand in the orthonormal basis \stepref{rgdiag} and
use $(\sqrt{t_k})^2=t_k$ from \stepref{rgnonneg}.~\qed\end{proof}
\step{rgentries}{$(\Zv_i,\Zv_j)=Q(i,j)$.}
\begin{proof}\pf\ Specialize \stepref{rgform} to $\Zu_i,\Zu_j$.~\qed\end{proof}
\step{rgspan}{The vectors $\Zv_i$ span $V$.}
\begin{proof}\pf\ Apply the linear map $T$ to the coordinate expansion
$\Zx=\sum_i\Zx(i)\Zu_i$.~\qed\end{proof}
\step{rgrelations}{For any family $(\Zw_i)$ with Gram matrix $Q$,
$\sum_i x(i)\Zw_i=0$ if and only if $Qx=0$.}
\begin{proof}
\step{rgnorm}{$\|\sum_i x(i)\Zw_i\|^2=x^TQx$.}
\begin{proof}\pf\ Bilinearity in \ref{facts:inner} and the Gram entries. ~\qed\end{proof}
\step{rgzero}{The sum is zero if and only if $x^TQx=0$.}
\begin{proof}\pf\ Apply definiteness in \ref{facts:inner} to
\stepref{rgnorm}.~\qed\end{proof}
\qedstep
\begin{proof}\pf\ Lemma~\ref{lem:rem-null} and substitution of $Qx=0$
give the equivalence in \stepref{rgzero}.~\qed\end{proof}
\end{proof}
\step{rgdimension}{For every realizing family $(\Zw_i)$,
$\dim\operatorname{span}\{\Zw_i:i\in I\}=|I|-\dim\ker Q$.}
\begin{proof}\pf\ Apply rank-nullity from \ref{facts:rem-dimension}
to $x\mapsto\sum_i x(i)\Zw_i$: its range is the span by
\ref{facts:inner}, and its kernel is \stepref{rgrelations}.~\qed\end{proof}
\qedstep
\begin{proof}\pf\ Combine \stepref{rgentries}, \stepref{rgspan},
\stepref{rgrelations}, \stepref{rgdimension}.~\qed\end{proof}
\end{proof}

\begin{cgss}{Theorem}{3.8}\label{thm:3.8}\label{cor:D}
\lean{SimpleGraph.smith}\lean{SimpleGraph.smith_gram_delete_basis}
\lean{SimpleGraph.smith_exists_gram}
\lean{SimpleGraph.smithA_maxRayleigh}\lean{SimpleGraph.smithD_maxRayleigh}
\lean{SimpleGraph.smithE₆_maxRayleigh}\lean{SimpleGraph.smithE₇_maxRayleigh}
\lean{SimpleGraph.smithE₈_maxRayleigh}
\uses{def:affine,thm:dichotomy,lem:affine,lem:monotonicity,lem:rem-gram,lem:rem-edge-constant,lem:rem-rayleigh-reindex,lem:rayleigh,facts:rem-dimension,def:rem-weighted}
\assume{$X$ is a finite connected simple graph on $m$ vertices and $B$
is its real adjacency matrix.}
\prove{\textup{(Corollary D of the ordinary proof.)}
$\lambda(B)=2$ if and only if $X$ is isomorphic to a cycle
$\Cyc_s$ ($s\geq3$), $\widetilde{\typ D}_\ell$ ($\ell\geq4$), or one of
$\widetilde{\typ E}_6,\widetilde{\typ E}_7,\widetilde{\typ E}_8$.
For any such graph, $2I-B$ is the Gram matrix of spanning vectors of norm
$\sqrt2$ in dimension $m-1$. In every realization, deleting any one
vector leaves a basis of their span.}
\end{cgss}
\noindent\repair\ The ambient space in the deletion assertion is the span.
The assertion includes every deletion, thus includes the marked vertices
of [CGSS, Figure 1]. Its $A_n$ is our cycle $\Cyc_{n+1}$.
\begin{proof}
\pfsketch\ Strict monotonicity gives equality only at the unweighted
obstructions; a positive null vector controls every deletion.
\step{rDclassification}{$\lambda(B)=2$ exactly for graphs isomorphic
to the listed unweighted obstructions.}
\begin{proof}
\step{rDunit}{Only unweighted obstructions can be contained in $X$.}
\begin{proof}\pf\ A weight $\sqrt2$ or $\sqrt3$ cannot be at most an
adjacency entry 0 or 1, by Definition~\ref{def:rem-weighted}.~\qed\end{proof}
\step{rDforward}{\assume{$\lambda(B)=2$.}
\prove{$X$ is isomorphic to a listed unweighted obstruction.}}
\begin{proof}
\step{rDdynkin}{$X$ has no listed Dynkin labelling.}
\begin{proof}\pf\ Such a labelling would give $\lambda(B)<2$ by
Theorem~\ref{thm:dichotomy} and Lemma~\ref{lem:rem-rayleigh-reindex},
contradicting \asm{2}.~\qed\end{proof}
\step{rDobstruction}{\pick\ an obstruction $Y$ and a containment
$e:Y\preccurlyeq X$.}
\begin{proof}\pf\ Theorem~\ref{thm:dichotomy} and
\stepref{rDdynkin} give the obstruction alternative.~\qed\end{proof}
\step{rDobstructionUnit}{$Y$ is one of the listed unweighted graphs.}
\begin{proof}\pf\ Apply \stepref{rDunit} to \stepref{rDobstruction}.~\qed\end{proof}
\step{rDproper}{The containment $e$ is not proper.}
\begin{proof}\pf\ Otherwise Lemmas~\ref{lem:affine}, \ref{lem:monotonicity}
give $2=\lambda(Y)<\lambda(B)=2$, contradicting \asm{2}.~\qed\end{proof}
\step{rDequal}{$e$ is a graph isomorphism from $Y$ onto $X$.}
\begin{proof}\pf\ By Definition~\ref{def:rem-weighted},
\stepref{rDproper} says $e$ is onto and every entry agrees.
Both graphs have unit edge weights by \stepref{rDobstructionUnit}.~\qed\end{proof}
\qedstep
\begin{proof}\pf\ Use \stepref{rDobstructionUnit}, \stepref{rDequal}.~\qed\end{proof}
\end{proof}
\step{rDbackward}{\assume{$X$ is isomorphic to a listed unweighted graph $Y$.}
\prove{$\lambda(B)=2$.}}
\begin{proof}
\step{rDlisted}{$\lambda(Y)=2$.}
\begin{proof}\pf\ Apply Lemma~\ref{lem:affine}.~\qed\end{proof}
\step{rDtransport}{$\lambda(B)=\lambda(Y)$.}
\begin{proof}\raggedright\pf\ The graph isomorphism in \asm{2}
preserves unit adjacency entries; apply Lemma~\ref{lem:rem-rayleigh-reindex}.~\qed\end{proof}
\qedstep
\begin{proof}\pf\ Combine \stepref{rDlisted}, \stepref{rDtransport}.~\qed\end{proof}
\end{proof}
\qedstep
\begin{proof}\pf\ Steps \stepref{rDforward}, \stepref{rDbackward}
are the two directions of the equivalence.~\qed\end{proof}
\end{proof}
\step{rDgeometry}{\assume{$X$ is isomorphic to a listed graph.}
\prove{The stated realization and every-deletion basis assertion hold.}}
\begin{proof}
\step{rDpositivevector}{\pick\ $\Zw>0$ with $B\Zw=2\Zw$.}
\begin{proof}\pf\ Apply Lemma~\ref{lem:rayleigh} to the connected
adjacency matrix, with $\lambda(B)=2$ from \stepref{rDclassification}
and \asm{1}.~\qed\end{proof}
\pflet{$Q=2I-B$.}
\step{rDform}{For every real vector $\Zx$,
\begin{equation}\label{eq:affine-form}
 \Zx^TQ\Zx=\sum_{\{i,j\}\in E(X)}w(i)w(j)
 \left(\frac{x(i)}{w(i)}-\frac{x(j)}{w(j)}\right)^2.
\end{equation}}
\begin{proof}
\step{rDcross}{The mixed term for each edge is $-2x(i)x(j)$.}
\begin{proof}\pf\ Expand that edge's square in \eqref{eq:affine-form}.~\qed\end{proof}
\step{rDdiagonal}{The coefficient of $x(i)^2$ on the right is 2.}
\begin{proof}\pf\ It is $\sum_{j\sim i}w(j)/w(i)=2$ by the chosen
eigenvector equation. ~\qed\end{proof}
\qedstep
\begin{proof}\pf\ Compare the coefficients in \stepref{rDcross},
\stepref{rDdiagonal} with the quadratic form of $2I-B$.~\qed\end{proof}
\end{proof}
\step{rDpsd}{The form of $Q$ is nonnegative.}
\begin{proof}\pf\ Every summand in \stepref{rDform} is a nonnegative
square times a positive number. ~\qed\end{proof}
\step{rDkernel}{$\ker Q=\mathbb R\Zw$.}
\begin{proof}
\step{rDratios}{\assume{$Q\Zx=0$.}
\prove{$x(i)/w(i)=x(j)/w(j)$ on every edge.}}
\begin{proof}\pf\ Each nonnegative summand in the zero sum
\stepref{rDform} must vanish. ~\qed\end{proof}
\step{rDconstant}{\assume{$Q\Zx=0$.}
\prove{$\Zx$ is a scalar multiple of $\Zw$.}}
\begin{proof}\pf\ Apply
Lemma~\ref{lem:rem-edge-constant} to $i\mapsto x(i)/w(i)$ on the
connected adjacency graph $X$, using \stepref{rDratios}.~\qed\end{proof}
\step{rDreverse}{Every multiple of $\Zw$ is in $\ker Q$.}
\begin{proof}\pf\ Use $Q\Zw=2\Zw-B\Zw=0$.~\qed\end{proof}
\qedstep
\begin{proof}\pf\ Combine \stepref{rDconstant}, \stepref{rDreverse}.~\qed\end{proof}
\end{proof}
\step{rDrealize}{There is a spanning realization of $Q$ in dimension
$m-1$, and each realizing vector has squared norm 2.}
\begin{proof}
\step{rDgram}{A realization exists with dimension $m-\dim\ker Q$.}
\begin{proof}\pf\ Apply Lemma~\ref{lem:rem-gram} to \stepref{rDpsd}.~\qed\end{proof}
\step{rDrank}{$\dim\ker Q=1$.}
\begin{proof}\pf\ Use \stepref{rDkernel}, $\Zw\ne0$, and
\path{finrank_span_singleton}.~\qed\end{proof}
\step{rDnorm}{Every realizing vector has squared norm 2.}
\begin{proof}\pf\ Its squared norm is the diagonal entry $Q(i,i)=2$.~\qed\end{proof}
\qedstep
\begin{proof}\pf\ Combine \stepref{rDgram}, \stepref{rDrank},
\stepref{rDnorm}.~\qed\end{proof}
\end{proof}
\step{rDdelete}{For any realization and any index $i_0$, the remaining
vectors form a basis of its span.}
\begin{proof}
\step{rDrelation}{A relation omitting $i_0$ has coefficient vector $t\Zw$
and zero $i_0$-coordinate.}
\begin{proof}\pf\ Use Lemma~\ref{lem:rem-gram} and \stepref{rDkernel}.~\qed\end{proof}
\step{rDindependent}{All coefficients of this relation vanish.}
\begin{proof}\pf\ In \stepref{rDrelation}, $tw(i_0)=0$ and $w(i_0)>0$
force $t=0$.~\qed\end{proof}
\step{rDomittedspan}{The omitted vector $\Zv_{i_0}$ is in the span
of the remaining vectors.}
\begin{proof}\pf\ By Lemma~\ref{lem:rem-gram} and $Q\Zw=0$,
$\sum_i w(i)\Zv_i=0$; solving with $w(i_0)>0$ gives
$\Zv_{i_0}=-\sum_{i\ne i_0}(w(i)/w(i_0))\Zv_i$.~\qed\end{proof}
\qedstep
\begin{proof}\pf\ The independent remaining family of
\stepref{rDindependent} spans the full realizing span by
\stepref{rDomittedspan}, so it is a basis.~\qed\end{proof}
\end{proof}
\qedstep
\begin{proof}\pf\ Combine \stepref{rDrealize}, \stepref{rDdelete}.~\qed\end{proof}
\end{proof}
\qedstep
\begin{proof}\pf\ Combine \stepref{rDclassification}, \stepref{rDgeometry}.~\qed\end{proof}
\end{proof}

\subsubsection{Coordinate systems and their bases}\label{sec:rem-models}

\begin{namedresult}{Definition}{(classical coordinate models)}\label{def:rem-classical}
\lean{RootSetModel.An}\lean{RootSetModel.Bn}\lean{RootSetModel.Cn}\lean{RootSetModel.Dn}
\uses{def:roots,def:rem-bases}
\assume{$\ell\geq1$ for $\typ A$, $\ell\geq2$ for $\typ B,\typ C$,
and $\ell\geq4$ for $\typ D$.}
In standard coordinates put
\[
\begin{aligned}
\Phi(\typ A_\ell)&=\{\Zu_i-\Zu_j:1\leq i\ne j\leq\ell+1\},&
 V(\typ A_\ell)&=\{\Zx:\sum_i x(i)=0\},\\
\Phi(\typ B_\ell)&=\{\pm\Zu_i\}\cup\{\pm\Zu_i\pm\Zu_j:i<j\},&
 V(\typ B_\ell)&=\mathbb R^\ell,\\
\Phi(\typ C_\ell)&=\{\pm2\Zu_i\}\cup\{\pm\Zu_i\pm\Zu_j:i<j\},&
 V(\typ C_\ell)&=\mathbb R^\ell,\\
\Phi(\typ D_\ell)&=\{\pm\Zu_i\pm\Zu_j:i<j\},&
 V(\typ D_\ell)&=\mathbb R^\ell.
\end{aligned}
\]
The indices in the last three rows range from 1 to $\ell$.
The proposed bases are
\begin{equation}\label{eq:classical-bases}
\begin{aligned}
\typ A:\quad&\alpha_i=\Zu_{i+1}-\Zu_{i+2}&&(0\leq i<\ell);\\
\typ B,\typ C,\typ D:\quad&\alpha_i=\Zu_{i+1}-\Zu_{i+2}&&(0\leq i<\ell-1),
\end{aligned}
\end{equation}
with last root $\Zu_\ell$, $2\Zu_\ell$, or $\Zu_{\ell-1}+\Zu_\ell$,
respectively. We retain the $\typ C_2$ model for comparison with $\typ B_2$.
\end{namedresult}

\begin{namedresult}{Lemma}{(signed coordinate reflections)}\label{lem:rem-coordinate-reflections}
\lean{RootSetModel.coordinate_reflection}
\uses{lem:rem-reflection,facts:inner,def:roots}
\assume{$\Zu_1,\ldots,\Zu_m$ is the standard orthonormal basis of
$\mathbb R^m$, $i\ne j$, and $c\ne0$.}
\prove{Reflection in $\Zu_i-\Zu_j$ interchanges coordinates $i,j$;
reflection in $\Zu_i+\Zu_j$ interchanges and negates both;
reflection in $c\Zu_i$ negates coordinate $i$ and fixes the others.
Signed coordinate permutations preserve each set consisting of all axis
roots of a fixed length, all two-coordinate roots with coefficients
$\pm1$, or all half-coordinate sign vectors. Coordinate permutations
preserve all differences and the zero-sum subspace.}
\end{namedresult}
\begin{proof}
\pfsketch\ Evaluate the reflection formula in each coordinate.
\step{rcoorddiff}{The difference reflection has entries $x(j),x(i)$
at $i,j$ and $x(k)$ elsewhere.}
\begin{proof}\pf\ Substitute $(\Zx,\Zu_i-\Zu_j)=x(i)-x(j)$ and
$\|\Zu_i-\Zu_j\|^2=2$ in Definition~\ref{def:roots}.~\qed\end{proof}
\step{rcoordsum}{The sum reflection has entries $-x(j),-x(i)$
at $i,j$ and $x(k)$ elsewhere.}
\begin{proof}\pf\ Substitute $(\Zx,\Zu_i+\Zu_j)=x(i)+x(j)$ and
squared norm 2 in the reflection formula. ~\qed\end{proof}
\step{rcoordaxis}{The axis reflection has entry $-x(i)$ at $i$
and $x(k)$ elsewhere.}
\begin{proof}\pf\ Substitute $(\Zx,c\Zu_i)=cx(i)$ and squared norm
$c^2$ in the reflection formula. ~\qed\end{proof}
\step{rcoordpatterns}{Signed permutations preserve the three asserted
coordinate patterns.}
\begin{proof}\pf\ They preserve support size and the absolute value
of every nonzero coordinate. ~\qed\end{proof}
\step{rcoordpermutation}{Coordinate permutations preserve differences
and the zero-sum subspace.}
\begin{proof}\pf\ Reindex their two nonzero coordinates and the coordinate sum. ~\qed\end{proof}
\qedstep
\begin{proof}\pf\ Combine \stepref{rcoorddiff}--\stepref{rcoordpermutation}.~\qed\end{proof}
\end{proof}

\begin{namedresult}{Lemma}{(classical Root Systems)}\label{lem:classical}
\lean{RootSetModel.classical_isRootSet}\lean{RootSetModel.classical_base}
\lean{RootSetModel.classical_cartanMatrix}
\uses{def:rem-classical,def:rem-matrices,lem:rem-coordinate-reflections,lem:rem-reflection,facts:inner,facts:rem-dimension,def:rem-bases,def:roots}
\assume{The parameters, spaces, root sets and candidate bases are those
of Definition~\ref{def:rem-classical}.}
\prove{Each displayed set is a Root System in its specified space.
The displayed family is a root base and has exactly Mathlib's Cartan
matrix of that type. The dimension of each space is $\ell$.}
\end{namedresult}
\begin{proof}
\pfsketch\ Check reflection closure by coordinate operations and prove
the base property by nonnegative telescoping expansions.
\step{rclfinite}{Each root set is finite, contains no zero, and is reduced.}
\begin{proof}
\step{rclfiniteonly}{Each set is finite and excludes zero.}
\begin{proof}\pf\ Its finite index and sign choices in
Definition~\ref{def:rem-classical} have a specified nonzero coordinate. ~\qed\end{proof}
\step{rclaxisreduce}{On an axis the only listed vectors are one sign pair.}
\begin{proof}\pf\ Read the axis clause in Definition~\ref{def:rem-classical}.~\qed\end{proof}
\step{rclpairreduce}{On a two-coordinate line the absolute values of
both nonzero coordinates are 1.}
\begin{proof}\pf\ Read the two-coordinate clause in
Definition~\ref{def:rem-classical}.~\qed\end{proof}
\qedstep
\begin{proof}\pf\ Proportional vectors have identical supports, and
\stepref{rclaxisreduce}, \stepref{rclpairreduce} force the scalar to be
$1$ or $-1$; combine with \stepref{rclfiniteonly}.~\qed\end{proof}
\end{proof}
\step{rclreflection}{Each set is closed under all its root reflections.}
\begin{proof}
\step{rclrefA}{In $\typ A$ each reflection permutes coordinates.}
\begin{proof}\pf\ Apply the difference case of
Lemma~\ref{lem:rem-coordinate-reflections}.~\qed\end{proof}
\step{rclrefB}{In $\typ B$ every reflection is a signed permutation
preserving both its patterns.}
\begin{proof}\pf\ Apply the axis, sum and difference formulas of
Lemma~\ref{lem:rem-coordinate-reflections} to its listed roots. ~\qed\end{proof}
\step{rclrefC}{In $\typ C$ every reflection is a signed permutation
preserving both its patterns.}
\begin{proof}\pf\ Apply that lemma with axis scalar $c=2$, and its sum
and difference formulas, to the $\typ C$ list. ~\qed\end{proof}
\step{rclrefD}{In $\typ D$ every reflection is a signed permutation
preserving its two-coordinate pattern.}
\begin{proof}\pf\ Apply the sum and difference formulas of that lemma
to the $\typ D$ list. ~\qed\end{proof}
\qedstep
\begin{proof}\pf\ Use \stepref{rclrefA}--\stepref{rclrefD} and
$s_{-\alpha}=s_\alpha$ from Lemma~\ref{lem:rem-reflection}.~\qed\end{proof}
\end{proof}
\step{rclinteger}{Every root pairing is integral.}
\begin{proof}
\step{rclpairlengthtwo}{Pairing against a two-coordinate root is an
integer inner product.}
\begin{proof}\pf\ All coordinates are integers and its squared norm is
2 in \eqref{eq:pairing}.~\qed\end{proof}
\step{rclpairB}{Pairing against $\pm\Zu_i$ is $\pm2x(i)\in\mathbb Z$.}
\begin{proof}\pf\ Substitute squared norm 1 in \eqref{eq:pairing}.~\qed\end{proof}
\step{rclpairC}{Pairing against $\pm2\Zu_i$ is $\pm x(i)\in\mathbb Z$.}
\begin{proof}\pf\ Substitute squared norm 4 in \eqref{eq:pairing}.~\qed\end{proof}
\qedstep
\begin{proof}\pf\ The denominator-root cases
\stepref{rclpairlengthtwo}, \stepref{rclpairB}, \stepref{rclpairC}
exhaust all four root lists. ~\qed\end{proof}
\end{proof}
\step{rclexpansions}{Every listed root has coefficients of one sign in
the proposed family.}
\begin{proof}
\step{rclexpdiff}{For $i<j$ in all four types,
$\Zu_i-\Zu_j=\sum_{k=i-1}^{j-2}\alpha_k$.}
\begin{proof}\pf\ Telescope the consecutive differences in
\eqref{eq:classical-bases}.~\qed\end{proof}
\step{rclexpBaxis}{In $\typ B$, $\Zu_i=\sum_{k=i-1}^{\ell-1}\alpha_k$.}
\begin{proof}\pf\ Telescope to the endpoint $\alpha_{\ell-1}=\Zu_\ell$.~\qed\end{proof}
\step{rclexpBsum}{In $\typ B$, for $i<j$,
$\Zu_i+\Zu_j=\sum_{k=i-1}^{j-2}\alpha_k+
2\sum_{k=j-1}^{\ell-1}\alpha_k$.}
\begin{proof}\pf\ Add \stepref{rclexpdiff} to twice \stepref{rclexpBaxis}
at $j$.~\qed\end{proof}
\step{rclexpCaxis}{In $\typ C$,
$2\Zu_i=2\sum_{k=i-1}^{\ell-2}\alpha_k+\alpha_{\ell-1}$.}
\begin{proof}\pf\ Telescope to $\alpha_{\ell-1}=2\Zu_\ell$.~\qed\end{proof}
\step{rclexpCsum}{In $\typ C$, for $i<j$,
$\Zu_i+\Zu_j=\sum_{k=i-1}^{j-2}\alpha_k+
2\sum_{k=j-1}^{\ell-2}\alpha_k+\alpha_{\ell-1}$.}
\begin{proof}\pf\ Add \stepref{rclexpdiff} to \stepref{rclexpCaxis} at $j$.~\qed\end{proof}
\step{rclexpDsum}{In $\typ D$, for $i<j\leq\ell-1$,
\[
\Zu_i+\Zu_j=\sum_{k=i-1}^{j-2}\alpha_k+
2\sum_{k=j-1}^{\ell-3}\alpha_k+\alpha_{\ell-2}+\alpha_{\ell-1}.
\]}
\begin{proof}\pf\ Telescope using
$\alpha_{\ell-2}+\alpha_{\ell-1}=2\Zu_{\ell-1}$
from \eqref{eq:classical-bases}.~\qed\end{proof}
\step{rclexpDlast}{In $\typ D$, for $i<\ell$,
$\Zu_i+\Zu_\ell=\sum_{k=i-1}^{\ell-3}\alpha_k+\alpha_{\ell-1}$.}
\begin{proof}\pf\ Telescope using only
$\alpha_{\ell-1}=\Zu_{\ell-1}+\Zu_\ell$.~\qed\end{proof}
\step{rclexpnegative}{The negative of each displayed root has the
negative of its displayed coefficient vector.}
\begin{proof}\pf\ Negate each linear combination. ~\qed\end{proof}
\qedstep
\begin{proof}\pf\ The root patterns in Definition~\ref{def:rem-classical}
are exactly the vectors in \stepref{rclexpdiff}--\stepref{rclexpDlast}
and their negatives; all displayed coefficients are nonnegative integers. ~\qed\end{proof}
\end{proof}
\step{rclspan}{The roots span the specified $\ell$-dimensional spaces.}
\begin{proof}
\step{rclspanA}{The differences $\Zu_i-\Zu_{\ell+1}$,
$1\leq i\leq\ell$, form a real basis of the zero-sum space.}
\begin{proof}\pf\ The unique coefficients of a zero-sum vector are its
first $\ell$ coordinates, by coordinate substitution. ~\qed\end{proof}
\step{rclspanB}{The $\typ B$ axis roots span $\mathbb R^\ell$.}
\begin{proof}\pf\ They include its standard basis. ~\qed\end{proof}
\step{rclspanC}{The $\typ C$ axis roots span $\mathbb R^\ell$.}
\begin{proof}\pf\ Each standard vector is half a listed axis root. ~\qed\end{proof}
\step{rclspanD}{The $\typ D$ roots span $\mathbb R^\ell$.}
\begin{proof}\pf\ For each $i$, choose $j\ne i$ and use
$2\Zu_i=(\Zu_i+\Zu_j)+(\Zu_i-\Zu_j)$.~\qed\end{proof}
\qedstep
\begin{proof}\pf\ Combine \stepref{rclspanA}--\stepref{rclspanD}.~\qed\end{proof}
\end{proof}
\step{rclbase}{The proposed family is a root base.}
\begin{proof}
\step{rclsimplein}{Every proposed simple root belongs to its root list.}
\begin{proof}\pf\ Substitute the simple-root formulas in
Definition~\ref{def:rem-classical}.~\qed\end{proof}
\step{rclsimplesspan}{The $\ell$ simple roots span their ambient space.}
\begin{proof}\pf\ Combine \stepref{rclexpansions} with \stepref{rclspan}.~\qed\end{proof}
\step{rclindependent}{They are linearly independent.}
\begin{proof}\pf\ Apply the spanning-cardinality lemma of
\ref{facts:rem-dimension} to \stepref{rclsimplesspan} and dimension
$\ell$ in \stepref{rclspan}.~\qed\end{proof}
\qedstep
\begin{proof}\pf\ Use the base definition with \stepref{rclsimplein},
\stepref{rclindependent}, \stepref{rclexpansions}.~\qed\end{proof}
\end{proof}
\step{rclcartan}{The base matrices have the asserted entries.}
\begin{proof}
\step{rcldiffgram}{Consecutive differences have squared norm 2,
inner product $-1$ at consecutive indices, and 0 at nonconsecutive indices.}
\begin{proof}\pf\ Expand
$(\Zu_{i+1}-\Zu_{i+2},\Zu_{j+1}-\Zu_{j+2})$
using orthonormality. ~\qed\end{proof}
\step{rclBgram}{The $\typ B$ last root has squared norm 1, inner product
$-1$ with its predecessor, and zero with the other differences.}
\begin{proof}\pf\ Substitute $\alpha_{\ell-1}=\Zu_\ell$.~\qed\end{proof}
\step{rclCgram}{The $\typ C$ last root has squared norm 4, inner product
$-2$ with its predecessor, and zero with the other differences.}
\begin{proof}\pf\ Substitute $\alpha_{\ell-1}=2\Zu_\ell$.~\qed\end{proof}
\step{rclDgram}{The $\typ D$ last root has squared norm 2, product $-1$
with $\alpha_{\ell-3}$, and zero with all other simple roots.}
\begin{proof}\pf\ Expand the product with $\Zu_{\ell-1}+\Zu_\ell$
using orthonormality, including $1-1=0$ at $\alpha_{\ell-2}$.~\qed\end{proof}
\qedstep
\begin{proof}\pf\ Divide the columns in \stepref{rcldiffgram}--\stepref{rclDgram}
by half their squared lengths, as in \eqref{eq:pairing}, and compare
Definition~\ref{def:rem-matrices}.~\qed\end{proof}
\end{proof}
\qedstep
\begin{proof}\pf\ Assemble Definition~\ref{def:roots} from
\stepref{rclfinite}, \stepref{rclreflection}, \stepref{rclinteger},
\stepref{rclspan}, and add \stepref{rclbase}, \stepref{rclcartan}.~\qed\end{proof}
\end{proof}

\begin{namedresult}{Definition}{(exceptional coordinate data)}\label{def:rem-exceptional}
\lean{RootSetModel.E₆}\lean{RootSetModel.E₇}\lean{RootSetModel.E₈}
\lean{RootSetModel.F₄}\lean{RootSetModel.G₂}
\uses{def:roots,def:rem-bases}
\assume{Coordinates have the standard real inner product.}
In $\mathbb R^8$ define
\begin{equation}\label{eq:E-lattice}
 L=\{\Zx\in\mathbb Z^8:\sum_i x(i)\in2\mathbb Z\}
 \ \cup\
 \{\Zx\in(\mathbb Z+\tfrac12)^8:\sum_i x(i)\in2\mathbb Z\}.
\end{equation}
Put
\[
\begin{aligned}
\Phi(\typ E_8)&=\{\pm\Zu_i\pm\Zu_j:i<j\}\\
&\quad\cup\{\tfrac12\sum_{i=1}^8\epsilon(i)\Zu_i:
                \epsilon(i)\in\{1,-1\},\ \prod_i\epsilon(i)=1\},\\
V(\typ E_7)&=\{\Zx:x(7)+x(8)=0\},\\
V(\typ E_6)&=\{\Zx:x(6)=x(7)=-x(8)\},\\
\Phi(\typ E_j)&=\Phi(\typ E_8)\cap V(\typ E_j)\quad(j=6,7),
\qquad V(\typ E_8)=\mathbb R^8.
\end{aligned}
\]
The ordered $\typ E_8$ family is
\begin{equation}\label{eq:E-base}
\begin{split}
\alpha_0&=\tfrac12(\Zu_1+\Zu_8-\Zu_2-\Zu_3-\Zu_4-\Zu_5-\Zu_6-\Zu_7),\\
\alpha_1&=\Zu_1+\Zu_2,\qquad
\alpha_i=\Zu_i-\Zu_{i-1}\quad(2\leq i\leq7).
\end{split}
\end{equation}
Use its first $j$ vectors for $\typ E_j$.
In $\mathbb R^4$ put
\begin{equation}\label{eq:F-roots}
\Phi(\typ F_4)=\{\pm\Zu_i\pm\Zu_j:i<j\}\cup\{\pm\Zu_i\}
\cup\{\tfrac12\sum_{i=1}^4\epsilon(i)\Zu_i:\epsilon(i)\in\{1,-1\}\}
\end{equation}
with family
\begin{equation}\label{eq:F-base}
(\alpha_0,\alpha_1,\alpha_2,\alpha_3)=
(\Zu_2-\Zu_3,\ \Zu_3-\Zu_4,\ \Zu_4,\
\tfrac12(\Zu_1-\Zu_2-\Zu_3-\Zu_4)).
\end{equation}
In the zero-sum plane $V\subset\mathbb R^3$ put
\begin{equation}\label{eq:G-roots}
\Phi(\typ G_2)=\{\Zu_i-\Zu_j:i\ne j\}\cup
\{\pm(2\Zu_i-\Zu_j-\Zu_k):\{i,j,k\}=\{1,2,3\}\},
\end{equation}
with family
\begin{equation}\label{eq:G-base}
\alpha_0=\Zu_1-2\Zu_2+\Zu_3,\qquad \alpha_1=\Zu_2-\Zu_3.
\end{equation}
All sign choices are independent except for the displayed $\typ E_8$
product condition.
\end{namedresult}

\begin{namedresult}{Lemma}{(the even lattice)}\label{lem:even-lattice}
\lean{RootSetModel.E₈_lattice_even}\lean{RootSetModel.E₈_eq_norm_two}
\uses{def:rem-exceptional,facts:inner}
\assume{$L\subset\mathbb R^8$ is the set in \eqref{eq:E-lattice}.}
\prove{$L$ is an additive subgroup, all its inner products are integers,
all its squared norms are even integers, and its squared-norm-2 vectors
are exactly $\Phi(\typ E_8)$. Every nonzero element has squared norm at least 2.}
\end{namedresult}
\begin{proof}
\pfsketch\ Describe two cosets of the integer even-sum subgroup and
minimize the sum of the eight coordinate squares.
\pflet{$L_0=\{\Zx\in\mathbb Z^8:\sum_i x(i)\in2\mathbb Z\}$ and
$\Zv=\tfrac12\sum_i\Zu_i$.}
\step{relcosets}{$L=L_0\cup(L_0+\Zv)$ and $2\Zv\in L_0$.}
\begin{proof}\pf\ Subtracting $\Zv$ changes the coordinate sum by 4;
$2\Zv$ has integer coordinates and sum 8. ~\qed\end{proof}
\step{relgroup}{$L$ is an additive subgroup.}
\begin{proof}
\step{rel00}{$L_0+L_0\subseteq L_0$.}
\begin{proof}\pf\ Addition preserves integer coordinates and even sums. ~\qed\end{proof}
\step{rel01}{$L_0+(L_0+\Zv)\subseteq L_0+\Zv$.}
\begin{proof}\pf\ Translate the inclusion \stepref{rel00} by $\Zv$.~\qed\end{proof}
\step{rel11}{$(L_0+\Zv)+(L_0+\Zv)\subseteq L_0$.}
\begin{proof}\pf\ Use \stepref{rel00} and $2\Zv\in L_0$ from
\stepref{relcosets}.~\qed\end{proof}
\step{relneg0}{$-L_0=L_0$.}
\begin{proof}\pf\ Negation preserves integer coordinates and even sums. ~\qed\end{proof}
\step{relneg1}{$-(L_0+\Zv)=L_0+\Zv$.}
\begin{proof}\pf\ Write $-\Zx-\Zv=(-\Zx-2\Zv)+\Zv$ using
\stepref{relcosets}, \stepref{relneg0}.~\qed\end{proof}
\qedstep
\begin{proof}\pf\ The two-coset cases \stepref{rel00}--\stepref{relneg1},
commutativity of addition, and $0\in L_0$ give the subgroup laws. ~\qed\end{proof}
\end{proof}
\step{relintegral}{All inner products in $L$ are integers.}
\begin{proof}
\step{relint00}{For $\Zx,\Zy\in L_0$, $(\Zx,\Zy)\in\mathbb Z$.}
\begin{proof}\pf\ Sum products of their integer coordinates. ~\qed\end{proof}
\step{relintv}{For $\Zx\in L_0$, $(\Zx,\Zv)=\sum_i x(i)/2\in\mathbb Z$.}
\begin{proof}\pf\ Use the even-sum definition of $L_0$.~\qed\end{proof}
\step{relintvv}{$\rsq\Zv=2\in\mathbb Z$.}
\begin{proof}\pf\ Sum eight copies of $1/4$.~\qed\end{proof}
\qedstep
\begin{proof}\pf\ Expand
$(\Zx+a\Zv,\Zy+b\Zv)$ for $a,b\in\{0,1\}$ using
\stepref{relint00}, \stepref{relintv}, \stepref{relintvv}.~\qed\end{proof}
\end{proof}
\step{releven}{Every squared norm in $L$ is an even integer.}
\begin{proof}
\step{releven0}{For $\Zx\in L_0$, $\sum_i x(i)^2$ is even.}
\begin{proof}\pf\ Each integer satisfies $x(i)^2\equiv x(i)\pmod2$,
and $\sum_i x(i)$ is even. ~\qed\end{proof}
\step{releven1}{For $\Zx\in L_0$,
$\|\Zx+\Zv\|^2=\sum_i x(i)(x(i)+1)+2$ is even.}
\begin{proof}\pf\ Expand the eight coordinate squares; a product of
consecutive integers is even. ~\qed\end{proof}
\qedstep
\begin{proof}\pf\ Use the coset partition \stepref{relcosets} with
\stepref{releven0}, \stepref{releven1}.~\qed\end{proof}
\end{proof}
\step{relminimum}{A nonzero element of $L$ has squared norm at least 2.}
\begin{proof}\pf\ A positive even integer from \stepref{releven}
is at least 2, using definiteness in \ref{facts:inner}.~\qed\end{proof}
\step{relroots0}{The squared-norm-2 elements of $L_0$ are precisely
the two-coordinate roots $\pm\Zu_i\pm\Zu_j$, $i<j$.}
\begin{proof}
\step{relintbound}{Every nonzero integer coordinate contributes at least
1, and any coordinate of absolute value at least 2 contributes at least 4.}
\begin{proof}\pf\ Square the corresponding integer inequalities. ~\qed\end{proof}
\step{relintsupport}{A sum of integer squares equal to 2 has exactly
two nonzero coordinates, both of absolute value 1.}
\begin{proof}\pf\ Apply \stepref{relintbound} to the finite sum of squares. ~\qed\end{proof}
\qedstep
\begin{proof}\pf\ The patterns in \stepref{relintsupport} have even
coordinate sum and squared norm 2 by substitution. ~\qed\end{proof}
\end{proof}
\step{relroots1}{The squared-norm-2 elements in $L_0+\Zv$ are precisely
the half-coordinate roots with an even number of minus signs.}
\begin{proof}
\step{relhalfbound}{Each of their eight absolute coordinates is at least $1/2$.}
\begin{proof}\pf\ Each coordinate is an integer plus $1/2$.~\qed\end{proof}
\step{relhalfequal}{Squared norm 2 forces all absolute coordinates to be $1/2$.}
\begin{proof}\pf\ The eight lower bounds $1/4$ from
\stepref{relhalfbound} already sum to 2. ~\qed\end{proof}
\step{relhalfparity}{For $t$ minus signs, the coordinate sum is $4-t$.}
\begin{proof}\pf\ Sum $8-t$ copies of $1/2$ and $t$ copies of $-1/2$.~\qed\end{proof}
\qedstep
\begin{proof}\pf\ The sum in \stepref{relhalfparity} is even exactly for
even $t$; all these patterns have squared norm 2 by \stepref{relhalfequal}.~\qed\end{proof}
\end{proof}
\qedstep
\begin{proof}\pf\ Combine \stepref{relgroup}, \stepref{relintegral},
\stepref{releven}, \stepref{relminimum}, \stepref{relroots0},
\stepref{relroots1} with the root list in Definition~\ref{def:rem-exceptional}.~\qed\end{proof}
\end{proof}

\begin{namedresult}{Lemma}{(minimal vectors and an obtuse basis)}\label{lem:minimal-base}
\lean{RootSetModel.coefficients_one_sign_of_minimal}
\uses{facts:inner,def:rem-bases}
\assume{\begin{pfenum}
\item $L$ is an additive subgroup of a finite-dimensional real inner
 product space, $b>0$, and every nonzero $\Zx\in L$ has $\rsq\Zx\geq b$.
\item $(\alpha_i)_{i<m}$ is a real basis generating $L$ over $\mathbb Z$,
 and $(\alpha_i,\alpha_j)\leq0$ for $i\ne j$.
\end{pfenum}}
\prove{Every element of $L$ of squared norm $b$ has integer basis
coefficients all nonnegative or all nonpositive.}
\end{namedresult}
\begin{proof}
\pfsketch\ Split a coefficient vector into disjoint positive and negative parts.
\pflet{For $\Zx=\sum_i c(i)\alpha_i\in L$ with $\rsq\Zx=b$, put
$\Zv=\sum_i\max(c(i),0)\alpha_i$ and
$\Zw=\sum_i\max(-c(i),0)\alpha_i$.}
\step{rmbsplit}{$\Zx=\Zv-\Zw$, with $\Zv,\Zw\in L$.}
\begin{proof}\pf\ Use the integer generation in \asm[2]{0} and
$c=\max(c,0)-\max(-c,0)$.~\qed\end{proof}
\step{rmbcross}{$(\Zv,\Zw)\leq0$.}
\begin{proof}\pf\ Expand the inner product; same-index coefficients
have product zero, and all other summands are nonpositive by \asm[2]{0}.~\qed\end{proof}
\step{rmbone}{At least one of $\Zv,\Zw$ is zero.}
\begin{proof}\pf\ If both were nonzero, \asm[1]{0}, \stepref{rmbsplit},
\stepref{rmbcross} would give
$b=\rsq\Zv+\rsq\Zw-2(\Zv,\Zw)\geq2b>b$.~\qed\end{proof}
\qedstep
\begin{proof}\pf\ Independence in \asm[2]{0} makes every coefficient
of the zero part in \stepref{rmbone} vanish. ~\qed\end{proof}
\end{proof}

\begin{namedresult}{Lemma}{(the three E systems)}\label{lem:E}
\lean{RootSetModel.E_isRootSet}\lean{RootSetModel.E_base}
\lean{RootSetModel.E_cartanMatrix}
\uses{def:rem-exceptional,def:rem-matrices,lem:even-lattice,lem:minimal-base,lem:rem-reflection,facts:inner,def:rem-bases,def:roots}
\assume{$j\in\{6,7,8\}$, with $L$, $V(\typ E_j)$, $\Phi(\typ E_j)$,
and the first $j$ vectors of \eqref{eq:E-base} as in
Definition~\ref{def:rem-exceptional}.}
\prove{$\Phi(\typ E_j)$ is a Root System in $V(\typ E_j)$; those $j$
vectors are a root base and their Cartan matrix is
\path{CartanMatrix.E} at rank $j$. They are also an integer basis for
$L\cap V(\typ E_j)$, and $\dim V(\typ E_j)=j$.}
\end{namedresult}
\begin{proof}
\pfsketch\ Invert the coordinate matrix, then apply the minimum-norm
lemma in each subspace.
\step{rEcoefficient}{The coefficients of every $\Zx\in\mathbb R^8$ are
\begin{equation}\label{eq:E-coeff}
\begin{split}
c(0)&=2x(8),\qquad c(1)=\tfrac12\sum_{i=1}^8x(i)+2x(8),\\
c(2)&=c(1)+x(8)-x(1),\\
c(k)&=\sum_{i=k}^7x(i)+(8-k)x(8)\quad(3\leq k\leq7).
\end{split}
\end{equation}
In particular these formulas are the inverse of the coordinate map.}
\begin{proof}
\step{rEcoord8}{The eighth coordinate of $\sum_i c(i)\alpha_i$ is $x(8)$.}
\begin{proof}\pf\ It is $c(0)/2$ by \eqref{eq:E-base}; use \eqref{eq:E-coeff}.~\qed\end{proof}
\step{rEcoord1}{Its first coordinate is $x(1)$.}
\begin{proof}\pf\ Substitute in $c(0)/2+c(1)-c(2)$.~\qed\end{proof}
\step{rEcoord2}{Its second coordinate is $x(2)$.}
\begin{proof}\pf\ Substitute in $-c(0)/2+c(1)+c(2)-c(3)$.~\qed\end{proof}
\step{rEcoordmiddle}{For $3\leq i\leq6$, its $i$th coordinate is $x(i)$.}
\begin{proof}\pf\ Telescope $-c(0)/2+c(i)-c(i+1)$ using \eqref{eq:E-coeff}.~\qed\end{proof}
\step{rEcoord7}{Its seventh coordinate is $x(7)$.}
\begin{proof}\pf\ Substitute in $-c(0)/2+c(7)$.~\qed\end{proof}
\step{rEunique}{These coefficients are unique.}
\begin{proof}\pf\ The same coordinate equations solve successively for
$c(0),c(7),c(6),c(5),c(4),c(3)$ and then $c(1),c(2)$.~\qed\end{proof}
\qedstep
\begin{proof}\pf\ Combine the exhaustive coordinate checks
\stepref{rEcoord8}--\stepref{rEcoord7} and \stepref{rEunique}.~\qed\end{proof}
\end{proof}
\step{rElattice}{For $\Zx\in L$, all coefficients in \eqref{eq:E-coeff}
are integers.}
\begin{proof}
\step{rEint0}{$c(0)$ is an integer.}
\begin{proof}\pf\ Twice an integer or half-integer is integral. ~\qed\end{proof}
\step{rEint1}{$c(1)$ is an integer.}
\begin{proof}\pf\ Half the even coordinate sum is integral, and so is
$2x(8)$ by \stepref{rEint0}.~\qed\end{proof}
\step{rEint2}{$c(2)$ is an integer.}
\begin{proof}\pf\ Use \stepref{rEint1} and the integer difference
$x(8)-x(1)$ of two coordinates in the same coset. ~\qed\end{proof}
\step{rEintk}{For $3\leq k\leq7$, $c(k)$ is an integer.}
\begin{proof}\pf\ Write it as $\sum_{i=k}^7(x(i)+x(8))$;
each summand is an integer in either coset of \eqref{eq:E-lattice}.~\qed\end{proof}
\qedstep
\begin{proof}\pf\ The index cases \stepref{rEint0}--\stepref{rEintk}
are exhaustive. ~\qed\end{proof}
\end{proof}
\step{rEsubspaces}{All displayed simple roots lie in $L$ and have squared
norm 2. The first $j$ form a real basis of $V(\typ E_j)$ and generate
$L\cap V(\typ E_j)$ over $\mathbb Z$.}
\begin{proof}
\step{rEseven}{$x(7)+x(8)=0$ if and only if $c(7)=0$.}
\begin{proof}\pf\ Substitute $k=7$ in \eqref{eq:E-coeff}.~\qed\end{proof}
\step{rEsix}{The two $\typ E_6$ subspace equations hold if and only if
$c(6)=c(7)=0$.}
\begin{proof}\pf\ Use $c(6)=x(6)+x(7)+2x(8)$ and
\stepref{rEseven}.~\qed\end{proof}
\step{rEsimple}{Every $\alpha_i$ lies in $L$ and has squared norm 2.}
\begin{proof}\pf\ For $\alpha_0$, its six minus signs and eight
half-coordinates satisfy \eqref{eq:E-lattice}; each other vector has
two unit coordinates and even sum. ~\qed\end{proof}
\qedstep
\begin{proof}\pf\ Restrict the inverse bijection \stepref{rEcoefficient}
by \stepref{rEseven}, \stepref{rEsix}; use \stepref{rElattice},
\stepref{rEsimple}, and subgroup closure for integer generation. ~\qed\end{proof}
\end{proof}
\step{rEgram}{The simple Gram matrix is $\typ E_j$.}
\begin{proof}
\step{rEgram01}{$(\alpha_0,\alpha_1)=0$.}
\begin{proof}\pf\ The contributions in coordinates 1 and 2 are
$1/2-1/2=0$.~\qed\end{proof}
\step{rEgram02}{$(\alpha_0,\alpha_2)=-1$.}
\begin{proof}\pf\ The contributions are $-1/2-1/2=-1$.~\qed\end{proof}
\step{rEgram0k}{For $3\leq k\leq7$, $(\alpha_0,\alpha_k)=0$.}
\begin{proof}\pf\ The two coefficients of $\alpha_0$ are both $-1/2$
and cancel in the difference $\alpha_k$.~\qed\end{proof}
\step{rEgram12}{$(\alpha_1,\alpha_2)=0$.}
\begin{proof}\pf\ Expand $(\Zu_1+\Zu_2,\Zu_2-\Zu_1)=1-1$.~\qed\end{proof}
\step{rEgram13}{$(\alpha_1,\alpha_3)=-1$.}
\begin{proof}\pf\ Expand $(\Zu_1+\Zu_2,\Zu_3-\Zu_2)=-1$.~\qed\end{proof}
\step{rEgram1k}{For $k\geq4$, $(\alpha_1,\alpha_k)=0$.}
\begin{proof}\pf\ The supports are disjoint. ~\qed\end{proof}
\step{rEgramjk}{For distinct $i,k\geq2$, the product is $-1$ when
$|i-k|=1$ and zero otherwise.}
\begin{proof}\pf\ Expand $(\Zu_i-\Zu_{i-1},\Zu_k-\Zu_{k-1})$.~\qed\end{proof}
\qedstep
\begin{proof}\pf\ Combine \stepref{rEgram01}--\stepref{rEgramjk},
symmetry, and squared norms 2 from \stepref{rEsubspaces} with
Definition~\ref{def:rem-matrices}.~\qed\end{proof}
\end{proof}
\step{rEbase}{Every norm-2 vector in $L\cap V(\typ E_j)$ has integer
coefficients of one sign.}
\begin{proof}\pf\ Apply Lemma~\ref{lem:minimal-base} with $b=2$, using
Lemma~\ref{lem:even-lattice}, \stepref{rEsubspaces}, \stepref{rEgram}.~\qed\end{proof}
\step{rEreflection}{For $\alpha,\beta\in\Phi(\typ E_j)$,
$s_\alpha\beta\in\Phi(\typ E_j)$.}
\begin{proof}
\step{rEreflectL}{$s_\alpha\beta=\beta-(\alpha,\beta)\alpha\in L$.}
\begin{proof}\pf\ Use integer inner products and subgroup closure in
Lemma~\ref{lem:even-lattice}, since $\rsq\alpha=2$.~\qed\end{proof}
\step{rEreflectV}{$s_\alpha\beta\in V(\typ E_j)$.}
\begin{proof}\pf\ The reflection formula is a linear combination of
$\alpha,\beta$ in that subspace. ~\qed\end{proof}
\step{rEreflectnorm}{$\|s_\alpha\beta\|^2=2$.}
\begin{proof}\pf\ Apply Lemma~\ref{lem:rem-reflection}.~\qed\end{proof}
\qedstep
\begin{proof}\pf\ Lemma~\ref{lem:even-lattice} identifies the vectors in
\stepref{rEreflectL}--\stepref{rEreflectnorm} with the stated root list. ~\qed\end{proof}
\end{proof}
\step{rErootset}{The displayed set is a Root System.}
\begin{proof}
\step{rEfinite}{It is finite and contains no zero.}
\begin{proof}\pf\ Its finite coordinate patterns all have squared norm 2
by Lemma~\ref{lem:even-lattice}.~\qed\end{proof}
\step{rEreduced}{It is reduced.}
\begin{proof}\pf\ Equal positive squared norms force a proportionality
scalar to satisfy $t^2=1$.~\qed\end{proof}
\step{rEintegral}{Its pairings are integers.}
\begin{proof}\pf\ Equation~\eqref{eq:pairing} equals the integer inner
product of Lemma~\ref{lem:even-lattice}.~\qed\end{proof}
\qedstep
\begin{proof}\pf\ Assemble Definition~\ref{def:roots} using
\stepref{rEfinite}, \stepref{rEreduced}, \stepref{rEintegral},
\stepref{rEreflection}, and the spanning roots of \stepref{rEsubspaces}.~\qed\end{proof}
\end{proof}
\qedstep
\begin{proof}\pf\ Combine \stepref{rErootset}, \stepref{rEsubspaces},
\stepref{rEbase}, \stepref{rEgram}, and denominator 2 in \eqref{eq:pairing}.~\qed\end{proof}
\end{proof}

\begin{namedresult}{Lemma}{(the F system)}\label{lem:F}
\lean{RootSetModel.F₄_isRootSet}\lean{RootSetModel.F₄_base}
\lean{RootSetModel.F₄_cartanMatrix}
\uses{def:rem-exceptional,def:rem-matrices,lem:rem-coordinate-reflections,lem:rem-reflection,facts:inner,def:rem-bases}
\assume{$\Phi(\typ F_4)\subset\mathbb R^4$ and $(\alpha_i)_{i<4}$ are
the set and family in \eqref{eq:F-roots}, \eqref{eq:F-base}.}
\prove{This is a Root System with the displayed root base and Cartan
matrix $\typ F_4$. Its first two simple roots have squared length 2
and its last two squared length 1.}
\end{namedresult}
\begin{proof}
\pfsketch\ Only half-coordinate reflections need a new closure check;
an inverse coordinate formula proves the base property.
\step{rFbasic}{The set is finite, nonzero, reduced, spans $\mathbb R^4$,
and its two-coordinate roots have squared length 2, all others length squared 1.}
\begin{proof}
\step{rFfinite}{The set is finite and has no zero.}
\begin{proof}\pf\ Its three finite coordinate patterns have a nonzero
coordinate in \eqref{eq:F-roots}.~\qed\end{proof}
\step{rFlength}{The squared lengths are $1+1=2$, $1$, and $4/4=1$.}
\begin{proof}\pf\ Sum coordinate squares in the three patterns. ~\qed\end{proof}
\step{rFreduced}{Proportional roots differ only by sign.}
\begin{proof}\pf\ Support sizes 2, 1, 4 separate the patterns; within
each pattern the fixed absolute nonzero coordinates force $|t|=1$.~\qed\end{proof}
\step{rFspan}{The roots span $\mathbb R^4$.}
\begin{proof}\pf\ The axis pattern contains the standard basis. ~\qed\end{proof}
\qedstep
\begin{proof}\pf\ Combine \stepref{rFfinite}--\stepref{rFspan}.~\qed\end{proof}
\end{proof}
\step{rFreflection}{All root reflections preserve the set.}
\begin{proof}
\step{rFintegerref}{Reflections in an axis or two-coordinate root preserve it.}
\begin{proof}\pf\ Apply Lemma~\ref{lem:rem-coordinate-reflections} to
the three complete signed patterns in \eqref{eq:F-roots}.~\qed\end{proof}
\pflet{$\Zv=(\Zu_1+\Zu_2+\Zu_3+\Zu_4)/2$.}
\step{rFvformula}{$s_\Zv\Zx=\Zx-(\sum_i x(i))\Zv$.}
\begin{proof}\pf\ Substitute $\rsq\Zv=1$ in the reflection formula. ~\qed\end{proof}
\step{rFvaxis}{For an axis root $\epsilon\Zu_i$, its image is
$\epsilon(\Zu_i-\Zv)$, a half-coordinate root.}
\begin{proof}\pf\ Substitute its coordinate sum $\epsilon$ in
\stepref{rFvformula}.~\qed\end{proof}
\step{rFvopposite}{For $\epsilon(\Zu_i-\Zu_j)$ its image is itself.}
\begin{proof}\pf\ Its coordinate sum is zero in \stepref{rFvformula}.~\qed\end{proof}
\step{rFvequal}{For $\epsilon(\Zu_i+\Zu_j)$ its image is
$-\epsilon(\Zu_k+\Zu_l)$ on the complementary pair.}
\begin{proof}\pf\ Substitute its coordinate sum $2\epsilon$ in
\stepref{rFvformula}.~\qed\end{proof}
\step{rFvhalf0}{A half-coordinate root with zero plus signs maps to its negative.}
\begin{proof}\pf\ It is $-\Zv$, with sum $-2$, in \stepref{rFvformula}.~\qed\end{proof}
\step{rFvhalf1}{A half-coordinate root with one plus sign maps to the
positive axis at that sign.}
\begin{proof}\pf\ Its sum is $-1$, so \stepref{rFvformula} adds $\Zv$.~\qed\end{proof}
\step{rFvhalf2}{A half-coordinate root with two plus signs is fixed.}
\begin{proof}\pf\ Its sum is zero in \stepref{rFvformula}.~\qed\end{proof}
\step{rFvhalf3}{A half-coordinate root with three plus signs maps to
the negative axis at its minus sign.}
\begin{proof}\pf\ Its sum is 1, so \stepref{rFvformula} subtracts $\Zv$.~\qed\end{proof}
\step{rFvhalf4}{A half-coordinate root with four plus signs maps to its negative.}
\begin{proof}\pf\ It is $\Zv$, with sum 2, in \stepref{rFvformula}.~\qed\end{proof}
\step{rFvclosed}{$s_\Zv$ preserves the whole set.}
\begin{proof}\pf\ The cases \stepref{rFvaxis}--\stepref{rFvhalf4}
exhaust \eqref{eq:F-roots}.~\qed\end{proof}
\step{rFhalfref}{Reflection in any half-coordinate root preserves the set.}
\begin{proof}
\step{rFsignmap}{A diagonal sign isometry taking $\Zv$ to that root
permutes $\Phi(\typ F_4)$.}
\begin{proof}\pf\ Apply the signed-pattern assertion of
Lemma~\ref{lem:rem-coordinate-reflections}.~\qed\end{proof}
\qedstep
\begin{proof}\pf\ Conjugate \stepref{rFvclosed} by \stepref{rFsignmap}
using Lemma~\ref{lem:rem-reflection}.~\qed\end{proof}
\end{proof}
\qedstep
\begin{proof}\pf\ Combine \stepref{rFintegerref}, \stepref{rFhalfref}
over the three root patterns. ~\qed\end{proof}
\end{proof}
\step{rFintegrality}{Every root pairing is an integer.}
\begin{proof}
\step{rFpairlongint}{Pairing an integer root against a long root is integral.}
\begin{proof}\pf\ Its value is an inner product of integer coordinates
by \eqref{eq:pairing} and \stepref{rFbasic}.~\qed\end{proof}
\step{rFpairlonghalf}{Pairing a half-coordinate root against a long
root is integral.}
\begin{proof}\pf\ The value is a sum of two signed halves by \eqref{eq:pairing}.~\qed\end{proof}
\step{rFpairaxis}{Pairing any root against an axis root is integral.}
\begin{proof}\pf\ Its value is twice a signed coordinate by
\eqref{eq:pairing}; all coordinates are half-integers. ~\qed\end{proof}
\step{rFpairhalfint}{Pairing an integer root against a half-coordinate
root is integral.}
\begin{proof}\pf\ Twice the inner product is a signed sum of the
integer root's coordinates. ~\qed\end{proof}
\step{rFpairhalfhalf}{Pairing two half-coordinate roots is integral.}
\begin{proof}\pf\ If their signs differ in $t$ places, the value is
$2(4-2t)/4=2-t$.~\qed\end{proof}
\qedstep
\begin{proof}\pf\ The denominator and numerator patterns in
\stepref{rFpairlongint}--\stepref{rFpairhalfhalf} are exhaustive. ~\qed\end{proof}
\end{proof}
\step{rFcoeff}{The inverse coefficients in the displayed family are
\begin{equation}\label{eq:F-coeff}
\begin{aligned}
c(0)&=x(1)+x(2),&c(1)&=2x(1)+x(2)+x(3),\\
c(2)&=3x(1)+x(2)+x(3)+x(4),&c(3)&=2x(1).
\end{aligned}
\end{equation}}
\begin{proof}
\step{rFcoord1}{The first coordinate of the combination is $c(3)/2=x(1)$.}
\begin{proof}\pf\ Substitute \eqref{eq:F-base}, \eqref{eq:F-coeff}.~\qed\end{proof}
\step{rFcoord2}{The second is $c(0)-c(3)/2=x(2)$.}
\begin{proof}\pf\ Substitute \eqref{eq:F-coeff}.~\qed\end{proof}
\step{rFcoord3}{The third is $-c(0)+c(1)-c(3)/2=x(3)$.}
\begin{proof}\pf\ Substitute \eqref{eq:F-coeff}.~\qed\end{proof}
\step{rFcoord4}{The fourth is $-c(1)+c(2)-c(3)/2=x(4)$.}
\begin{proof}\pf\ Substitute \eqref{eq:F-coeff}.~\qed\end{proof}
\qedstep
\begin{proof}\pf\ These four triangular equations uniquely solve for
$c(3),c(0),c(1),c(2)$, giving the asserted inverse. ~\qed\end{proof}
\end{proof}
\step{rFbase}{The displayed family is a root base.}
\begin{proof}
\step{rFbasehalf}{Orient a half-coordinate root to have $x(1)=1/2$.
Its coefficients are integers in $[0,1],[0,2],[0,3],\{1\}$.}
\begin{proof}\pf\ Substitute the three independent signed halves in
\eqref{eq:F-coeff}.~\qed\end{proof}
\step{rFbasefirst}{For an integer root with $x(1)\ne0$, orient it with
$x(1)=1$; all four coefficients are nonnegative integers.}
\begin{proof}\pf\ In \eqref{eq:F-coeff} at most one of $x(2),x(3),x(4)$
is nonzero, and it is $1$ or $-1$.~\qed\end{proof}
\step{rFbasezero}{For $x(1)=0$, orient the first nonzero coordinate
among 2,3,4 positively; the four coefficients are nonnegative integers.}
\begin{proof}
\step{rFprefix}{Every prefix sum of coordinates 2,3,4 is nonnegative.}
\begin{proof}\pf\ It begins with zeros, then 1, and has at most one
further nonzero term, bounded below by $-1$.~\qed\end{proof}
\qedstep
\begin{proof}\pf\ For $x(1)=0$, \eqref{eq:F-coeff} consists of those
three prefix sums and zero; apply \stepref{rFprefix}.~\qed\end{proof}
\end{proof}
\step{rFbaseindependent}{The four simple roots are linearly independent.}
\begin{proof}\pf\ Their coordinate map has the inverse \stepref{rFcoeff}.~\qed\end{proof}
\qedstep
\begin{proof}\pf\ The root patterns in \stepref{rFbasehalf},
\stepref{rFbasefirst}, \stepref{rFbasezero} exhaust the roots up to sign;
use \stepref{rFbaseindependent} and their membership in \eqref{eq:F-roots}.~\qed\end{proof}
\end{proof}
\step{rFmatrix}{The simple squared lengths are $(2,2,1,1)$ and their
Cartan matrix is $\typ F_4$.}
\begin{proof}
\step{rFgram01}{$(\alpha_0,\alpha_1)=-1$.}
\begin{proof}\pf\ Expand the product of $\Zu_2-\Zu_3$ and $\Zu_3-\Zu_4$.~\qed\end{proof}
\step{rFgram12}{$(\alpha_1,\alpha_2)=-1$.}
\begin{proof}\pf\ Expand the product of $\Zu_3-\Zu_4$ and $\Zu_4$.~\qed\end{proof}
\step{rFgram23}{$(\alpha_2,\alpha_3)=-1/2$.}
\begin{proof}\pf\ Read the fourth coordinate of $\alpha_3$.~\qed\end{proof}
\step{rFgram02}{$(\alpha_0,\alpha_2)=0$.}
\begin{proof}\pf\ Their supports are disjoint. ~\qed\end{proof}
\step{rFgram03}{$(\alpha_0,\alpha_3)=0$.}
\begin{proof}\pf\ The equal second and third coordinates of $\alpha_3$
cancel. ~\qed\end{proof}
\step{rFgram13}{$(\alpha_1,\alpha_3)=0$.}
\begin{proof}\pf\ The equal third and fourth coordinates of $\alpha_3$
cancel. ~\qed\end{proof}
\qedstep
\begin{proof}\pf\ Use \stepref{rFbasic} for diagonal entries and
\stepref{rFgram01}--\stepref{rFgram13} in \eqref{eq:pairing};
compare Definition~\ref{def:rem-matrices}.~\qed\end{proof}
\end{proof}
\qedstep
\begin{proof}\pf\ Assemble the Root System conditions from
\stepref{rFbasic}, \stepref{rFreflection}, \stepref{rFintegrality};
add \stepref{rFbase}, \stepref{rFmatrix}.~\qed\end{proof}
\end{proof}

\begin{namedresult}{Lemma}{(the G system)}\label{lem:G}
\lean{RootSetModel.G₂_isRootSet}\lean{RootSetModel.G₂_base}
\lean{RootSetModel.G₂_cartanMatrix}
\uses{def:rem-exceptional,def:rem-matrices,lem:rem-coordinate-reflections,lem:rem-reflection,facts:inner,def:rem-bases,def:roots}
\assume{$V$ is the zero-sum plane in $\mathbb R^3$, with root set
\eqref{eq:G-roots} and family \eqref{eq:G-base}.}
\prove{This is a Root System with the displayed root base and Cartan
matrix $\typ G_2$. Its simple roots 0 and 1 have squared lengths 6 and 2,
respectively.}
\end{namedresult}
\begin{proof}
\pfsketch\ A long-root reflection is a permutation followed by global
negation. Six coefficient pairs give one root from every sign pair.
\step{rGbasic}{There are six short roots of squared length 2 and six long
roots of squared length 6; they are finite, nonzero and reduced and span $V$.}
\begin{proof}
\step{rGshort}{The six distinct-coordinate differences have squared norm 2.}
\begin{proof}\pf\ There are $3\cdot2$ ordered pairs and each squared
norm is $1+1$.~\qed\end{proof}
\step{rGlong}{The six signed roots $(2,-1,-1)$ under coordinate
permutations have squared norm 6.}
\begin{proof}\pf\ Choose the doubled coordinate and a global sign;
their squared norm is $4+1+1$.~\qed\end{proof}
\step{rGreduced}{Proportional listed vectors differ by sign.}
\begin{proof}\pf\ Short and long supports have different sizes; in
either pattern the ratios of nonzero coordinate magnitudes are fixed. ~\qed\end{proof}
\step{rGspan}{The differences span the zero-sum plane.}
\begin{proof}\pf\ A vector in that plane is
$x(1)(\Zu_1-\Zu_3)+x(2)(\Zu_2-\Zu_3)$.~\qed\end{proof}
\qedstep
\begin{proof}\pf\ Combine \stepref{rGshort}--\stepref{rGspan}.~\qed\end{proof}
\end{proof}
\pflet{$\beta=2\Zu_1-\Zu_2-\Zu_3$.}
\step{rGlongformula}{For $\Zx\in V$,
$(\Zx,\beta)=3x(1)$ and
$s_\beta\Zx=(-x(1),-x(3),-x(2))$.}
\begin{proof}
\step{rGinner}{$2x(1)-x(2)-x(3)=3x(1)$.}
\begin{proof}\pf\ Substitute the zero-sum equation of $V$.~\qed\end{proof}
\qedstep
\begin{proof}\pf\ Substitute \stepref{rGinner}, $\rsq\beta=6$ in the
reflection formula, and use $x(1)+x(2)+x(3)=0$.~\qed\end{proof}
\end{proof}
\step{rGreflection}{The set is closed under every root reflection.}
\begin{proof}
\step{rGshortref}{Each short reflection preserves it.}
\begin{proof}\pf\ It is a coordinate interchange by
Lemma~\ref{lem:rem-coordinate-reflections}; both lists are permutation-invariant. ~\qed\end{proof}
\step{rGbetaref}{$s_\beta$ preserves it.}
\begin{proof}\pf\ Its formula in \stepref{rGlongformula} is a
coordinate interchange followed by global negation, preserving both lists. ~\qed\end{proof}
\step{rGlongref}{Each long reflection preserves it.}
\begin{proof}\pf\ Every long root is a signed coordinate permutation
of $\beta$ using only a global sign; conjugate \stepref{rGbetaref}
by Lemma~\ref{lem:rem-reflection}.~\qed\end{proof}
\qedstep
\begin{proof}\pf\ The short and long cases \stepref{rGshortref},
\stepref{rGlongref} exhaust \eqref{eq:G-roots}.~\qed\end{proof}
\end{proof}
\step{rGinteger}{All pairings are integral.}
\begin{proof}
\step{rGshortpair}{Against a short root the pairing is an integer inner product.}
\begin{proof}\pf\ All root coordinates are integers and its norm squared
is 2 by \stepref{rGbasic}.~\qed\end{proof}
\step{rGbetapair}{Against $\beta$ the pairing is $x(1)\in\mathbb Z$.}
\begin{proof}\pf\ Divide twice the inner product in
\stepref{rGlongformula} by 6. ~\qed\end{proof}
\step{rGlongpair}{Against every long root it is integral.}
\begin{proof}\pf\ Transport \stepref{rGbetapair} by the global sign
and coordinate permutation carrying $\beta$ to that root. ~\qed\end{proof}
\qedstep
\begin{proof}\pf\ Combine \stepref{rGshortpair}, \stepref{rGlongpair}.~\qed\end{proof}
\end{proof}
\step{rGcoeff}{The six positive coefficient pairs are
\begin{equation}\label{eq:G-coeff}
(0,1),\ (1,1),\ (1,2),\ (1,0),\ (1,3),\ (2,3).
\end{equation}
They give all roots up to sign.}
\begin{proof}
\step{rGc01}{$\alpha_1=\Zu_2-\Zu_3$.}
\begin{proof}\pf\ Use \eqref{eq:G-base}.~\qed\end{proof}
\step{rGc11}{$\alpha_0+\alpha_1=\Zu_1-\Zu_2$.}
\begin{proof}\pf\ Add the coordinates in \eqref{eq:G-base}.~\qed\end{proof}
\step{rGc12}{$\alpha_0+2\alpha_1=\Zu_1-\Zu_3$.}
\begin{proof}\pf\ Add one more $\alpha_1$ to \stepref{rGc11}.~\qed\end{proof}
\step{rGc10}{$\alpha_0=\Zu_1-2\Zu_2+\Zu_3$.}
\begin{proof}\pf\ Use \eqref{eq:G-base}.~\qed\end{proof}
\step{rGc13}{$\alpha_0+3\alpha_1=\Zu_1+\Zu_2-2\Zu_3$.}
\begin{proof}\pf\ Substitute \eqref{eq:G-base}.~\qed\end{proof}
\step{rGc23}{$2\alpha_0+3\alpha_1=2\Zu_1-\Zu_2-\Zu_3$.}
\begin{proof}\pf\ Substitute \eqref{eq:G-base}.~\qed\end{proof}
\qedstep
\begin{proof}\pf\ The first three vectors and their negatives enumerate
the short pattern; the last three and their negatives enumerate the
long pattern of \eqref{eq:G-roots}.~\qed\end{proof}
\end{proof}
\step{rGmatrix}{The simple Gram matrix is
$\left(\begin{smallmatrix}6&-3\\-3&2\end{smallmatrix}\right)$ and
the Cartan matrix is $\typ G_2$.}
\begin{proof}\pf\ Sum coordinate products in \eqref{eq:G-base} and
divide columns by 3 and 1 according to \eqref{eq:pairing}.~\qed\end{proof}
\step{rGindependent}{The two simple roots are independent.}
\begin{proof}\pf\ Their Gram determinant from \stepref{rGmatrix} is
$12-9=3\ne0$.~\qed\end{proof}
\qedstep
\begin{proof}\pf\ Assemble Definition~\ref{def:roots} from
\stepref{rGbasic}, \stepref{rGreflection}, \stepref{rGinteger};
\stepref{rGcoeff}, \stepref{rGindependent} give the base, and
\stepref{rGmatrix} gives its matrix. ~\qed\end{proof}
\end{proof}

\begin{namedresult}{Lemma}{(indecomposability from a base)}\label{lem:rem-base-connected}
\lean{IsRootSet.isIndecomposable_of_connected_base}
\uses{def:roots,def:rem-bases,lem:rem-edge-constant,facts:inner}
\assume{\begin{pfenum}
\item $\Phi$ is a Root System in $V$ with a nonempty root base $\Delta$.
\item The graph on $\Delta$ joining distinct nonorthogonal roots is connected.
\end{pfenum}}
\prove{$\Phi$ is indecomposable.}
\end{namedresult}
\begin{proof}
\pfsketch\ An orthogonal partition puts all simple roots on the same
side, leaving a vector orthogonal to a spanning set on the other side.
\step{rbcpartition}{\assume{$\Phi=P\sqcup Q$ is a partition into
mutually orthogonal parts.}
\prove{$P=\varnothing$ or $Q=\varnothing$.}}
\begin{proof}
\step{rbcside}{Adjacent base roots belong to the same part.}
\begin{proof}\pf\ Roots in different parts are orthogonal by the partition
hypothesis. ~\qed\end{proof}
\step{rbcwhole}{All base roots belong to one part.}
\begin{proof}\pf\ Apply
Lemma~\ref{lem:rem-edge-constant} to the function recording the part
containing a simple root, on the connected nonorthogonality graph
of $\Delta$, using \stepref{rbcside}.~\qed\end{proof}
\step{rbcother}{The other part is empty.}
\begin{proof}
\step{rbcperp}{A root in the other part is orthogonal to every vector in $V$.}
\begin{proof}\pf\ Expand any vector in the spanning base, using
\stepref{rbcwhole} and bilinearity. ~\qed\end{proof}
\qedstep
\begin{proof}\pf\ Pair such a root with itself in \stepref{rbcperp};
definiteness contradicts its nonzero status. ~\qed\end{proof}
\end{proof}
\qedstep
\begin{proof}\pf\ Step \stepref{rbcother} makes one part empty.~\qed\end{proof}
\end{proof}
\qedstep
\begin{proof}\pf\ The nonempty base and \stepref{rbcpartition}
give indecomposability in Definition~\ref{def:roots}.~\qed\end{proof}
\end{proof}

\begin{namedresult}{Lemma}{(the models are indecomposable)}\label{lem:models-irreducible}
\lean{RootSetModel.isIndecomposable}\lean{RootSetModel.Cn_isIndecomposable}
\uses{lem:classical,lem:E,lem:F,lem:G,lem:rem-model-supports,lem:rem-base-connected}
\assume{The model is $\typ A_\ell$ ($\ell\geq1$), $\typ B_\ell$
or $\typ C_\ell$ ($\ell\geq2$), $\typ D_\ell$ ($\ell\geq4$),
$\typ E_j$ ($j=6,7,8$), $\typ F_4$, or $\typ G_2$, in its specified space.}
\prove{It is indecomposable.}
\end{namedresult}
\begin{proof}
\pfsketch\ Apply a single general base-connectivity lemma to each model,
so this proof never reuses local steps for a second Root System.
\step{rmiA}{The $\typ A$ base graph is a connected nonempty path.}
\begin{proof}\pf\ Its Cartan matrix is $\typ A_\ell$ by
Lemma~\ref{lem:classical}; Lemma~\ref{lem:rem-model-supports}
identifies its connected support with the base nonorthogonality graph.~\qed\end{proof}
\step{rmiBC}{Each $\typ B,\typ C$ base graph is a connected nonempty path.}
\begin{proof}\pf\ Use the two matrix identities in
Lemma~\ref{lem:classical} and the path supports of
Lemma~\ref{lem:rem-model-supports}, including $\typ C_2$.~\qed\end{proof}
\step{rmiD}{The $\typ D$ base graph is the connected tree $T(1,1,\ell-3)$.}
\begin{proof}\pf\ Combine Lemma~\ref{lem:classical} with $h_D$ and
connectedness in Lemma~\ref{lem:rem-model-supports}.~\qed\end{proof}
\step{rmiE}{Each $\typ E$ base graph is the connected tree $T(1,2,j-4)$.}
\begin{proof}\pf\ Combine Lemma~\ref{lem:E} with $h_E$ and
connectedness in Lemma~\ref{lem:rem-model-supports}.~\qed\end{proof}
\step{rmiF}{The $\typ F_4$ base graph is a connected four-vertex path.}
\begin{proof}\pf\ Use the Cartan matrix in Lemma~\ref{lem:F}
and its connected support in Lemma~\ref{lem:rem-model-supports}.~\qed\end{proof}
\step{rmiG}{The $\typ G_2$ base graph is a connected single edge.}
\begin{proof}\pf\ Use the Cartan matrix in Lemma~\ref{lem:G}
and its connected support in Lemma~\ref{lem:rem-model-supports}.~\qed\end{proof}
\qedstep
\begin{proof}\pf\ Apply Lemma~\ref{lem:rem-base-connected}
to the nonempty bases and connected graphs in
\stepref{rmiA}--\stepref{rmiG}.~\qed\end{proof}
\end{proof}

\begin{namedresult}{Lemma}{(root, length and line counts)}\label{lem:counts}
\lean{RootSetModel.root_counts}\lean{RootSetModel.length_counts}
\lean{RootSetModel.numLines}\lean{RootSetModel.one_length_iff}
\lean{RootSetModel.lengthRatio}\lean{RootSetModel.maxLengthSq_eq_lengthRatio_mul}
\uses{def:rem-classical,def:rem-exceptional,lem:classical,lem:E,lem:F,lem:G,lem:rem-reflection}
\assume{The coordinate models are in their construction ranges:
$\ell\geq1$ for $\typ A$, $\ell\geq2$ for $\typ B,\typ C$,
$\ell\geq4$ for $\typ D$, and $j\in\{6,7,8\}$ for $\typ E_j$.}
\prove{The counts are those in the following table. The two-length
models have squared-length ratio 2 in types $\typ B,\typ C,\typ F$ and
3 in type $\typ G$; all other models have one length.}
\begin{center}
\begin{tabular}{c|r|r|r|r}
\label{tab:roots}
Type&one length&short&long&lines\\ \hline
$\typ A_\ell$&$\ell(\ell+1)$&&&$\ell(\ell+1)/2$\\
$\typ B_\ell$&&$2\ell$&$2\ell(\ell-1)$&$\ell^2$\\
$\typ C_\ell$&&$2\ell(\ell-1)$&$2\ell$&$\ell^2$\\
$\typ D_\ell$&$2\ell(\ell-1)$&&&$\ell(\ell-1)$\\
$\typ E_6$&72&&&36\\
$\typ E_7$&126&&&63\\
$\typ E_8$&240&&&120\\
$\typ F_4$&&24&24&24\\
$\typ G_2$&&6&6&6
\end{tabular}
\end{center}
\end{namedresult}
\begin{proof}
\pfsketch\ Count disjoint coordinate patterns, then divide by the
two roots on each line.
\step{rctA}{The $\typ A_\ell$ set has $\ell(\ell+1)$ roots.}
\begin{proof}\pf\ The ordered distinct-coordinate pair is recovered
from its coefficients $1,-1$; count $(\ell+1)\ell$ pairs. ~\qed\end{proof}
\step{rctpairs}{There are $4\binom\ell2=2\ell(\ell-1)$ signed
two-coordinate roots.}
\begin{proof}\pf\ Their support and the two signs give a bijection with
two-element coordinate subsets times $\{1,-1\}^2$.~\qed\end{proof}
\step{rctaxes}{There are $2\ell$ axis roots in $\typ B$ and in $\typ C$.}
\begin{proof}\pf\ The coordinate and sign are uniquely recovered from
each such root in Definition~\ref{def:rem-classical}.~\qed\end{proof}
\step{rctE8int}{The $\typ E_8$ integer pattern has 112 roots.}
\begin{proof}\pf\ Evaluate \stepref{rctpairs} at $\ell=8$.~\qed\end{proof}
\step{rctE8half}{Its half-coordinate pattern has 128 roots.}
\begin{proof}\pf\ Choose seven signs freely; their product determines
the eighth in Definition~\ref{def:rem-exceptional}.~\qed\end{proof}
\step{rctE7int}{The $\typ E_7$ integer pattern has $60+2=62$ roots.}
\begin{proof}
\step{rctE7first}{There are 60 with both coordinates among the first six.}
\begin{proof}\pf\ Evaluate \stepref{rctpairs} at $\ell=6$.~\qed\end{proof}
\step{rctE7last}{There are two using coordinates 7 or 8.}
\begin{proof}\pf\ The equation $x(7)+x(8)=0$ forces the pair
$\pm(\Zu_7-\Zu_8)$ in a two-coordinate root. ~\qed\end{proof}
\qedstep
\begin{proof}\pf\ The disjoint cases \stepref{rctE7first},
\stepref{rctE7last} exhaust the integer pattern. ~\qed\end{proof}
\end{proof}
\step{rctE7half}{The $\typ E_7$ half-coordinate pattern has 64 roots.}
\begin{proof}\pf\ Its last two signs are opposite, the first six have
product $-1$, and five of the first six signs and sign 8 are free. ~\qed\end{proof}
\step{rctE6int}{The $\typ E_6$ integer pattern has 40 roots.}
\begin{proof}
\step{rctE6support}{Its nonzero coordinates are among the first five.}
\begin{proof}\pf\ A nonzero coordinate among 6,7,8 would force all
three to be nonzero, contrary to support size 2. ~\qed\end{proof}
\qedstep
\begin{proof}\pf\ Apply \stepref{rctpairs} at $\ell=5$ to
\stepref{rctE6support}.~\qed\end{proof}
\end{proof}
\step{rctE6half}{The $\typ E_6$ half-coordinate pattern has 32 roots.}
\begin{proof}\pf\ The first five signs are free; their product is
sign 8, and signs 6,7 are its negative. ~\qed\end{proof}
\step{rctFlong}{There are 24 long $\typ F_4$ roots.}
\begin{proof}\pf\ Apply \stepref{rctpairs} at $\ell=4$.~\qed\end{proof}
\step{rctFshort}{There are $8+16=24$ short $\typ F_4$ roots.}
\begin{proof}\pf\ The disjoint axis and half-coordinate patterns have
$2\cdot4$ and $2^4$ independent index/sign choices. ~\qed\end{proof}
\step{rctG}{There are six roots of each length in $\typ G_2$.}
\begin{proof}\pf\ Count $3\cdot2$ ordered distinct differences and
$3\cdot2$ choices of doubled coordinate and global sign in \eqref{eq:G-roots}.~\qed\end{proof}
\step{rctlengths}{The squared lengths are 2 for $\typ A,\typ D,\typ E$;
1 and 2 for $\typ B,\typ F$; 2 and 4 for $\typ C$; 2 and 6 for $\typ G$.}
\begin{proof}\pf\ Sum the coordinate squares in
Definitions~\ref{def:rem-classical}, \ref{def:rem-exceptional}.~\qed\end{proof}
\step{rctlines}{Each line contains exactly two roots.}
\begin{proof}\pf\ Reducedness and closure under negation, from the
model lemmas and Lemma~\ref{lem:rem-reflection}, give precisely a sign pair. ~\qed\end{proof}
\qedstep
\begin{proof}\pf\ Add disjoint pattern counts
\stepref{rctA}--\stepref{rctG}, separate by \stepref{rctlengths},
and divide by 2 using \stepref{rctlines}.~\qed\end{proof}
\end{proof}

\subsubsection{The Mathlib bridge and classification}

\begin{namedresult}{Lemma}{(the finite-root-set pairing)}\label{lem:pairing-construction}
\lean{IsRootSet.toRootPairing}\lean{IsRootSet.isIrreducible_iff}
\uses{def:roots,lem:rem-reflection,facts:inner,facts:rem-library,def:rem-bases}
\assume{$\Phi$ is a Root System in a finite-dimensional real inner
product space $V$.}
\prove{It defines a finite reduced crystallographic root pairing $P$
over $\mathbb R$ with root and coroot spaces both $V$, and $P$ is a
root system in Mathlib's sense. Its pairing is \eqref{eq:pairing}.
If $\Phi\ne\varnothing$, its indecomposability is equivalent to
$P$ being Mathlib-irreducible.}
\end{namedresult}
\begin{proof}
\pfsketch\ Construct the structure on the prescribed index set $\Phi$.
For irreducibility a reflection-invariant subspace assigns each root
to itself or to its orthogonal complement.
\pflet{Use root indices $\Phi$, root embedding $\alpha\mapsto\alpha$,
coroot map $\alpha^\vee=2\alpha/\rsq\alpha$, and bilinear form the real inner product.}
\step{rpcperfect}{This bilinear pairing is perfect.}
\begin{proof}
\step{rpcinjective}{The map from $V$ to its dual is injective.}
\begin{proof}\pf\ If $(\Zx-\Zy,\Zz)=0$ for every $\Zz$, choose
$\Zz=\Zx-\Zy$ and apply definiteness in \ref{facts:inner}.~\qed\end{proof}
\step{rpcflip}{The flipped map is injective.}
\begin{proof}\pf\ Rewrite it as the same map by real inner-product
symmetry and use \stepref{rpcinjective}.~\qed\end{proof}
\qedstep
\begin{proof}\pf\ Apply \path{LinearMap.IsPerfPair.of_injective} to
\stepref{rpcinjective}, \stepref{rpcflip}, using finite dimensionality. ~\qed\end{proof}
\end{proof}
\step{rpccoroot}{The coroot map is injective.}
\begin{proof}
\step{rpctwice}{$2\alpha^\vee/\rsq{\alpha^\vee}=\alpha$ for $\alpha\ne0$.}
\begin{proof}\pf\ Substitute $\rsq{\alpha^\vee}=4/\rsq\alpha$
using bilinearity and $\rsq\alpha>0$.~\qed\end{proof}
\qedstep
\begin{proof}\pf\ A map with the left inverse \stepref{rpctwice} is injective. ~\qed\end{proof}
\end{proof}
\step{rpcself}{$(\alpha,\alpha^\vee)=2$.}
\begin{proof}\pf\ Cancel $\rsq\alpha>0$ in the coroot formula. ~\qed\end{proof}
\step{rpcreflection}{The restricted map $s_\alpha$ is a permutation of
the root indices.}
\begin{proof}\pf\ Apply Lemma~\ref{lem:rem-reflection} and reflection
closure in Definition~\ref{def:roots}.~\qed\end{proof}
\step{rpccorootidentity}{$(s_\alpha\beta)^\vee
=\beta^\vee-(\alpha,\beta^\vee)\alpha^\vee$.}
\begin{proof}
\step{rpcnorm}{$(s_\alpha\beta)^\vee=s_\alpha(\beta^\vee)$.}
\begin{proof}\pf\ The orthogonal linear map of
Lemma~\ref{lem:rem-reflection} commutes with normalization by squared norm. ~\qed\end{proof}
\qedstep
\begin{proof}\pf\ Substitute the reflection formula and
$\alpha^\vee=2\alpha/\rsq\alpha$ in \stepref{rpcnorm}.~\qed\end{proof}
\end{proof}
\step{rpcstructure}{These data define a root pairing $P$ with
$P.p(i,j)=2(\alpha_i,\alpha_j)/\rsq{\alpha_j}$.}
\begin{proof}\pf\ Fill the structure fields of \path{RootPairing}
using \stepref{rpcperfect}, root inclusion, \stepref{rpccoroot},
\stepref{rpcself}, \stepref{rpcreflection}, \stepref{rpccorootidentity}
and the defining root reflection identity. ~\qed\end{proof}
\step{rpcreflectionmaps}{The root and coroot reflection at index
$\alpha$ are both $s_\alpha$ on every vector of $V$.}
\begin{proof}\pf\ The root formula subtracts
$(\Zx,\alpha^\vee)\alpha$; the coroot formula subtracts
$(\alpha,\Zx)\alpha^\vee$. Both equal
$2(\alpha,\Zx)\alpha/\rsq\alpha$ by the coroot definition. ~\qed\end{proof}
\step{rpcproperties}{$P$ is finite, crystallographic, reduced and spanning.}
\begin{proof}
\step{rpcfinite}{The index type is finite.}
\begin{proof}\pf\ It is the finite set $\Phi$ of Definition~\ref{def:roots}.~\qed\end{proof}
\step{rpccrystal}{All its pairings are integral.}
\begin{proof}\pf\ Apply crystallographic integrality in
Definition~\ref{def:roots} with the two roots interchanged to
\stepref{rpcstructure}.~\qed\end{proof}
\step{rpcreduced}{A dependent pair of roots is an equal or opposite pair.}
\begin{proof}
\step{rpcdependence}{For a dependent pair $\alpha,\beta$, choose
$(a,b)\ne(0,0)$ with $a\alpha+b\beta=0$.}
\begin{proof}\pf\ Unfold linear dependence for the two-vector family. ~\qed\end{proof}
\step{rpcsecond}{$b\ne0$.}
\begin{proof}\pf\ Otherwise \stepref{rpcdependence} and $\alpha\ne0$
would force $a=0$.~\qed\end{proof}
\step{rpcproportional}{$\beta=(-a/b)\alpha$.}
\begin{proof}\pf\ Solve the relation \stepref{rpcdependence} using
\stepref{rpcsecond}.~\qed\end{proof}
\qedstep
\begin{proof}\pf\ Apply reducedness in Definition~\ref{def:roots}
to the proportionality \stepref{rpcproportional}.~\qed\end{proof}
\end{proof}
\step{rpcrootspan}{The roots span $V$.}
\begin{proof}\pf\ This is the spanning condition in Definition~\ref{def:roots}.~\qed\end{proof}
\step{rpccorootspan}{The coroots span $V$.}
\begin{proof}\pf\ Each root is a nonzero scalar multiple of its coroot,
so \stepref{rpcrootspan} lies in the coroot span. ~\qed\end{proof}
\qedstep
\begin{proof}\pf\ Assemble the Mathlib predicates from
\stepref{rpcfinite}--\stepref{rpccorootspan} using \ref{facts:rem-library}.~\qed\end{proof}
\end{proof}
\step{rpcirreducible}{\assume{$\Phi$ is nonempty and indecomposable.}
\prove{$P$ is irreducible.}}
\begin{proof}
\step{rpcnontrivial}{$V$ is nontrivial.}
\begin{proof}\pf\ A chosen root is nonzero by Definition~\ref{def:roots}.~\qed\end{proof}
\step{rpcinvariant}{Every reflection-invariant subspace $W$ is $0$ or $V$.}
\begin{proof}
\step{rpcalternative}{For every root $\alpha$, either $\alpha\in W$
or $\alpha\in W^\perp$.}
\begin{proof}
\step{rpcdetect}{\assume{some $\Zw\in W$ has $(\Zw,\alpha)\ne0$.}
\prove{$\alpha\in W$.}}
\begin{proof}\pf\ The vector $\Zw-s_\alpha\Zw$ is a nonzero scalar
multiple of $\alpha$ in $W$.~\qed\end{proof}
\step{rpcundetect}{\assume{no such $\Zw$ exists.}
\prove{$\alpha\in W^\perp$.}}
\begin{proof}\pf\ Unfold the definition of the orthogonal complement. ~\qed\end{proof}
\qedstep
\begin{proof}\pf\ Excluded middle for the detecting vector gives
\stepref{rpcdetect} or \stepref{rpcundetect}.~\qed\end{proof}
\end{proof}
\step{rpcpartition}{The roots in $W$ and the roots in $W^\perp$
partition $\Phi$ into orthogonal parts.}
\begin{proof}\pf\ Use \stepref{rpcalternative} and
$W\cap W^\perp=\{0\}$, excluding zero roots. ~\qed\end{proof}
\step{rpconeempty}{One of those parts is empty.}
\begin{proof}\pf\ Apply indecomposability to \stepref{rpcpartition}.~\qed\end{proof}
\step{rpcallW}{\assume{all roots lie in $W$.}
\prove{$W=V$.}}
\begin{proof}\pf\ Use root spanning in Definition~\ref{def:roots}.~\qed\end{proof}
\step{rpcallperp}{\assume{all roots lie in $W^\perp$.}
\prove{$W=0$.}}
\begin{proof}\pf\ Each $\Zw\in W$ is orthogonal to the spanning root
set, thus to itself; use definiteness. ~\qed\end{proof}
\qedstep
\begin{proof}\pf\ Use \stepref{rpconeempty} to select
\stepref{rpcallW} or \stepref{rpcallperp}.~\qed\end{proof}
\end{proof}
\qedstep
\begin{proof}\pf\ Both copies of $V$ have the same orthogonal reflection
maps by \stepref{rpcreflectionmaps}; \stepref{rpcnontrivial},
\stepref{rpcinvariant} supply both irreducibility fields. ~\qed\end{proof}
\end{proof}
\step{rpcconverse}{\assume{$P$ is irreducible and
$\Phi=A\sqcup B$ is a partition into nonempty mutually orthogonal parts.}
\prove{\textsc{False}.}}
\begin{proof}
\pflet{$W=\operatorname{span}A$.}
\step{rpcWnonzero}{$W\ne0$.}
\begin{proof}\pf\ A member of nonempty $A$ is a nonzero vector in $W$.~\qed\end{proof}
\step{rpcWproper}{$W\ne V$.}
\begin{proof}\pf\ A nonzero root in $B$ is orthogonal to $W$ and
cannot be orthogonal to itself. ~\qed\end{proof}
\step{rpcWA}{A reflection at a root of $A$ preserves $W$.}
\begin{proof}\pf\ Its formula subtracts a scalar multiple of an element
of $A$ from a vector of $W$.~\qed\end{proof}
\step{rpcWB}{A reflection at a root of $B$ fixes $W$ pointwise.}
\begin{proof}\pf\ Orthogonality to $A$ extends linearly to $W$ in the
reflection formula. ~\qed\end{proof}
\qedstep
\begin{proof}\pf\ The invariant nonzero proper subspace in
\stepref{rpcWnonzero}--\stepref{rpcWB} contradicts irreducibility. ~\qed\end{proof}
\end{proof}
\qedstep
\begin{proof}\pf\ Combine \stepref{rpcstructure}, \stepref{rpcproperties},
\stepref{rpcirreducible}, \stepref{rpcconverse}.~\qed\end{proof}
\end{proof}
\noindent\new\ We use the direct structure constructor.
\path{RootPairing.ofBilinear} takes all reflective vectors of a form,
and would not select this prescribed finite root set over $\mathbb R$.

\begin{namedresult}{Lemma}{(one common similarity factor)}\label{lem:similarity}
\lean{IsRootSet.similarity_of_cartanMatrix_eq}
\uses{def:rem-bases,lem:rem-edge-constant,facts:inner,facts:rem-library,def:roots}
\assume{\begin{pfenum}
\item $V,W$ are finite-dimensional real inner product spaces and
 $(\alpha_i)_{i\in I}$, $(\beta_i)_{i\in I}$ are real vector-space bases
 indexed by the same finite nonempty set $I$.
\item The matrices $p_\alpha(i,j)=2(\alpha_i,\alpha_j)/\rsq{\alpha_j}$
 and $p_\beta(i,j)=2(\beta_i,\beta_j)/\rsq{\beta_j}$ agree, and their
 common nonzero off-diagonal support is connected.
\end{pfenum}}
\prove{The linear equivalence $f:V\to W$ taking $\alpha_i$ to $\beta_i$
satisfies $(f\Zx,f\Zy)=a(\Zx,\Zy)$ for a single $a>0$.
For any Root Systems $\Phi\subset V$, $\Psi\subset W$ with
$f(\Phi)=\Psi$, this is a similarity of those Root Systems.}
\end{namedresult}
\begin{proof}
\pfsketch\ Ratios of squared simple-root lengths propagate along the
connected support; bilinearity extends the resulting Gram identity.
\pflet{$q(i)=\rsq{\alpha_i}$, $r(i)=\rsq{\beta_i}$, and $M=p_\alpha=p_\beta$.}
\step{rsimpositive}{$q(i)>0$ and $r(i)>0$ for every $i$.}
\begin{proof}\pf\ A vector in a linearly independent family is nonzero;
use the bases in \asm[1]{0} and definiteness. ~\qed\end{proof}
\step{rsimbalance}{$M(i,j)q(j)=M(j,i)q(i)$ and
$M(i,j)r(j)=M(j,i)r(i)$.}
\begin{proof}\pf\ Use symmetry of each inner product in \eqref{eq:pairing}.~\qed\end{proof}
\step{rsimedge}{On a support edge, $r(i)/q(i)=r(j)/q(j)$.}
\begin{proof}\pf\ Divide the two equalities in \stepref{rsimbalance}
by their nonzero entries and the positive lengths of \stepref{rsimpositive}.~\qed\end{proof}
\step{rsimconstant}{There is $a>0$ with $r(i)=a q(i)$ for all $i$.}
\begin{proof}\pf\ Apply
Lemma~\ref{lem:rem-edge-constant} to $i\mapsto r(i)/q(i)$ on the common
Cartan support, using \stepref{rsimedge} and connectedness in \asm[2]{0};
\stepref{rsimpositive} gives a positive value at any vertex. ~\qed\end{proof}
\step{rsimgram}{$(\beta_i,\beta_j)=a(\alpha_i,\alpha_j)$ for every $i,j$.}
\begin{proof}\pf\ Substitute \stepref{rsimconstant} in
$(\beta_i,\beta_j)=M(i,j)r(j)/2$ and \eqref{eq:pairing}.~\qed\end{proof}
\step{rsimmap}{The basis correspondence extends to a linear equivalence $f$.}
\begin{proof}\pf\ Apply \path{Module.Basis.equiv} to the two real
bases and the identity index equivalence. ~\qed\end{proof}
\step{rsimall}{$(f\Zx,f\Zy)=a(\Zx,\Zy)$ for all vectors.}
\begin{proof}\pf\ Expand both vectors in their finite bases and apply
\stepref{rsimgram} term by term. ~\qed\end{proof}
\qedstep
\begin{proof}\pf\ Combine \stepref{rsimconstant}, \stepref{rsimmap},
\stepref{rsimall}; root-set equality supplies exactly the last clause
of the similarity definition. ~\qed\end{proof}
\end{proof}

\begin{namedresult}{Lemma}{(orthogonal components)}\label{lem:rem-components}
\lean{IsRootSet.orthogonal_components}\lean{IsRootSet.component_isRootSet}
\lean{IsRootSet.isIndecomposable_iff_connected}\lean{IsRootSet.subsingleton_of_eq_empty}
\lean{IsRootSet.eq_empty_iff_subsingleton}
\lean{IsRootSet.isEmpty_connectedComponent_of_eq_empty}
\uses{def:roots,lem:rem-reflection,facts:inner}
\assume{$\Phi$ is a Root System in $V$. Join distinct roots when their
inner product is nonzero.}
\prove{Each connected component is an indecomposable Root System in its
own span. Different component spans are orthogonal, and $V$ is their
orthogonal direct sum. In particular, a nonempty $\Phi$ is indecomposable
if and only if this graph is connected. The empty Root System has $V=0$
and no components.}
\end{namedresult}
\begin{proof}
\pfsketch\ Reflections preserve the component of their reflecting root;
orthogonality separates different components and their spans.
\step{roorth}{Roots in distinct components are orthogonal.}
\begin{proof}\pf\ A nonzero product would give an edge between the
components by the graph definition. ~\qed\end{proof}
\step{roreflect}{A component $A$ is closed under reflections in its roots.}
\begin{proof}
\pflet{Fix $\alpha,\beta\in A$.}
\step{rozero}{\assume{$(\alpha,\beta)=0$.}
\prove{$s_\alpha\beta=\beta\in A$.}}
\begin{proof}\pf\ Substitute zero in the reflection formula. ~\qed\end{proof}
\step{rononzero}{\assume{$(\alpha,\beta)\ne0$.}
\prove{$s_\alpha\beta\in A$.}}
\begin{proof}
\step{roproduct}{$(\alpha,s_\alpha\beta)=-(\alpha,\beta)\ne0$.}
\begin{proof}\pf\ Expand the reflection formula by bilinearity. ~\qed\end{proof}
\qedstep
\begin{proof}\pf\ The reflected root equals $\alpha$ or is joined to it
by \stepref{roproduct}, so it lies in the same component. ~\qed\end{proof}
\end{proof}
\qedstep
\begin{proof}\pf\ The cases \stepref{rozero}, \stepref{rononzero} exhaust
the inner-product values. ~\qed\end{proof}
\end{proof}
\step{rorootset}{Each component is a Root System in its span.}
\begin{proof}\pf\ Finiteness, nonzero roots, integrality and reducedness
are inherited from $\Phi$; \stepref{roreflect} gives closure and its
own span gives the spanning field of Definition~\ref{def:roots}.~\qed\end{proof}
\step{roindecomp}{Each component is indecomposable.}
\begin{proof}\pf\ An orthogonal partition could have no graph edge
crossing its parts by definition, contrary to connectedness of the component. ~\qed\end{proof}
\step{rospans}{Distinct component spans are orthogonal.}
\begin{proof}\pf\ Expand finite linear combinations using
\stepref{roorth} and bilinearity. ~\qed\end{proof}
\step{rodirect}{The sum of those spans is direct and equals $V$.}
\begin{proof}
\step{roallspan}{The sum equals $V$.}
\begin{proof}\pf\ Every root lies in a component, and $\Phi$ spans $V$.~\qed\end{proof}
\step{roindependent}{\assume{$\sum_A \Zx_A=0$ with each $\Zx_A$ in its
component span.}
\prove{Every $\Zx_A=0$.}}
\begin{proof}\pf\ Pair the relation with $\Zx_A$;
\stepref{rospans} leaves $\rsq{\Zx_A}=0$, so definiteness applies. ~\qed\end{proof}
\qedstep
\begin{proof}\pf\ Combine \stepref{roallspan}, \stepref{roindependent}.~\qed\end{proof}
\end{proof}
\step{roconnected}{For nonempty $\Phi$, indecomposability is equivalent
to connectedness of its graph.}
\begin{proof}\pf\ Multiple components give an orthogonal partition by
\stepref{roorth}; one component is indecomposable by \stepref{roindecomp}.~\qed\end{proof}
\step{roempty}{\assume{$\Phi$ is empty.}
\prove{$V=0$.}}
\begin{proof}\pf\ Use its spanning condition and the span of the empty set. ~\qed\end{proof}
\qedstep
\begin{proof}\pf\ Combine \stepref{rorootset}, \stepref{roindecomp},
\stepref{rospans}, \stepref{rodirect}, \stepref{roconnected}, \stepref{roempty}.~\qed\end{proof}
\end{proof}

\begin{namedresult}{Lemma}{(indecomposable classification)}\label{lem:rem-indecomp-classification}
\lean{IsRootSet.classification_of_isIndecomposable}\lean{RootSetModel.two_lengths_iff}
\lean{RootSetModel.squaredLengths_eq}\lean{RootSetModel.B₂_C₂_similarity}
\uses{def:roots,def:rem-bases,lem:rem-similarity-laws,lem:pairing-construction,thm:A,lem:classical,lem:E,lem:F,lem:G,lem:connected,lem:similarity,lem:counts,facts:rem-library}
\assume{$\Phi$ is a nonempty indecomposable Root System in a
finite-dimensional real inner product space $V$.}
\prove{After one common positive rescaling, $\Phi$ is isometric to exactly
one coordinate Root System in \eqref{eq:list}. Exactly $\typ B,\typ C,
\typ F,\typ G$ have two lengths. The additional $\typ C_2$ coordinate
model is similar to $\typ B_2$.}
\end{namedresult}
\begin{proof}
\pfsketch\ Apply Theorem A to a Mathlib base. Construct the basis map
and use the library equivalence of root pairings. Counts distinguish
the models without a conjugacy theorem for bases.
\step{rBpairing}{The pairing $P$ constructed from $\Phi$ is finite,
spanning, reduced, crystallographic and irreducible over $\mathbb R$.}
\begin{proof}\pf\ Apply Lemma~\ref{lem:pairing-construction} to \asm{0}.~\qed\end{proof}
\step{rBbase}{\pick\ a Mathlib base $b$ for $P$.}
\begin{proof}\pf\ Apply \path{RootPairing.nonempty_base} to
\stepref{rBpairing}; $\mathbb R$ is a characteristic zero field. ~\qed\end{proof}
\step{rBsupport}{The finite base support is nonempty.}
\begin{proof}\pf\ Apply \path{RootPairing.Base.support_nonempty} to
\stepref{rBbase}, the nonempty root set and $2\ne0$ in $\mathbb R$.~\qed\end{proof}
\step{rBfinite}{Its integer Cartan matrix is finite Cartan.}
\begin{proof}\pf\ Apply \path{RootPairing.Base.cartanMatrix_isFiniteCartan}
using \stepref{rBpairing} and classically chosen decidable equality. ~\qed\end{proof}
\step{rBindecomp}{That matrix is indecomposable.}
\begin{proof}\pf\ Apply \path{RootPairing.Base.cartanMatrix_isIndecomposable}
using \stepref{rBpairing}.~\qed\end{proof}
\step{rBtype}{\pick\ a type, rank $\ell$, and reindexing from its
model indices to the support of $b$.}
\begin{proof}\pf\ Apply Theorem~\ref{thm:A} to
\stepref{rBsupport}, \stepref{rBfinite}, \stepref{rBindecomp}.~\qed\end{proof}
\step{rBmodel}{There is a coordinate Root System $\Psi$ with a root base
having precisely this reindexed matrix.}
\begin{proof}\pf\ Select the matching case of Lemmas~\ref{lem:classical},
\ref{lem:E}, \ref{lem:F}, \ref{lem:G}, using the rank restrictions
in \stepref{rBtype}.~\qed\end{proof}
\step{rBmodelpairing}{Its pairing $P_2$ is finite, spanning, reduced and
crystallographic over $\mathbb R$.}
\begin{proof}\pf\ Apply Lemma~\ref{lem:pairing-construction} to
the Root System in \stepref{rBmodel}.~\qed\end{proof}
\step{rBmodelbase}{The displayed root base gives a Mathlib base $b_2$.}
\begin{proof}
\step{rBmonoid}{Every model root or its negative lies in the additive
monoid of the selected simple roots.}
\begin{proof}\pf\ Use the nonnegative integer expansions in
\stepref{rBmodel} and \path{Submodule.span_nat_eq_addSubmonoidClosure}.~\qed\end{proof}
\qedstep
\begin{proof}\pf\ Apply \path{RootPairing.Base.mk'} using
\stepref{rBmodelpairing}, root independence in \stepref{rBmodel},
and \stepref{rBmonoid}.~\qed\end{proof}
\end{proof}
\pflet{Let $e:b.\mathrm{support}\to b_2.\mathrm{support}$ be the index
bijection induced by \stepref{rBtype} and the model base parametrization.}
\step{rBcartan}{For all base indices $i,j$,
$b_2.\mathrm{cartanMatrix}(e(i),e(j))=b.\mathrm{cartanMatrix}(i,j)$.}
\begin{proof}
\step{rBrealcartan}{The real pairings of the two bases agree through $e$.}
\begin{proof}\pf\ The entry equality in \stepref{rBtype} and the
model pairing matrix in \stepref{rBmodel} identify both real
pairings with the same cast named matrix, using
\path{RootPairing.Base.algebraMap_cartanMatrixIn_apply}.~\qed\end{proof}
\qedstep
\begin{proof}\pf\ Rewrite both real pairings in
\stepref{rBrealcartan} as casts of integer Cartan entries by
\path{RootPairing.Base.algebraMap_cartanMatrixIn_apply};
injectivity of $\mathbb Z\to\mathbb R$ gives the equality.~\qed\end{proof}
\end{proof}
\step{rBmap}{\pick\ the linear equivalence
$f=b.\mathrm{toWeightBasis.equiv}\ b_2.\mathrm{toWeightBasis}\ e$.}
\begin{proof}\pf\ Use the spanning pairings
\stepref{rBpairing}, \stepref{rBmodelpairing} in the weight-basis
and basis-equivalence declarations of \ref{facts:rem-library}.~\qed\end{proof}
\step{rBsimple}{$f(\alpha_i)=\beta_{e(i)}$ for every simple root.}
\begin{proof}\pf\ Apply \path{Module.Basis.equiv_apply} and
\path{RootPairing.Base.toWeightBasis_apply} to \stepref{rBmap}.~\qed\end{proof}
\step{rBequivalence}{There is a root-pairing equivalence
$P\simeq P_2$ whose map on root spaces is $f$.}
\begin{proof}\pf\ Apply \path{RootPairing.Base.equivOfCartanMatrixEq}
to $b,b_2,e$ and \stepref{rBcartan}, using the source hypotheses in
\stepref{rBpairing} and finite, spanning, reduced, crystallographic
target hypotheses in \stepref{rBmodelpairing}.
Its defining root-space map is exactly \stepref{rBmap}.~\qed\end{proof}
\step{rBonto}{$f(\Phi)=\Psi$.}
\begin{proof}\pf\ The index equivalence and root compatibility of
\stepref{rBequivalence} give a bijection of the root sets.~\qed\end{proof}
\step{rBconnected}{The common base matrix has connected nonzero support.}
\begin{proof}\pf\ Apply Lemma~\ref{lem:connected} to
\stepref{rBsupport}, \stepref{rBfinite}, \stepref{rBindecomp}.~\qed\end{proof}
\step{rBsimilarity}{$f$ is a similarity using one positive factor.}
\begin{proof}
\step{rBreindexedbases}{The families
$\alpha_i=b.\mathrm{toWeightBasis}(i)$ and
$\widehat\beta_i=b_2.\mathrm{toWeightBasis}(e(i))$, indexed by
$b.\mathrm{support}$, are real vector-space bases.}
\begin{proof}\pf\ Reindex the second weight basis along $e$;
the first is the weight basis from \stepref{rBmap}.~\qed\end{proof}
\step{rBmetric}{For some $a>0$, $(f\Zx,f\Zy)=a(\Zx,\Zy)$ for all vectors.}
\begin{proof}\pf\ Apply Lemma~\ref{lem:similarity} to
\stepref{rBreindexedbases}, \stepref{rBcartan} cast to real pairings,
and \stepref{rBconnected}; \stepref{rBsimple} identifies its basis
map with $f$.~\qed\end{proof}
\qedstep
\begin{proof}\pf\ Steps \stepref{rBmetric}, \stepref{rBonto}
are the similarity conditions of Definition~\ref{def:roots}.~\qed\end{proof}
\end{proof}
\step{rBunique}{No two distinct models in \eqref{eq:list} are similar.}
\begin{proof}
\step{rBinvariants}{Similarity preserves rank, root count, number of
lengths, their ratio, and the numbers of short and long roots.}
\begin{proof}\pf\ A bijection scaling every squared length by the same
$a>0$ preserves these definitions. ~\qed\end{proof}
\step{rBonevsTwo}{One-length and two-length models are distinguished.}
\begin{proof}\pf\ Use the lengths in Lemma~\ref{lem:counts} with
\stepref{rBinvariants}.~\qed\end{proof}
\step{rBAD}{At equal rank, the $\typ A$ and $\typ D$ root counts differ
throughout the permitted $\typ D$ range.}
\begin{proof}\pf\ The equality $\ell(\ell+1)=2\ell(\ell-1)$ with
$\ell>0$ implies $\ell=3$, excluded by $\ell\geq4$.~\qed\end{proof}
\step{rBE6}{At rank 6, the counts $42,60,72$ distinguish $\typ A,\typ D,\typ E$.}
\begin{proof}\pf\ Evaluate Lemma~\ref{lem:counts} at rank 6. ~\qed\end{proof}
\step{rBE7}{At rank 7, the counts $56,84,126$ distinguish these types.}
\begin{proof}\pf\ Evaluate Lemma~\ref{lem:counts} at rank 7. ~\qed\end{proof}
\step{rBE8}{At rank 8, the counts $72,112,240$ distinguish these types.}
\begin{proof}\pf\ Evaluate Lemma~\ref{lem:counts} at rank 8. ~\qed\end{proof}
\step{rBGunique}{Squared-length ratio 3 distinguishes $\typ G_2$.}
\begin{proof}\pf\ All other two-length models have ratio 2 by
Lemma~\ref{lem:counts}.~\qed\end{proof}
\step{rBFunique}{At rank 4, $\typ F_4$ has 48 roots whereas
$\typ B_4,\typ C_4$ have 32.}
\begin{proof}\pf\ Evaluate Lemma~\ref{lem:counts}.~\qed\end{proof}
\step{rBBCunique}{For $\ell\geq3$, $\typ B_\ell,\typ C_\ell$ have
different numbers of short roots.}
\begin{proof}\pf\ Their counts $2\ell$ and $2\ell(\ell-1)$ from
Lemma~\ref{lem:counts} differ when $\ell\geq3$.~\qed\end{proof}
\qedstep
\begin{proof}\pf\ The cases \stepref{rBonevsTwo}--\stepref{rBBCunique}
separate every pair at equal rank, and \stepref{rBinvariants} separates
unequal ranks. ~\qed\end{proof}
\end{proof}
\step{rBCtwo}{The map with matrix
$\left(\begin{smallmatrix}1&1\\-1&1\end{smallmatrix}\right)$
is a similarity from $\typ B_2$ onto $\typ C_2$.}
\begin{proof}
\step{rBCaxis1}{$\Zu_1$ maps to $\Zu_1-\Zu_2$.}
\begin{proof}\pf\ Read column 1 of the displayed matrix. ~\qed\end{proof}
\step{rBCaxis2}{$\Zu_2$ maps to $\Zu_1+\Zu_2$.}
\begin{proof}\pf\ Read column 2 of the displayed matrix. ~\qed\end{proof}
\step{rBCplus}{$\Zu_1+\Zu_2$ maps to $2\Zu_1$.}
\begin{proof}\pf\ Add \stepref{rBCaxis1}, \stepref{rBCaxis2}.~\qed\end{proof}
\step{rBCminus}{$\Zu_1-\Zu_2$ maps to $-2\Zu_2$.}
\begin{proof}\pf\ Subtract \stepref{rBCaxis2} from \stepref{rBCaxis1}.~\qed\end{proof}
\step{rBCmetric}{The transpose times the matrix is $2I$.}
\begin{proof}\pf\ Compute each entry of $T^TT$ as the dot product of two columns of the displayed matrix $T$. ~\qed\end{proof}
\qedstep
\begin{proof}\pf\ The four images and their negatives exhaust
$\Phi(\typ C_2)$, while \stepref{rBCmetric} gives scale 2 and invertibility. ~\qed\end{proof}
\end{proof}
\qedstep
\begin{proof}\pf\ Use \stepref{rBsimilarity}, \stepref{rBunique},
\stepref{rBCtwo} and the length assertion of
Lemma~\ref{lem:counts}. Two similarities from $\Phi$ to models
compose, after inversion of the first, to a similarity between the
models by Lemma~\ref{lem:rem-similarity-laws}; \stepref{rBunique}
then forces the same model.~\qed\end{proof}
\end{proof}

\begin{namedresult}{Theorem}{B}\label{thm:B}
\lean{IsRootSet.classification}
\uses{def:roots,lem:rem-indecomp-classification,lem:rem-components,lem:models-irreducible,def:rem-bases}
\assume{$\Phi$ is a Root System in a finite-dimensional real inner
product space $V$.}
\prove{If $\Phi$ is nonempty and indecomposable, it is similar, by one
common positive scale, to exactly one model in \eqref{eq:list}.
Exactly the $\typ B,\typ C,\typ F,\typ G$ models have two lengths;
$\typ C_2$ is similar to $\typ B_2$ and is omitted from that list.
In general $V$ is the orthogonal direct sum of the component spans, and
each nonzero component is similar to a listed indecomposable model,
with an independent scale on each component. $\Phi$ is empty exactly
when $V$ is the zero space, and then there are no components.}
\end{namedresult}
\begin{proof}
\pfsketch\ \new\ The indecomposable assertion is a separate lemma so
that applying it to each component does not reuse local proof steps.
\step{rBmain}{The nonempty indecomposable assertion and the model length
and rank-2 assertions hold.}
\begin{proof}\pf\ Apply Lemma~\ref{lem:rem-indecomp-classification}.~\qed\end{proof}
\step{rBdecompose}{Every nonempty component is an indecomposable Root
System in its span, and these spans form an orthogonal direct sum.}
\begin{proof}\pf\ Apply Lemma~\ref{lem:rem-components} to $\Phi$.~\qed\end{proof}
\step{rBclassifycomponents}{Each component is similar to a unique listed
model using a single positive scale for that component.}
\begin{proof}\pf\ Apply Lemma~\ref{lem:rem-indecomp-classification}
separately to each Root System supplied by \stepref{rBdecompose}.~\qed\end{proof}
\step{rBmodelindecomp}{All listed models are indecomposable.}
\begin{proof}\pf\ Apply Lemma~\ref{lem:models-irreducible}.~\qed\end{proof}
\qedstep
\begin{proof}\pf\ Combine \stepref{rBmain}--\stepref{rBmodelindecomp};
the empty case is the final clause of Lemma~\ref{lem:rem-components}, and
conversely the zero space has no nonzero vector, so no root.~\qed\end{proof}
\end{proof}

\subsubsection{Angles, lines and the normalized matrix}

\paragraph{Angle convention.}\label{facts:rem-angles}
\begingroup\raggedright
For two distinct lines the angle lies in $(0,90^\circ]$ and its cosine
is $|(\alpha,\beta)|/(\|\alpha\|\|\beta\|)$ for nonzero representatives.
We use \path{Real.cos_pi_div_three}, \path{Real.cos_pi_div_four},
\path{Real.cos_pi_div_six}, and \path{Real.injOn_cos}, all in
\path{Mathlib/Analysis/SpecialFunctions/Trigonometric/Basic.lean}.
Their values and injectivity on $[0,\pi]$ identify squared cosines
$1/4,1/2,3/4$ with $60^\circ,45^\circ,30^\circ$, respectively.
\par\endgroup

\begin{namedresult}{Lemma}{(angles and fixed-length subsystems)}\label{lem:angles}
\lean{IsRootSet.pairing_product_mem}\lean{IsRootSet.angle_length_cases}
\lean{IsRootSet.fixedLength_isRootSet}
\lean{IsRootSet.fixedLength_components_simplyLaced}
\uses{def:roots,lem:rem-reflection,thm:B,lem:counts,facts:inner,facts:rem-angles}
\assume{$\Phi$ is a Root System in a finite-dimensional real inner product space.}
\prove{For every $\alpha,\beta\in\Phi$ which are neither perpendicular
nor proportional,
\[
\frac{4(\alpha,\beta)^2}{\rsq\alpha\,\rsq\beta}\in\{1,2,3\}.
\]
The longer-to-shorter length ratio and acute line angle are exactly one of
$(1,60^\circ)$, $(\sqrt2,45^\circ)$, $(\sqrt3,30^\circ)$.
For every occurring length $a>0$, the roots of length $a$ form a Root
System in their own span, and it is an orthogonal sum of systems similar
to $\typ A_\ell$ ($\ell\geq1$), $\typ D_\ell$ ($\ell\geq4$),
$\typ E_6,\typ E_7,\typ E_8$.}
\end{namedresult}
\noindent\repair\ The hypotheses exclude proportional pairs; fixed-length
subsystems need not be indecomposable or span the original space.
\begin{proof}
\pfsketch\ Bound the product of two nonzero integer pairings by 4.
Reflections preserve length, then Theorem B classifies each component.
\step{rangpair}{\assume{$\alpha,\beta\in\Phi$ are neither perpendicular
nor proportional.}
\prove{Their pairing product, length ratio and angle have the stated alternatives.}}
\begin{proof}
\pflet{$q=\rsq\alpha>0$, $r=\rsq\beta>0$, $h=(\alpha,\beta)\ne0$,
$u=2h/q$, and $v=2h/r$.}
\step{ranginteger}{$u,v$ are nonzero integers of the same sign.}
\begin{proof}\pf\ Use the two integrality conditions in
Definition~\ref{def:roots} and the positive denominators $q,r$.~\qed\end{proof}
\step{rangstrict}{$h^2<qr$.}
\begin{proof}
\step{rangresidual}{$\beta-(h/q)\alpha\ne0$.}
\begin{proof}\pf\ Otherwise $\beta$ would be proportional to $\alpha$,
contrary to \asm{1}.~\qed\end{proof}
\step{rangnorm}{$\|\beta-(h/q)\alpha\|^2=r-h^2/q>0$.}
\begin{proof}
\step{rangnormidentity}{$\|\beta-(h/q)\alpha\|^2=r-h^2/q$.}
\begin{proof}\pf\ Expand by bilinearity in \ref{facts:inner} and
substitute the definitions of $q,r,h$.~\qed\end{proof}
\step{rangnormpositive}{$\|\beta-(h/q)\alpha\|^2>0$.}
\begin{proof}\pf\ Apply definiteness in \ref{facts:inner} to
\stepref{rangresidual}.~\qed\end{proof}
\qedstep
\begin{proof}\pf\ Combine \stepref{rangnormidentity},
\stepref{rangnormpositive}.~\qed\end{proof}
\end{proof}
\qedstep
\begin{proof}\pf\ Multiply \stepref{rangnorm} by $q>0$.~\qed\end{proof}
\end{proof}
\step{rangproducts}{$uv=4h^2/(qr)\in\{1,2,3\}$.}
\begin{proof}\pf\ It is a positive integer less than 4 by
\stepref{ranginteger}, \stepref{rangstrict}.~\qed\end{proof}
\step{rangcases}{The length and angle alternatives follow.}
\begin{proof}
\step{rangone}{\assume{$uv=1$.}
\prove{$|u|=|v|=1$, $q=r$, and the angle is $60^\circ$.}}
\begin{proof}
\step{rangoneinteger}{$|u|=|v|=1$.}
\begin{proof}\pf\ A product of two positive integers is 1 only when
both are 1, applied to \stepref{ranginteger} and $uv=1$.~\qed\end{proof}
\step{rangonelength}{$q=r$.}
\begin{proof}\pf\ Substitute \stepref{rangoneinteger} in
$u/v=r/q$, using the equal signs in \stepref{ranginteger}.~\qed\end{proof}
\step{rangoneangle}{The angle is $60^\circ$.}
\begin{proof}\pf\ The squared cosine is $h^2/(qr)=uv/4=1/4$;
apply \ref{facts:rem-angles}.~\qed\end{proof}
\qedstep
\begin{proof}\pf\ Combine \stepref{rangoneinteger},
\stepref{rangonelength}, \stepref{rangoneangle}.~\qed\end{proof}
\end{proof}
\step{rangtwo}{\assume{$uv=2$.}
\prove{$\{|u|,|v|\}=\{1,2\}$, the squared-length
ratio is 2, and the angle is $45^\circ$.}}
\begin{proof}
\step{rangtwointeger}{$(|u|,|v|)$ is $(1,2)$ or $(2,1)$.}
\begin{proof}\pf\ The only positive divisors of 2 are 1 and 2,
applied to \stepref{ranginteger} and $uv=2$.~\qed\end{proof}
\step{rangtwolength}{$\max(q,r)/\min(q,r)=2$.}
\begin{proof}\pf\ Substitute the two alternatives
\stepref{rangtwointeger} in $u/v=r/q$, with $q,r>0$.~\qed\end{proof}
\step{rangtwoangle}{The angle is $45^\circ$.}
\begin{proof}\pf\ The squared cosine is $h^2/(qr)=uv/4=1/2$;
apply \ref{facts:rem-angles}.~\qed\end{proof}
\qedstep
\begin{proof}\pf\ Combine \stepref{rangtwointeger},
\stepref{rangtwolength}, \stepref{rangtwoangle}.~\qed\end{proof}
\end{proof}
\step{rangthree}{\assume{$uv=3$.}
\prove{$\{|u|,|v|\}=\{1,3\}$, the squared-length
ratio is 3, and the angle is $30^\circ$.}}
\begin{proof}
\step{rangthreeinteger}{$(|u|,|v|)$ is $(1,3)$ or $(3,1)$.}
\begin{proof}\pf\ The only positive divisors of 3 are 1 and 3,
applied to \stepref{ranginteger} and $uv=3$.~\qed\end{proof}
\step{rangthreelength}{$\max(q,r)/\min(q,r)=3$.}
\begin{proof}\pf\ Substitute the two alternatives
\stepref{rangthreeinteger} in $u/v=r/q$, with $q,r>0$.~\qed\end{proof}
\step{rangthreeangle}{The angle is $30^\circ$.}
\begin{proof}\pf\ The squared cosine is $h^2/(qr)=uv/4=3/4$;
apply \ref{facts:rem-angles}.~\qed\end{proof}
\qedstep
\begin{proof}\pf\ Combine \stepref{rangthreeinteger},
\stepref{rangthreelength}, \stepref{rangthreeangle}.~\qed\end{proof}
\end{proof}
\qedstep
\begin{proof}\pf\ Take positive square roots of the squared-length
ratios in the exhaustive cases \stepref{rangone}--\stepref{rangthree}.~\qed\end{proof}
\end{proof}
\qedstep
\begin{proof}\pf\ Combine \stepref{rangproducts}, \stepref{rangcases}.~\qed\end{proof}
\end{proof}
\step{rangfixed}{For each occurring $a>0$, the set
$\Psi=\{\gamma\in\Phi:\|\gamma\|=a\}$ is a Root System in its span.}
\begin{proof}
\step{rangfixedbasic}{$\Psi$ is finite, nonzero, reduced, integral,
and spans its own span.}
\begin{proof}\pf\ Restrict the corresponding conditions of
Definition~\ref{def:roots} to $\Psi$.~\qed\end{proof}
\step{rangfixedroot}{For $\gamma,\delta\in\Psi$, $s_\gamma\delta\in\Phi$.}
\begin{proof}\pf\ Apply reflection closure of $\Phi$.~\qed\end{proof}
\step{rangfixednorm}{$\|s_\gamma\delta\|=a$.}
\begin{proof}\pf\ Apply Lemma~\ref{lem:rem-reflection} and $\|\delta\|=a$.~\qed\end{proof}
\step{rangfixedspan}{$s_\gamma\delta$ lies in $\operatorname{span}\Psi$.}
\begin{proof}\pf\ Its defining formula is a linear combination of
$\gamma,\delta\in\Psi$.~\qed\end{proof}
\qedstep
\begin{proof}\pf\ Assemble Definition~\ref{def:roots} from
\stepref{rangfixedbasic}--\stepref{rangfixedspan}.~\qed\end{proof}
\end{proof}
\step{rangADE}{Every nonempty component of this subsystem has type
$\typ A,\typ D$ or $\typ E$.}
\begin{proof}\pf\ Apply Theorem~\ref{thm:B} to \stepref{rangfixed};
its other models have two lengths by Lemma~\ref{lem:counts}, impossible
in the fixed-length subset. ~\qed\end{proof}
\qedstep
\begin{proof}\pf\ Combine \stepref{rangpair},
\stepref{rangfixed}, \stepref{rangADE}.~\qed\end{proof}
\end{proof}

\begin{namedresult}{Lemma}{(unit-line orthogonal projections)}\label{lem:rem-projectors}
\lean{IsRootSet.lineProjector_identities}
\uses{facts:inner}
\assume{$V$ is a finite-dimensional real inner product space,
$\Zv$ is a unit vector, and $P_\Zv\Zx=(\Zv,\Zx)\Zv$ is the rank-one orthogonal projection onto $\mathbb R\Zv$.}
\prove{$P_{-\Zv}=P_\Zv$; $P_\Zv$ is symmetric; its trace is 1;
and for every orthogonal linear bijection $f$,
$fP_\Zv f^{-1}=P_{f\Zv}$.}
\end{namedresult}
\begin{proof}
\pfsketch\ Expand the rank-one operator and its diagonal.
\step{rprsign}{$P_{-\Zv}\Zx=P_\Zv\Zx$.}
\begin{proof}\pf\ The two minus signs cancel by bilinearity. ~\qed\end{proof}
\step{rprsymmetry}{$(P_\Zv\Zx,\Zy)=(\Zx,P_\Zv\Zy)$.}
\begin{proof}\pf\ Both sides are $(\Zv,\Zx)(\Zv,\Zy)$ by
\ref{facts:inner}.~\qed\end{proof}
\step{rprtrace}{In an orthonormal basis the trace is $\sum_i v(i)^2=1$.}
\begin{proof}\pf\ The $i$th diagonal entry is $(\Zv,\Zu_i)^2$;
sum and use $\|\Zv\|=1$.~\qed\end{proof}
\step{rprconjugate}{$fP_\Zv f^{-1}\Zx=(f\Zv,\Zx)f\Zv$.}
\begin{proof}\pf\ Substitute the orthogonal projection formula and the isometry
identity $(\Zv,f^{-1}\Zx)=(f\Zv,\Zx)$.~\qed\end{proof}
\qedstep
\begin{proof}\pf\ Operator extensionality applied to
\stepref{rprsign}, \stepref{rprsymmetry}, \stepref{rprtrace},
\stepref{rprconjugate} gives all four identities. ~\qed\end{proof}
\end{proof}

\begin{namedresult}{Lemma}{(the line frame)}\label{lem:tight-frame}
\lean{IsRootSet.frame_identity}
\uses{def:roots,lem:rem-reflection,lem:rem-components,lem:rem-edge-constant,lem:rem-projectors,facts:inner}
\assume{\begin{pfenum}
\item $\Phi$ is a nonempty indecomposable Root System in an
 $n$-dimensional real inner product space $V$.
\item Its line set is $\mathcal L$, $|\mathcal L|=N$, and $\Zv_L$ is
 any unit representative of each line.
\end{pfenum}}
\prove{$n>0$ and
\begin{equation}\label{eq:frame}
\mathcal F:=\sum_{L\in\mathcal L}P_{\Zv_L}=\frac Nn I_V.
\end{equation}
The identity is independent of all representative signs.}
\end{namedresult}
\begin{proof}
\pfsketch\ The symmetric frame operator commutes with every root
reflection. It acts by one scalar on all roots, and its trace determines it.
\step{rfrdimension}{$n>0$.}
\begin{proof}\pf\ The space contains a nonzero root from \asm[1]{0}.~\qed\end{proof}
\step{rfrsign}{The operator $\mathcal F$ is independent of signs and symmetric.}
\begin{proof}\pf\ Sum the sign and symmetry identities of
Lemma~\ref{lem:rem-projectors}.~\qed\end{proof}
\step{rfrcommute}{$\mathcal F$ commutes with each root reflection $s_\alpha$.}
\begin{proof}
\step{rfrpermute}{$s_\alpha$ permutes $\mathcal L$.}
\begin{proof}\pf\ It is an orthogonal root permutation by
Lemma~\ref{lem:rem-reflection}.~\qed\end{proof}
\step{rfrconjugation}{$s_\alpha\mathcal F s_\alpha^{-1}=\mathcal F$.}
\begin{proof}\pf\ Reindex the sum by \stepref{rfrpermute}, using
projector conjugation and sign independence in Lemma~\ref{lem:rem-projectors}.~\qed\end{proof}
\qedstep
\begin{proof}\pf\ Multiply \stepref{rfrconjugation} on the right by
$s_\alpha$.~\qed\end{proof}
\end{proof}
\step{rfreigen}{For every root $\alpha$ there is a unique scalar
$t_\alpha$ with $\mathcal F\alpha=t_\alpha\alpha$.}
\begin{proof}
\step{rfrminus}{$s_\alpha(\mathcal F\alpha)=-\mathcal F\alpha$.}
\begin{proof}\pf\ Use \stepref{rfrcommute} and
$s_\alpha\alpha=-\alpha$ from Lemma~\ref{lem:rem-reflection}.~\qed\end{proof}
\qedstep
\begin{proof}\pf\ The $-1$ eigenspace in
Lemma~\ref{lem:rem-reflection} is $\mathbb R\alpha$; uniqueness uses
$\alpha\ne0$.~\qed\end{proof}
\end{proof}
\step{rfredge}{Nonorthogonal roots have the same scalar.}
\begin{proof}\pf\ Symmetry in \stepref{rfrsign} gives
$(t_\alpha-t_\beta)(\alpha,\beta)=0$ by \stepref{rfreigen};
cancel the nonzero inner product. ~\qed\end{proof}
\step{rfrconstant}{There is a single real $t$ with $t_\alpha=t$
for all roots.}
\begin{proof}\pf\ Apply Lemma~\ref{lem:rem-edge-constant} to
$\alpha\mapsto t_\alpha$ on the nonorthogonality graph of all roots,
connected by Lemma~\ref{lem:rem-components}, using \stepref{rfredge}.~\qed\end{proof}
\step{rfrscalar}{$\mathcal F=tI_V$.}
\begin{proof}\pf\ Expand a vector in the spanning roots and use
\stepref{rfreigen}, \stepref{rfrconstant}.~\qed\end{proof}
\step{rfrtrace}{$\operatorname{tr}\mathcal F=N$.}
\begin{proof}\pf\ Sum the diagonal entries and the unit traces of
Lemma~\ref{lem:rem-projectors}.~\qed\end{proof}
\step{rfrvalue}{$t=N/n$.}
\begin{proof}\pf\ The trace of \stepref{rfrscalar} is $nt$;
use \stepref{rfrtrace} and divide by $n>0$ from \stepref{rfrdimension}.~\qed\end{proof}
\qedstep
\begin{proof}\pf\ Substitute \stepref{rfrvalue} in \stepref{rfrscalar}
and use \stepref{rfrdimension}, \stepref{rfrsign}.~\qed\end{proof}
\end{proof}

\begin{namedresult}{Lemma}{(the normalized matrix convention)}\label{lem:C-convention}
\lean{IsRootSet.signedAdj_quadratic_iff_frame}\lean{IsRootSet.lineFrame_signs}
\lean{IsRootSet.signedAdj_signs}
\uses{lem:rem-projectors,facts:inner}
\assume{\begin{pfenum}
\item $N$ lines span an $n$-dimensional real inner product space, $n>0$.
\item In an orthonormal coordinate basis, $U$ has a chosen unit vector on
 each line as its columns, $G=U^TU$, and $\theta\in\mathbb R$.
\end{pfenum}}
\pflet{\begin{equation}\label{eq:C-convention}C=2G-2I_N.\end{equation}}
\prove{$(C+2I_N)(C-\theta I_N)=0$ if and only if
$UU^T=(N/n)I_n$ and $\theta=2N/n-2$.
Both conditions are independent of the choices of signs.}
\end{namedresult}
\noindent\repair\ This extends the squared-norm-2 representative
convention preceding Theorem 3.7. Unequal-length roots are normalized
line by line for this matrix; that operation is not asserted to preserve
the Root System. Entries may be $\pm\sqrt2,\pm\sqrt3$.
\begin{proof}
\pfsketch\ Reduce the polynomial to $G^2=tG$, then use surjectivity of
$U$ to cancel on both sides and use the trace.
\pflet{$t=(\theta+2)/2$.}
\step{rCexpand}{$(C+2I)(C-\theta I)=4(G^2-tG)$.}
\begin{proof}\pf\ Expand using \eqref{eq:C-convention} and the definition
of $t$.~\qed\end{proof}
\step{rCfactor}{$G^2-tG=U^T(UU^T-tI)U$.}
\begin{proof}\pf\ Substitute $G=U^TU$ and use matrix associativity. ~\qed\end{proof}
\step{rCsurjective}{$U$ is surjective onto the ambient coordinate space.}
\begin{proof}\pf\ Its columns span by \asm[1]{0}.~\qed\end{proof}
\step{rCforward}{\assume{$G^2-tG=0$.}
\prove{$UU^T=tI$.}}
\begin{proof}
\step{rCtest}{For any $\Zx,\Zy$ there exist $\Zv,\Zw$ with
$U\Zv=\Zx$, $U\Zw=\Zy$.}
\begin{proof}\pf\ Apply \stepref{rCsurjective} to each of $\Zx,\Zy$.~\qed\end{proof}
\step{rCbilinear}{$\Zx^T(UU^T-tI)\Zy=0$ for every $\Zx,\Zy$.}
\begin{proof}\pf\ Substitute \stepref{rCtest} into \stepref{rCfactor}
under $G^2-tG=0$.~\qed\end{proof}
\qedstep
\begin{proof}\pf\ Evaluate \stepref{rCbilinear} on every pair of
coordinate unit vectors to obtain each zero entry. ~\qed\end{proof}
\end{proof}
\step{rCreverse}{\assume{$UU^T=tI$.}
\prove{$G^2-tG=0$.}}
\begin{proof}\pf\ Substitute in \stepref{rCfactor}.~\qed\end{proof}
\step{rCtrace}{\assume{$UU^T=tI$.}
\prove{$t=N/n$.}}
\begin{proof}
\step{rCunittrace}{$\operatorname{tr}(UU^T)=N$.}
\begin{proof}\pf\ Sum the squared coordinates of each unit column in
\asm[2]{0}, using Lemma~\ref{lem:rem-projectors}.~\qed\end{proof}
\qedstep
\begin{proof}\pf\ The assumed equality gives $nt=N$ by
\stepref{rCunittrace}; cancel $n>0$.~\qed\end{proof}
\end{proof}
\step{rCsigns}{Changing column signs leaves $UU^T$ unchanged and
conjugates both $G$ and $C$ by a diagonal sign matrix.}
\begin{proof}
\pflet{$D_\epsilon=\operatorname{diag}(\epsilon)$ with $\epsilon(i)^2=1$.}
\step{rCsignsquare}{$D_\epsilon^T=D_\epsilon$ and $D_\epsilon^2=I$.}
\begin{proof}\pf\ Multiply diagonal entries. ~\qed\end{proof}
\step{rCsignframe}{$(UD_\epsilon)(UD_\epsilon)^T=UU^T$.}
\begin{proof}\pf\ Apply associativity and \stepref{rCsignsquare}.~\qed\end{proof}
\step{rCsignGram}{$(UD_\epsilon)^T(UD_\epsilon)=D_\epsilon G D_\epsilon$.}
\begin{proof}\pf\ Use transpose of a product and $G=U^TU$.~\qed\end{proof}
\qedstep
\begin{proof}\pf\ Substitute \stepref{rCsignGram} in
\eqref{eq:C-convention}, and use \stepref{rCsignsquare},
\stepref{rCsignframe}.~\qed\end{proof}
\end{proof}
\qedstep
\begin{proof}\pf\ Combine \stepref{rCexpand}, \stepref{rCforward},
\stepref{rCreverse}, \stepref{rCtrace}, $t=(\theta+2)/2$, and
the conjugation invariance in \stepref{rCsigns}.~\qed\end{proof}
\end{proof}

\begin{namedresult}{Lemma}{(three lines at forty-five degrees)}\label{lem:rem-three-lines}
\lean{IsRootSet.not_contains_three_equiangular_lines}\lean{IsRootSet.fortyFiveTriple_properties}
\uses{def:roots,lem:angles,facts:inner,facts:rem-angles}
\assume{$\Zu_1,\Zu_2,\Zu_3$ are orthonormal in $\mathbb R^3$ and $c=1/\sqrt2$.}
\pflet{$\Zv_1=\Zu_1$, $\Zv_2=c\Zu_1+c\Zu_2$, and
$\Zv_3=c\Zu_1+(1-c)\Zu_2+\sqrt{2c-1}\Zu_3$.}
\prove{These are unit vectors on three distinct lines pairwise at
$45^\circ$. There are no Root System $\Phi$ in any Euclidean space
and $\alpha,\beta,\gamma\in\Phi$ satisfying all three equations
\[
\begin{aligned}
4(\alpha,\beta)^2&=2\rsq\alpha\,\rsq\beta,\\
4(\alpha,\gamma)^2&=2\rsq\alpha\,\rsq\gamma,\\
4(\beta,\gamma)^2&=2\rsq\beta\,\rsq\gamma.
\end{aligned}
\]}
\end{namedresult}
\begin{proof}
\pfsketch\ Check their Gram matrix, then choose a shortest root on the
three lines in a putative containing Root System.
\step{rtlradicand}{$c^2=1/2$, $1/2<c<1$, and $2c-1>0$.}
\begin{proof}\pf\ Use $1<\sqrt2<2$ and $(\sqrt2)^2=2$ in the definition of $c$.~\qed\end{proof}
\step{rtlnorm1}{$\rsq{\Zv_1}=1$.}
\begin{proof}\pf\ Use orthonormality. ~\qed\end{proof}
\step{rtlnorm2}{$\rsq{\Zv_2}=2c^2=1$.}
\begin{proof}\pf\ Sum its coordinate squares using \stepref{rtlradicand}.~\qed\end{proof}
\step{rtlnorm3}{$\rsq{\Zv_3}=c^2+(1-c)^2+(2c-1)=1$.}
\begin{proof}\pf\ Sum its coordinate squares and use \stepref{rtlradicand}.~\qed\end{proof}
\step{rtlip12}{$(\Zv_1,\Zv_2)=c$.}
\begin{proof}\pf\ Read the first coordinate of $\Zv_2$.~\qed\end{proof}
\step{rtlip13}{$(\Zv_1,\Zv_3)=c$.}
\begin{proof}\pf\ Read the first coordinate of $\Zv_3$.~\qed\end{proof}
\step{rtlip23}{$(\Zv_2,\Zv_3)=c^2+c(1-c)=c$.}
\begin{proof}\pf\ Sum their coordinate products. ~\qed\end{proof}
\step{rtlangle}{The three lines are distinct and pairwise at $45^\circ$.}
\begin{proof}\pf\ Use the unit norms \stepref{rtlnorm1}--\stepref{rtlnorm3},
the inner products \stepref{rtlip12}--\stepref{rtlip23}, and
$0<c<1$ in \ref{facts:rem-angles}.~\qed\end{proof}
\step{rtlimpossible}{\assume{A Root System $\Phi$ and roots
$\alpha,\beta,\gamma\in\Phi$ satisfy the three equations in the statement.}
\prove{\textsc{False}.}}
\begin{proof}
\step{rtlpairs}{The three root pairs are nonperpendicular, nonproportional,
and have line angle $45^\circ$.}
\begin{proof}\pf\ Their squared cosines are $1/2$ by \asm{1}
and the positive squared norms of nonzero roots in
Definition~\ref{def:roots}; use \ref{facts:rem-angles}.~\qed\end{proof}
\step{rtlminimum}{\pick\ a permutation of the three roots with $\alpha$
of minimum squared length $b=\rsq\alpha>0$.}
\begin{proof}\pf\ A finite nonempty set of three squared lengths
has a minimum; the three equations are unchanged by the permutation.~\qed\end{proof}
\step{rtlbeta}{$\rsq\beta=2b$.}
\begin{proof}\pf\ Apply Lemma~\ref{lem:angles} to the $45^\circ$ pair
$\alpha,\beta$ from \stepref{rtlpairs} and minimality in
\stepref{rtlminimum}.~\qed\end{proof}
\step{rtlgamma}{$\rsq\gamma=2b$.}
\begin{proof}\pf\ Apply Lemma~\ref{lem:angles} to the $45^\circ$ pair
$\alpha,\gamma$ from \stepref{rtlpairs} and minimality in
\stepref{rtlminimum}.~\qed\end{proof}
\qedstep
\begin{proof}\pf\ The equal lengths in \stepref{rtlbeta},
\stepref{rtlgamma} contradict Lemma~\ref{lem:angles} for the
$45^\circ$ pair $\beta,\gamma$.~\qed\end{proof}
\end{proof}
\qedstep
\begin{proof}\pf\ Combine \stepref{rtlnorm1}--\stepref{rtlnorm3},
\stepref{rtlangle}, \stepref{rtlimpossible}.~\qed\end{proof}
\end{proof}

\begin{namedresult}{Theorem}{C}\label{thm:C}
\lean{IsRootSet.remark_lines}\lean{RootSetModel.Bn_Cn_same_lines}
\lean{RootSetModel.signedAdj_quadratic_Bn}\lean{RootSetModel.signedAdj_quadratic_Cn}
\lean{RootSetModel.signedAdj_quadratic_F₄}\lean{RootSetModel.signedAdj_quadratic_G₂}
\lean{IsRootSet.exists_rootRepresentatives}
\uses{lem:angles,lem:counts,lem:models-irreducible,lem:tight-frame,lem:C-convention,lem:rem-three-lines,def:rem-classical,def:rem-exceptional,facts:rem-angles,def:roots}
\assume{\begin{pfenum}
\item Root Systems are as in Definition~\ref{def:roots}; every line
 matrix is $C=2U^TU-2I$ with unit columns.
\item For the classical line assertions $n\geq2$, and the models are
 those of Definitions~\ref{def:rem-classical}, \ref{def:rem-exceptional}.
\end{pfenum}}
\prove{Nonperpendicular, nonproportional root pairs have exactly the
length-ratio/angle alternatives in Lemma~\ref{lem:angles}. Each occurring
fixed-length subset is a Root System in its own span and an orthogonal
sum of systems similar to $\typ A_\ell$ ($\ell\geq1$), $\typ D_\ell$
($\ell\geq4$), $\typ E_6,\typ E_7,\typ E_8$.
The $\typ B_n,\typ C_n$ models have the same $n^2$ lines in dimension $n$,
with allowed angles $90^\circ,60^\circ,45^\circ$, and
$(C+2I)(C-2(n-1)I)=0$.
Type $\typ F_4$ has 24 lines in dimension 4, with those three angles,
and $(C+2I)(C-10I)=0$.
Type $\typ G_2$ has 6 lines in dimension 2, with angles
$90^\circ,60^\circ,30^\circ$, and $(C+2I)(C-4I)=0$.
For $n=2$ the $60^\circ$ classical angle does not occur; for $n\geq3$
all three classical angles occur. All listed angles occur in $\typ F_4$
and $\typ G_2$. Three lines pairwise at $45^\circ$ exist in $\mathbb R^3$
and are contained in no Root System in any ambient dimension.}
In Lean the matrix equations hold for every choice of one root on each line, with a unit vector
on the same line. Such a choice of roots exists, indexed by $\{1,\dots,N\}$ for the number $N$ of
lines (\texttt{IsRootSet.exists\_rootRepresentatives}), so the equations are not vacuous.
\end{namedresult}
\begin{proof}
\pfsketch\ Use the common line frame for the matrix equations, and give
witness pairs for each claimed occurring angle.
\step{rCgeneral}{The root-pair and fixed-length assertions hold.}
\begin{proof}\pf\ Apply Lemma~\ref{lem:angles}.~\qed\end{proof}
\step{rCclassicalsame}{The $\typ B_n$ and $\typ C_n$ line sets coincide.}
\begin{proof}
\step{rCsameaxes}{Their axis roots lie on exactly the same axes.}
\begin{proof}\pf\ The vectors $\Zu_i$ and $2\Zu_i$ are nonzero
multiples in Definition~\ref{def:rem-classical}.~\qed\end{proof}
\step{rCsamepairs}{Their two-coordinate roots are identical.}
\begin{proof}\pf\ Compare their second clauses in
Definition~\ref{def:rem-classical}.~\qed\end{proof}
\qedstep
\begin{proof}\pf\ The two patterns in \stepref{rCsameaxes},
\stepref{rCsamepairs} exhaust both root sets. ~\qed\end{proof}
\end{proof}
\step{rCcounts}{The line counts and dimensions are
$(n^2,n)$, $(24,4)$, $(6,2)$ in the classical, $\typ F$, and $\typ G$ cases.}
\begin{proof}\pf\ Apply Lemma~\ref{lem:counts} and the coordinate-space
definitions. ~\qed\end{proof}
\step{rCallowed}{The allowed angles are the asserted ones.}
\begin{proof}
\step{rCallowedBF}{In $\typ B,\typ C,\typ F$ a nonorthogonal distinct-line
pair has angle $60^\circ$ or $45^\circ$.}
\begin{proof}\pf\ Lemma~\ref{lem:counts} gives equal lengths or ratio
$\sqrt2$; apply Lemma~\ref{lem:angles}.~\qed\end{proof}
\step{rCallowedG}{In $\typ G$ such a pair has angle $60^\circ$ or $30^\circ$.}
\begin{proof}\pf\ Lemma~\ref{lem:counts} gives equal lengths or ratio
$\sqrt3$; apply Lemma~\ref{lem:angles}.~\qed\end{proof}
\qedstep
\begin{proof}\pf\ Adjoin the perpendicular case $90^\circ$ to
\stepref{rCallowedBF}, \stepref{rCallowedG}.~\qed\end{proof}
\end{proof}
\step{rCoccurrence}{All claimed angles occur, with the stated rank-2 exception.}
\begin{proof}
\step{rCright}{The axes $\mathbb R\Zu_1,\mathbb R\Zu_2$ give $90^\circ$
in $\typ B_n,\typ C_n,\typ F_4$.}
\begin{proof}\pf\ Their product is zero, and both lines occur in the
coordinate lists. ~\qed\end{proof}
\step{rCfortyfive}{The lines of $\Zu_1,\Zu_1+\Zu_2$ give $45^\circ$
in those models.}
\begin{proof}\pf\ Their squared cosine is $1/(1\cdot2)=1/2$;
use \ref{facts:rem-angles}.~\qed\end{proof}
\step{rCsixty}{For $n\geq3$, the lines of $\Zu_1+\Zu_2,\Zu_1+\Zu_3$
give $60^\circ$ in $\typ B_n,\typ C_n$ and also in $\typ F_4$.}
\begin{proof}\pf\ Their squared cosine is $1/(2\cdot2)=1/4$.~\qed\end{proof}
\step{rCtwoangles}{For $n=2$ the classical models have no $60^\circ$ pair.}
\begin{proof}
\step{rCtwoaxes}{The two axes are perpendicular.}
\begin{proof}\pf\ Their inner product is zero. ~\qed\end{proof}
\step{rCtwodiagonals}{The two diagonal lines are perpendicular.}
\begin{proof}\pf\ $(\Zu_1+\Zu_2,\Zu_1-\Zu_2)=1-1=0$.~\qed\end{proof}
\step{rCtwomixed}{Every axis-diagonal pair has squared cosine $1/2$.}
\begin{proof}\pf\ Its absolute inner product is 1 and its squared norms
are 1 and 2. ~\qed\end{proof}
\qedstep
\begin{proof}\pf\ These three pair types exhaust the four lines
in Definition~\ref{def:rem-classical}.~\qed\end{proof}
\end{proof}
\step{rCGsixty}{The short roots $\Zu_1-\Zu_2,\Zu_1-\Zu_3$ give
$60^\circ$ in $\typ G_2$.}
\begin{proof}\pf\ Their squared cosine is $1/(2\cdot2)=1/4$.~\qed\end{proof}
\step{rCGright}{The roots $2\Zu_1-\Zu_2-\Zu_3,\Zu_2-\Zu_3$
give $90^\circ$.}
\begin{proof}\pf\ Their inner product is $-1+1=0$.~\qed\end{proof}
\step{rCGthirty}{The roots $2\Zu_1-\Zu_2-\Zu_3,\Zu_1-\Zu_2$
give $30^\circ$.}
\begin{proof}\pf\ Their squared cosine is $3^2/(6\cdot2)=3/4$.~\qed\end{proof}
\qedstep
\begin{proof}\pf\ Combine \stepref{rCright}--\stepref{rCGthirty}
using \ref{facts:rem-angles}.~\qed\end{proof}
\end{proof}
\step{rCequations}{All three asserted matrix equations hold.}
\begin{proof}
\step{rCframes}{The frame identity holds for each model.}
\begin{proof}\pf\ Apply Lemma~\ref{lem:tight-frame} to each
indecomposable Root System of Lemma~\ref{lem:models-irreducible}.~\qed\end{proof}
\step{rCeqBC}{For $\typ B_n$ and $\typ C_n$, $\theta=2(n-1)$.}
\begin{proof}\pf\ Substitute $N=n^2$ from \stepref{rCcounts}
and the frame \stepref{rCframes} in Lemma~\ref{lem:C-convention}.~\qed\end{proof}
\step{rCeqF}{For $\typ F_4$, $\theta=10$.}
\begin{proof}\pf\ Substitute $N=24$, dimension 4, and
\stepref{rCframes} in Lemma~\ref{lem:C-convention}.~\qed\end{proof}
\step{rCeqG}{For $\typ G_2$, $\theta=4$.}
\begin{proof}\pf\ Substitute $N=6$, dimension 2, and
\stepref{rCframes} in Lemma~\ref{lem:C-convention}.~\qed\end{proof}
\qedstep
\begin{proof}\pf\ Combine \stepref{rCeqBC}, \stepref{rCeqF}, \stepref{rCeqG}.~\qed\end{proof}
\end{proof}
\step{rCexample}{The asserted three-line example exists and cannot lie
in any Root System.}
\begin{proof}\pf\ Apply Lemma~\ref{lem:rem-three-lines}.~\qed\end{proof}
\qedstep
\begin{proof}\pf\ Combine \stepref{rCgeneral}, \stepref{rCclassicalsame},
\stepref{rCcounts}, \stepref{rCallowed}, \stepref{rCoccurrence},
\stepref{rCequations}, \stepref{rCexample}.~\qed\end{proof}
\end{proof}

\begin{namedresult}{Corollary}{(the equal-length matrix constants)}\label{cor:rem-constants}
\lean{RootSetModel.signedAdj_quadratic_simplyLaced}
\uses{lem:counts,lem:models-irreducible,lem:tight-frame,lem:C-convention}
\assume{The model is $\typ A_n$ ($n\geq1$), $\typ D_n$ ($n\geq4$),
or $\typ E_6,\typ E_7,\typ E_8$, in its span; $C=2U^TU-2I$
uses unit representatives of its lines.}
\prove{$(C+2I)(C-\theta I)=0$, with respective constants
$n-1,2(n-2),10,16,28$.}
\end{namedresult}
\begin{proof}
\pfsketch\ \repair\ The last constant is 28, correcting the printed
18 in Theorem 3.7.
\step{reqframe}{Each model satisfies the line frame identity.}
\begin{proof}\pf\ Apply Lemma~\ref{lem:tight-frame} using
Lemma~\ref{lem:models-irreducible}.~\qed\end{proof}
\step{reqA}{For $\typ A_n$, the second constant is $n-1$.}
\begin{proof}\pf\ Use $N=n(n+1)/2$ from Lemma~\ref{lem:counts}
in Lemma~\ref{lem:C-convention} with \stepref{reqframe}.~\qed\end{proof}
\step{reqD}{For $\typ D_n$, it is $2(n-2)$.}
\begin{proof}\pf\ Substitute $N=n(n-1)$ from Lemma~\ref{lem:counts}
in Lemma~\ref{lem:C-convention} with \stepref{reqframe}.~\qed\end{proof}
\step{reqE6}{For $\typ E_6$, it is $2(36/6)-2=10$.}
\begin{proof}\pf\ Substitute the $\typ E_6$ count in
Lemma~\ref{lem:C-convention} with \stepref{reqframe}.~\qed\end{proof}
\step{reqE7}{For $\typ E_7$, it is $2(63/7)-2=16$.}
\begin{proof}\pf\ Substitute the $\typ E_7$ count in
Lemma~\ref{lem:C-convention} with \stepref{reqframe}.~\qed\end{proof}
\step{reqE8}{For $\typ E_8$, it is $2(120/8)-2=28$.}
\begin{proof}\pf\ Substitute the $\typ E_8$ count in
Lemma~\ref{lem:C-convention} with \stepref{reqframe}.~\qed\end{proof}
\qedstep
\begin{proof}\pf\ The cases \stepref{reqA}--\stepref{reqE8}
exhaust the listed models. ~\qed\end{proof}
\end{proof}

\subsubsection{Finite certificates from the ordinary proof}

\begin{namedresult}{Lemma}{(leaf determinant expansion)}\label{lem:rem-leaf-det}
\lean{Matrix.det_eq_of_leaf}
\uses{facts:rem-library}
\assume{$M$ is a square matrix over a commutative ring, on a finite
index set; $v\ne w$, $M(v,v)=2$, and every off-diagonal entry in
row or column $v$ is zero except possibly the entry at $w$.}
\prove{Writing $p=M(v,w)M(w,v)$ and subscripts for principal deletions,
\begin{equation}\label{eq:leaf-det}
\det M=2\det M_{-v}-p\det M_{-v,-w}.
\end{equation}}
\end{namedresult}
\begin{proof}
\pfsketch\ Separate permutations fixing $v$ from those interchanging
$v,w$. This applies to a leaf of any support, not only a path.
\step{rldfixed}{Permutation terms fixing $v$ sum to $2\det M_{-v}$.}
\begin{proof}\pf\ Restrict those permutations to the complement of $v$
in \path{Matrix.det_apply'}; its diagonal factor is 2. ~\qed\end{proof}
\step{rldmoving}{A nonzero permutation term not fixing $v$ interchanges
$v$ and $w$.}
\begin{proof}
\step{rldcolumn}{Column $v$ must choose row $w$.}
\begin{proof}\pf\ All other available off-diagonal entries of column
$v$ vanish by \asm{0}.~\qed\end{proof}
\step{rldrow}{Row $v$ must choose column $w$.}
\begin{proof}\pf\ Column $v$ is already used by row $w$ in
\stepref{rldcolumn}, and all remaining entries of row $v$ except at $w$ vanish. ~\qed\end{proof}
\qedstep
\begin{proof}\pf\ Combine \stepref{rldcolumn}, \stepref{rldrow}.~\qed\end{proof}
\end{proof}
\step{rldsum}{The moving terms sum to $-p\det M_{-v,-w}$.}
\begin{proof}\pf\ Delete the transposition in \stepref{rldmoving}
from each permutation; it contributes sign $-1$ and entry product $p$
in \path{Matrix.det_apply'}.~\qed\end{proof}
\qedstep
\begin{proof}\pf\ Add the exhaustive permutation classes
\stepref{rldfixed}, \stepref{rldsum}.~\qed\end{proof}
\end{proof}

\begin{namedresult}{Lemma}{(leading determinants)}\label{comp:rem-minors}
\lean{CartanMatrix.model_leading_determinants}
\uses{def:rem-matrices,rem:small,lem:rem-leaf-det,def:rem-bases,facts:rem-library}
\assume{$M$ is a named matrix in \eqref{eq:list};
$\Delta_k$ is its leading $k$-row principal determinant and $\Delta_0=1$.}
\prove{The complete lists are
\begin{equation}\label{eq:minors}
\begin{array}{c|l}
\typ A_\ell&\Delta_k=k+1\\
\typ B_\ell,\typ C_\ell&\Delta_k=k+1\ (k<\ell),\quad\Delta_\ell=2\\
\typ D_\ell&\Delta_k=k+1\ (k<\ell),\quad\Delta_\ell=4\\
\typ E_6&(2,4,6,5,4,3)\\
\typ E_7&(2,4,6,5,4,3,2)\\
\typ E_8&(2,4,6,5,4,3,2,1)\\
\typ F_4&(2,3,2,1)\\
\typ G_2&(2,1).
\end{array}
\end{equation}}
\end{namedresult}
\begin{proof}
\pfsketch\ Apply the leaf formula with the actual deleted neighbour,
then solve the path recurrence.
\step{rmdpath}{For a path extension with last edge product $p$,
$\Delta_k=2\Delta_{k-1}-p\Delta_{k-2}$.}
\begin{proof}\pf\ Apply Lemma~\ref{lem:rem-leaf-det} to its last vertex. ~\qed\end{proof}
\step{rmdA}{The unit-path determinants are $k+1$.}
\begin{proof}
\step{rmdAbase}{$\Delta_0=1$ and $\Delta_1=2$.}
\begin{proof}\pf\ Evaluate the empty and one-entry determinants. ~\qed\end{proof}
\step{rmdAstep}{\assume{$\Delta_{k-1}=k$ and $\Delta_{k-2}=k-1$.}
\prove{$\Delta_k=k+1$.}}
\begin{proof}\pf\ Substitute in \stepref{rmdpath} with $p=1$.~\qed\end{proof}
\qedstep
\begin{proof}\pf\ Induction using \stepref{rmdAbase}, \stepref{rmdAstep}.~\qed\end{proof}
\end{proof}
\step{rmdBC}{The $\typ B,\typ C$ final determinants are 2.}
\begin{proof}\pf\ Their last edge product is 2, so
\stepref{rmdpath}, \stepref{rmdA} give $2\ell-2(\ell-1)=2$.~\qed\end{proof}
\step{rmdD}{The $\typ D_\ell$ final determinant is 4.}
\begin{proof}
\step{rmdDleaf}{Deleting $v=\ell-1$ leaves a unit path of determinant $\ell$.}
\begin{proof}\pf\ Use Definition~\ref{def:rem-matrices} and \stepref{rmdA}.~\qed\end{proof}
\step{rmdDneighbour}{Deleting also $w=\ell-3$ leaves a path of
$\ell-3$ vertices and an isolated vertex, of determinant $2(\ell-2)$.}
\begin{proof}\pf\ Factor the isolated diagonal entry 2 in the
determinant expansion and apply \stepref{rmdA}.~\qed\end{proof}
\qedstep
\begin{proof}\pf\ Lemma~\ref{lem:rem-leaf-det} with $p=1$,
\stepref{rmdDleaf}, \stepref{rmdDneighbour} gives $2\ell-2(\ell-2)=4$.~\qed\end{proof}
\end{proof}
\step{rmdE}{The $\typ E$ prefix determinants are $2,4,6,5,4,3,2,1$.}
\begin{proof}
\step{rmdE1}{The first determinant is 2.}
\begin{proof}\pf\ Evaluate the single diagonal entry. ~\qed\end{proof}
\step{rmdE2}{The second is 4.}
\begin{proof}\pf\ The leading two vertices are isolated, giving $2\cdot2$.~\qed\end{proof}
\step{rmdE3}{The third is 6.}
\begin{proof}\pf\ Its single unit edge contributes determinant 3 and
its isolated vertex contributes 2. ~\qed\end{proof}
\step{rmdE4}{The fourth is 5.}
\begin{proof}\pf\ The order in Remark~\ref{rem:small} makes it a
four-vertex unit path; apply \stepref{rmdA}.~\qed\end{proof}
\step{rmdErest}{The next four determinants are $4,3,2,1$.}
\begin{proof}\pf\ Its leading $k$-row matrix is exactly
$\typ E_k$ by \path{CartanMatrix.E}; for $k=5,6,7,8$, apply
\path{CartanMatrix.E_det}, which gives $9-k$ for every $k\geq3$.~\qed\end{proof}
\qedstep
\begin{proof}\pf\ Concatenate \stepref{rmdE1}--\stepref{rmdErest}.~\qed\end{proof}
\end{proof}
\step{rmdF}{The $\typ F$ determinants are $2,3,2,1$.}
\begin{proof}
\step{rmdFprefix}{The first three determinants are $2,3,2$.}
\begin{proof}\pf\ Substitute the edge products $1,2$ in
\stepref{rmdpath}, starting from $\Delta_0=1,\Delta_1=2$.~\qed\end{proof}
\step{rmdFfinal}{The fourth determinant is 1.}
\begin{proof}\pf\ Apply \texttt{CartanMatrix.F}\textsubscript{4}\texttt{\_det} in
\path{Mathlib/LinearAlgebra/Matrix/Cartan.lean}.~\qed\end{proof}
\qedstep
\begin{proof}\pf\ Combine \stepref{rmdFprefix}, \stepref{rmdFfinal}.~\qed\end{proof}
\end{proof}
\step{rmdG}{The $\typ G$ determinants are $2,1$.}
\begin{proof}\pf\ The first diagonal entry is 2; its final
determinant is 1 by \texttt{CartanMatrix.G}\textsubscript{2}\texttt{\_det} in
\path{Mathlib/LinearAlgebra/Matrix/Cartan.lean}.~\qed\end{proof}
\qedstep
\begin{proof}\pf\ Proper classical prefixes are unit paths; combine
\stepref{rmdA}--\stepref{rmdG}.~\qed\end{proof}
\end{proof}

\begin{namedresult}{Lemma}{(rational positivity certificates)}\label{comp:rem-ldl}
\lean{RootSetModel.exceptional_ldl}\lean{RootSetModel.symmetrizer_eq_modelSymmetrizer}
\uses{def:rem-matrices}
\assume{$M$ is $\typ E_6,\typ E_7,\typ E_8,\typ F_4$ or $\typ G_2$,
and $H=\operatorname{diag}(d)M$ with $d$ from \eqref{eq:symmetrizers}.}
\prove{There is a lower triangular $L$ with diagonal 1 such that
\begin{equation}\label{eq:ldl}
H=L\operatorname{diag}(q)L^T,\qquad
\Zx^TH\Zx=\sum_i q(i)((L^T\Zx)(i))^2,
\end{equation}
with every $q(i)>0$. The certificate data are:
\[
\begin{array}{c|l}
\typ E_8&q=(2,2,3/2,5/6,4/5,3/4,2/3,1/2)\\
\typ F_4&q=(2,3/2,4/3,1)\\
\typ G_2&q=(2,3/2).
\end{array}
\]
For $\typ E_8$ the nonzero strict lower entries are
\[
\begin{split}
L(2,0)&=-1/2,\quad L(3,1)=-1/2,\quad L(3,2)=-2/3,\\
L(4,3)&=-6/5,\quad L(5,4)=-5/4,\quad
L(6,5)=-4/3,\quad L(7,6)=-3/2.
\end{split}
\]
Use the leading six or seven entries and rows for $\typ E_6,\typ E_7$.
For $\typ F_4$ use $L(1,0)=-1/2$, $L(2,1)=-4/3$, $L(3,2)=-3/2$.
For $\typ G_2$ use $L(1,0)=-3/2$; all unspecified strict lower entries
are zero. These data prove $H$ positive definite.}
\end{namedresult}
\begin{proof}
\pfsketch\ These are bounded rational matrix identities, followed by
back substitution in a triangular matrix.
\step{rldlE}{The eight-row data satisfy
$H(i,j)=\sum_k L(i,k)q(k)L(j,k)$ for all $i,j<8$.}
\begin{proof}\pf\ Substitute the displayed rational entries and the
$\typ E_8$ entries of Definition~\ref{def:rem-matrices} in that finite sum. ~\qed\end{proof}
\step{rldlprefix}{The same equality holds for its six- and seven-row prefixes.}
\begin{proof}\pf\ Lower triangularity makes every omitted summand
zero in the prefix entries of \stepref{rldlE}.~\qed\end{proof}
\step{rldlF}{The equality holds for the four-row $\typ F$ data.}
\begin{proof}\pf\ Substitute those rational entries and $d=(1,1,2,2)$
in the finite entry sum. ~\qed\end{proof}
\step{rldlG}{The equality holds for the two-row $\typ G$ data.}
\begin{proof}\pf\ Substitute those rational entries and $d=(1,3)$
in the finite entry sum. ~\qed\end{proof}
\step{rldlform}{The quadratic identity in \eqref{eq:ldl} holds.}
\begin{proof}\pf\ Multiply the identities
\stepref{rldlE}--\stepref{rldlG} on the left and right by $\Zx^T,\Zx$.~\qed\end{proof}
\step{rldlinjective}{\assume{$L^T\Zx=0$.}
\prove{$\Zx=0$.}}
\begin{proof}\pf\ Backward induction on the coordinates of the upper
triangular matrix $L^T$ with diagonal 1. ~\qed\end{proof}
\qedstep
\begin{proof}\pf\ For $\Zx\ne0$, \stepref{rldlinjective} makes at least
one square in \stepref{rldlform} positive; every displayed $q(i)$ is
positive, so the sum is positive. ~\qed\end{proof}
\end{proof}

\begin{namedresult}{Lemma}{(rational coordinate frame)}\label{comp:rem-rational-frame}
\lean{RootSetModel.rational_frame}\lean{RootSetModel.exists_rat_simpleRootMatrix}
\lean{RootSetModel.exists_rat_inv_gram_simpleRootMatrix}
\uses{lem:classical,lem:E,lem:F,lem:G,lem:models-irreducible,lem:counts,lem:tight-frame,lem:rem-projectors,facts:inner,facts:rem-dimension}
\assume{A coordinate model from the construction ranges has $N$ root
lines and rank $n$; the matrix $A$ has its simple roots as columns in the
given coordinate space, which may be larger than its span.}
\prove{\begin{equation}\label{eq:rational-frame}
\frac12\sum_{\alpha\in\Phi}\frac{\alpha\alpha^T}{\rsq\alpha}
=\frac Nn\,A(A^TA)^{-1}A^T.
\end{equation}
All entries in this equation are rational.}
\end{namedresult}
\begin{proof}
\pfsketch\ Both sides act by $N/n$ on the root span and vanish on its
orthogonal complement.
\step{rrfprojection}{$P=A(A^TA)^{-1}A^T$ is the orthogonal projection
onto the column span.}
\begin{proof}
\step{rrfinvertible}{$A^TA$ is invertible.}
\begin{proof}\pf\ A null coefficient vector has combination of squared
norm zero, thus zero; the model bases are independent by
Lemmas~\ref{lem:classical}, \ref{lem:E}, \ref{lem:F}, \ref{lem:G}.~\qed\end{proof}
\step{rrffix}{$P(Ax)=Ax$.}
\begin{proof}\pf\ Cancel $(A^TA)^{-1}(A^TA)$ using \stepref{rrfinvertible}.~\qed\end{proof}
\step{rrfkill}{\assume{$A^T\Zy=0$.}
\prove{$P\Zy=0$.}}
\begin{proof}\pf\ Substitute in the definition of $P$.~\qed\end{proof}
\qedstep
\begin{proof}\pf\ The column span and its orthogonal complement are
the two summands tested in \stepref{rrffix}, \stepref{rrfkill}.~\qed\end{proof}
\end{proof}
\step{rrfleft}{The left side is the sum of the unit-line orthogonal projections.}
\begin{proof}\pf\ Each line has exactly its two opposite roots by
Lemma~\ref{lem:counts}; both summands are its orthogonal projection by
Lemma~\ref{lem:rem-projectors}.~\qed\end{proof}
\step{rrfequality}{The two operators agree on the ambient coordinate space.}
\begin{proof}\pf\ Lemma~\ref{lem:tight-frame} gives $(N/n)I$ on the
root span for \stepref{rrfleft}, and every root projection annihilates
the orthogonal complement; use \stepref{rrfprojection}.~\qed\end{proof}
\step{rrfrational}{All entries are rational.}
\begin{proof}\pf\ The model roots have rational coordinates and
positive rational squared norms; an invertible rational matrix has
rational inverse by \path{Matrix.inv_def}.~\qed\end{proof}
\qedstep
\begin{proof}\pf\ Combine \stepref{rrfequality}, \stepref{rrfrational}.~\qed\end{proof}
\end{proof}

\begin{cgss}{Remark}{3.11}\label{rem:3.11}
\uses{thm:A,thm:B,thm:C,cor:rem-constants}
Theorems A, B and C provide the finite-type classification, its passage
to Euclidean Root Systems, and all the unequal-length line assertions
of this remark. Corollary~\ref{cor:rem-constants} also recovers the
equal-length equations with the corrected $\typ E_8$ constant.
The classification does not use [CGSS, Theorem 3.10] as a premise.
Theorem 3.12 is proved separately (Section~\ref{sec:thm312}).
\end{cgss}
\endgroup
